% !TEX root = ../Elements.tex
% Greek and English text last checked 15/03/2013
%%%%%%
% BOOK 3
%%%%%%
\pdfbookmark[0]{Book 3}{book3}
\pagestyle{empty}
\begin{center}
{\Huge ELEMENTS BOOK 3}\\
\spa\spa\spa
{\huge\it Fundamentals of Plane Geometry Involving}\\[0.5ex] {\huge\it Circles}
\end{center}
%\vspace{2cm}
%\epsfysize=5in
%\centerline{\epsffile{Book03/front.eps}}
\newpage

%%%%%%%
% Definitions
%%%%%%%
\pdfbookmark[1]{Definitions}{def3}
\pagestyle{fancy}
\cfoot{{\greekfont τ̀ηεπαγε}}
\lhead{\large{\greekfont ΣΤΟΙΘΕΙΩΝ γ̀γν{3}.}}
\rhead{\large ELEMENTS BOOK 3}
\begin{Parallel}{}{} 
\ParallelLText{
\begin{center}
\large{{\greekfont Ὅροι}.}
\end{center}\vspace*{-7pt}

\ggn{1}.~{\greekfont Ἴσοι κύκλοι εἰσίν, ὧν αἱ διάμετροι ἴσαι εἰσίν,
ἢ ὧν αἱ ἐκ τῶν κέντρων ἴσαι εἰσίν.}

\ggn{2}.~{\greekfont Εὐθεῖα κύκλου ἐφάπτεσθαι λέγεται, ἥτις ἁπτομένη τοῦ
κύκλου καὶ ἐκβαλλομένη οὐ τέμνει τὸν κύκλον.}

\ggn{3}.~{\greekfont Κύκλοι ἐφάπτεσθαι ἀλλήλων λέγονται οἵτινες ἁπτό-μενοι ἀλλήλων οὐ τέμνουσιν ἀλλήλους.}

\ggn{4}.~{\greekfont Ἐν κύκλῳ ἴσον ἀπέθειν ἀπὸ τοῦ κέντρου εὐθεῖαι
λέγονται, ὅταν αἱ ἀπὸ τοῦ κέντρου ἐπ α ὐτὰς κάθετοι ἀγόμεναι
ἴσαι ὦσιν.}

\ggn{5}.~{\greekfont Μεῖζον δὲ ἀπέθειν λέγεται, ἐφ ἣν ἡ μείζων κάθετος πίπτει.}

\ggn{6}.~{\greekfont Τμῆμα κύκλου ἐστὶ τὸ περιεχόμενον σχῆμα ὑπό τε εὐθείας καὶ κύκλου περιφερείας.}

\ggn{7}.~{\greekfont Τμήματος δὲ γωνία ἐστὶν ἡ περιεθομένη ὑπό
τε εὐθείας καὶ κύκλου περιφερείας.}

\ggn{8}.~{\greekfont Ἐν τμήματι δὲ γωνία ἐστίν, ὅταν ἐπὶ τῆς περιφερείας τοῦ τμήματος ληφθῇ τι σημεῖον καὶ ἀπ α ὐτοῦ ἐπὶ τὰ πέρατα
τῆς εὐθείας, ἥ ἐστι βάσις τοῦ τμήματος, ἐπιζευθθῶσιν εὐθεῖαι,
ἡ περιεθομένη γωνία ὑπὸ τῶν ἐπιζευθθεισῶν εὐθειῶν.}

\ggn{9}.~{\greekfont Ὅταν δὲ αἱ περιέθουσαι τὴν γωνίαν εὐθεῖαι ἀπολαμβάνωσί τινα περιφέρειαν, ἐπ ἐκείνης λέγεται βεβηκέναι ἡ γωνία.}

\ggn{10}.~{\greekfont Τομεὺς δὲ κύκλου ἐστίν, ὅταν πρὸς τῷ κέντρῷ τοῦ
κύκλου συσταθῇ γωνία, τὸ περιεχόμενον σχῆμα ὑπό τε τῶν τὴν
γωνίαν περιεθουσῶν
εὐθειῶν καὶ τῆς ἀπολαμβανομένης ὑπ α ὐτῶν περιφερείας.}

\ggn{11}.~{\greekfont Ὅμοία τμήματα κύκλων ἐστὶ τὰ δεθόμενα γωνίας ἴσας, ἤ ἐν οἷς αἱ γωνίαι ἴσαι ἀλλήλαις εἰσίν.}}

\ParallelRText{
\begin{center}
{\large Definitions}
\end{center}

1.~Equal circles are (circles) whose diameters are equal,  or whose (distances)
from the centers (to the circumferences) are equal (i.e., whose radii are equal).

2.~A straight-line said to touch a circle  is any (straight-line) which, meeting the circle and being produced, does not cut the circle.

3.~Circles said to touch one another are any (circles) which, meeting one another, do not cut one another.

4.~In a circle, straight-lines are said to be equally far from the center
when the perpendiculars drawn to them from the center are equal.

5.~And (that straight-line) is said to be further (from the center) on
which the greater perpendicular falls (from the center).

6. ~A segment of a circle is the figure contained by a straight-line
and a circumference of a circle.

7.~And the angle of a segment is that contained by a straight-line
and a circumference of a circle.

8.~And the angle in a segment is the
angle contained by the joined straight-lines, when any
point  is taken on the circumference of a segment, 
and straight-lines are joined from it to the ends of the straight-line which is the base of the
segment.

9. ~And when the straight-lines containing an angle cut off some circumference, the angle is said to stand upon that (circumference).

10.~And a sector of a circle is the figure contained by the
straight-lines surrounding an angle, and the circumference cut off by them,
when the angle is constructed at the center of a circle.

11.~Similar segments of circles are those accepting equal angles, or
in which the angles are equal to one another.}
\end{Parallel}

%%%%%%
% Prop 3.1
%%%%%%
\pdfbookmark[1]{Proposition 3.1}{pdf3.1}
\begin{Parallel}{}{} 
\ParallelLText{
\begin{center}
{\large \ggn{1}.}
\end{center}\vspace*{-7pt}

{\greekfont Τοῦ δοθέντος κύκλου τὸ κέντρον εὑρεῖν.}

{\greekfont Ἔστω ὁ δοθεὶς κύκλος ὁ ΑΒΓ; δεῖ δὴ τοῦ ΑΒΓ κύκλου
τὸ κέντρον εὑρεῖν.}

{\greekfont Διήθθω τις εἰς α ὐτόν, ὡς ἔτυθεν, εὐθεῖα ἡ ΑΒ, καὶ τετμήσθω δίθα
κατὰ τὸ Δ σημεῖον, καὶ ἀπὸ τοῦ Δ τῇ ΑΒ πρὸς ὀρθὰς ἤθθω
ἡ ΔΓ καὶ διήθθω ἐπὶ τὸ Ε, καὶ τετμήσθω ἡ ΓΕ δίθα κατὰ τὸ Ζ;
λέγω, ὅτι τὸ Ζ κέντρον ἐστὶ τοῦ ΑΒΓ [κύκλου].}

{\greekfont μὴ γάρ, ἀλλ εἰ δυνατόν, ἔστω τὸ Η, καὶ ἐπεζεύθθωσαν αἱ ΗΑ, ΗΔ, ΗΒ. καὶ ἐπεὶ ἴση ἐστὶν ἡ ΑΔ τῇ ΔΒ, κοινὴ δὲ ἡ ΔΗ, δύο δὴ
αἱ ΑΔ, ΔΗ δύο ταῖς ΗΔ, ΔΒ ἴσαι εἰσὶν ἑκατέρα ἑκατέρᾳ;
καὶ βάσις ἡ ΗΑ βάσει τῇ ΗΒ ἐστιν ἴση; ἐκ κέντρου γάρ;
γωνία ἄρα ἡ ὑπὸ ΑΔΗ γωνίᾳ τῇ ὑπὸ ΗΔΒ ἴση ἐστίν.
ὅταν δὲ εὐθεῖα ἐπ εὐθεῖαν σταθεῖσα τὰς ἐφεχῆς γωνίας
ἴσας ἀλλήλαις ποιῇ, ὀρθὴ ἑκατέρα τῶν ἴσων γωνιῶν ἐστιν; ὀρθὴ
ἄρα ἐστὶν ἡ  ὑπὸ ΗΔΒ. ἐστὶ δὲ καὶ ἡ ὑπὸ ΖΔΒ ὀρθή;
ἴση ἄρα ἡ ὑπὸ ΖΔΒ τῇ ὑπὸ ΗΔΒ, ἡ μείζων τῇ ἐλάττονι;
ὅπερ ἐστὶν ἀδύνατον. οὐκ ἄρα τὸ Η κέντρον ἐστὶ τοῦ
ΑΒΓ κύκλου. ὁμοίως δὴ δείχομεν, ὅτι οὐδ ἄλλο
τι πλὴν τοῦ Ζ.}\\~\\~\\~\\

\epsfysize=2.2in
\centerline{\epsffile{Book03/fig01g.eps}}

{\greekfont Τὸ Ζ ἄρα σημεῖον κέντρον ἐστὶ τοῦ ΑΒΓ [κύκλου].}

\begin{center}\mbox{}\\
{\large {\greekfont Πόρισμα}.}
\end{center}\vspace*{-7pt}

{\greekfont Ἐκ δὴ τούτου φανερόν, ὅτι ἐὰν ἐν κύκλῳ εὐθεῖά
τις εὐθεῖάν τινα δίθα καὶ πρὸς ὀρθὰς τέμνῃ, ἐπὶ
τῆς τεμνούσης ἐστὶ τὸ κέντρον τοῦ κύκλου.} --- {\greekfont ὅπερ ἔδει
ποιῆσαι.}}

\ParallelRText{
\begin{center}
{\large Proposition 1}
\end{center}

To find the center of a given circle.

Let $ABC$ be the given circle. So it is required to find the center of circle $ABC$.

Let some straight-line $AB$ be drawn through ($ABC$), at random,
and let ($AB$) be cut in half at point $D$ 
[Prop.~1.9]. And let
$DC$ be drawn from $D$, at right-angles to $AB$ [Prop.~1.11].
And let ($CD$) be drawn through to $E$. And let $CE$ be cut in half
at $F$ [Prop.~1.9]. I say that (point) $F$ is the center of the [circle] $ABC$.

For (if) not (then), if possible, let $G$ (be the center of the circle), and
let $GA$, $GD$, and $GB$ be joined. And since $AD$ is equal to $DB$,
and $DG$ (is) common, the two (straight-lines) $AD$, $DG$
are equal to the two (straight-lines) $BD$, $DG$,$^\dag$ respectively. And the
base $GA$ is equal to the base $GB$. For (they are both) radii. Thus, angle
$ADG$ is equal to angle $GDB$ [Prop.~1.8]. And when a straight-line
stood upon (another) straight-line make  adjacent angles (which are)
equal to one another, each of the equal angles is a right-angle [Def.~1.10].
Thus, $GDB$ is a right-angle. And $FDB$ is also a right-angle. Thus,
$FDB$ (is) equal to $GDB$, the greater to the lesser. The very thing
is impossible. Thus, (point) $G$ is not the center of the circle $ABC$. 
So, similarly, we can show that neither is any other (point) except $F$.

\epsfysize=2.2in
\centerline{\epsffile{Book03/fig01e.eps}}

Thus, point $F$ is the center of the [circle] $ABC$.

\begin{center}\mbox{}\\
{\large Corollary}
\end{center}

So, from this, (it is) manifest that if any straight-line in a circle
cuts any (other) straight-line in half, and at right-angles, (then) the center
of the circle is on the former (straight-line). --- (Which is) the very thing it
was required to do.}
\end{Parallel}
{\footnotesize \noindent$^\dag$  The Greek text has ``$GD$, $DB$'', which is obviously a mistake.} 

%%%%%%
% Prop 3.2
%%%%%%
\pdfbookmark[1]{Proposition 3.2}{pdf3.2}
\begin{Parallel}{}{} 
\ParallelLText{
\begin{center}
{\large \ggn{2}.}
\end{center}\vspace*{-7pt}

{\greekfont Ἐὰν κύκλου ἐπὶ τῆς περιφερείας ληφθῇ δύο
τυθόντα σημεῖα, ἡ ἐπὶ τὰ σημεῖα ἐπιζευγνυμένη εὐθεῖα
ἐντὸς πεσεῖται τοῦ κύκλου.}

{\greekfont Ἔστω κύκλος ὁ ΑΒΓ, καὶ ἐπὶ τῆς περιφερείας α ὐτοῦ εἰλήφθω
δύο τυθόντα σημεῖα τὰ Α, Β; λέγω, ὅτι ἡ ἀπὸ τοῦ Α ἐπὶ τὸ Β
ἐπιζευγνυμένη εὐθεῖα ἐντὸς πεσεῖται τοῦ κύκλου.}

{\greekfont μὴ γάρ, ἀλλ εἰ δυνατόν, πιπτέτω ἐκτὸς ὡς ἡ ΑΕΒ, καὶ εἰλήφθω
τὸ κέντρον τοῦ ΑΒΓ κύκλου, καὶ ἔστω τὸ Δ, καὶ ἐπεζεύθθωσαν αἱ
ΔΑ, ΔΒ, καὶ διήθθω ἡ ΔΖΕ.}

{\greekfont Ἐπεὶ οὖν ἴση ἐστὶν ἡ ΔΑ τῇ ΔΒ, ἴση ἄρα
καὶ γωνία ἡ ὑπὸ ΔΑΕ τῇ ὑπὸ ΔΒΕ; καὶ ἐπεὶ τριγώνου
τοῦ ΔΑΕ μία πλευρὰ προσεκβέβληται ἡ ΑΕΒ, μείζων
ἄρα ἡ ὑπὸ ΔΕΒ γωνία τῆς ὑπὸ ΔΑΕ. ἴση δὲ ἡ ὑπὸ
ΔΑΕ τῇ  ὑπὸ ΔΒΕ; μείζων ἄρα ἡ ὑπὸ ΔΕΒ τῆς ὑπὸ
ΔΒΕ. ὑπὸ δὲ τὴν μείζονα γωνίαν ἡ μείζων πλευρὰ
ὑποτείνει; μείζων ἄρα ἡ ΔΒ τῆς ΔΕ. ἴση δὲ ἡ ΔΒ
τῇ ΔΖ. μείζων ἄρα ἡ ΔΖ τῆς ΔΕ ἡ ἐλάττων τῆς
μείζονος; ὅπερ ἐστὶν ἀδύνατον. οὐκ ἄρα ἡ ἀπὸ
τοῦ Α ἐπὶ τὸ Β ἐπιζευγνυμένη εὐθεῖα ἐκτὸς
πεσεῖται τοῦ κύκλου. ὁμοίως δὴ δείχομεν, ὅτι οὐδὲ ἐπ
α ὐτῆς τῆς περιφερείας; ἐντὸς ἄρα.}\\~\\~\\

\epsfysize=2.2in
\centerline{\epsffile{Book03/fig02g.eps}}

{\greekfont Ἐὰν ἄρα κύκλου ἐπὶ τῆς περιφερείας ληφθῇ δύο
τυθόντα σημεῖα, ἡ ἐπὶ τὰ σημεῖα ἐπιζευγνυμένη εὐθεῖα
ἐντὸς πεσεῖται τοῦ κύκλου; ὅπερ ἔδει δεῖχαι.}}

\ParallelRText{
\begin{center}
{\large Proposition 2}
\end{center}

If two points are taken at random on the circumference of a circle,  (then)
the straight-line joining the points will fall inside the circle.

Let $ABC$ be a circle, and let two points $A$ and $B$ be taken
at random on its circumference. I say that the straight-line joining
$A$ to $B$ will fall inside the circle.

For (if) not (then), if possible, let it fall outside (the circle), like $AEB$ (in the figure). And let the center of the circle $ABC$ be found [Prop.~3.1], and let it be (at point) $D$. And let $DA$ and $DB$ be joined, and let $DFE$ be
drawn through.

Therefore, since $DA$ is equal to $DB$, the angle $DAE$ (is) thus also equal to
$DBE$ [Prop.~1.5]. And since in triangle $DAE$ the one side, $AEB$, has been
produced, angle $DEB$ (is) thus greater than $DAE$ [Prop.~1.16]. And
$DAE$ (is) equal to $DBE$ [Prop.~1.5]. Thus, $DEB$ (is) greater than $DBE$.
And the greater angle is subtended by the greater side [Prop.~1.19].
Thus, $DB$ (is) greater than $DE$. And $DB$ (is) equal to $DF$. Thus,
$DF$ (is) greater than $DE$, the lesser than the greater. The very thing is impossible.
Thus, the straight-line joining $A$ to $B$ will not fall outside the
circle. So, similarly, we can show that neither (will it fall) on the
circumference itself. Thus, (it will fall) inside (the circle).

\epsfysize=2.2in
\centerline{\epsffile{Book03/fig02e.eps}}

Thus, if two points are taken at random on the circumference of a circle, (then)
the straight-line joining the points will fall inside the circle. (Which is) the
very thing it was required to show.}
\end{Parallel}

%%%%%%
% Prop 3.3
%%%%%%
\pdfbookmark[1]{Proposition 3.3}{pdf3.3}
\begin{Parallel}{}{} 
\ParallelLText{
\begin{center}
{\large \ggn{3}.}
\end{center}\vspace*{-7pt}

{\greekfont Ἐὰν ἐν κύκλῳ εὐθεῖά τις διὰ τοῦ κέντρου εὐθεῖάν τινα μὴ διὰ τοῦ κέντρου δίθα τέμνῃ, καὶ πρὸς ὀρθὰς α ὐτὴν τέμνει; καὶ ἐὰν πρὸς ὀρθὰς α ὐτὴν τέμνῃ, καὶ δίθα α ὐτὴν τέμνει.}

{\greekfont Ἔστω κύκλος ὁ ΑΒΓ, καὶ ἐν α ὐτῷ εὐθεῖά τις διὰ τοῦ κέντρου ἡ ΓΔ εὐθεῖάν τινα μὴ διὰ τοῦ κέντρου τὴν ΑΒ δίθα τεμνέτω κατὰ τὸ
Ζ σημεῖον; λέγω, ὅτι καὶ πρὸς ὀρθὰς α ὐτὴν τέμνει.}

{\greekfont Εἰλήφθω γὰρ τὸ κέντρον τοῦ ΑΒΓ κύκλου, καὶ ἔστω τὸ Ε, καὶ ἐπεζεύθθωσαν αἱ ΕΑ, ΕΒ.}

{\greekfont Καὶ ἐπεὶ ἴση ἐστὶν ἡ ΑΖ τῇ ΖΒ, κοινὴ δὲ ἡ ΖΕ, δύο δυσὶν ἴσαι [εἰσίν]; καὶ βάσις ἡ ΕΑ βάσει τῇ ΕΒ ἴση; γωνία ἄρα ἡ ὑπὸ ΑΖΕ γωνίᾳ τῇ ὑπὸ ΒΖΕ ἴση ἐστίν. ὅταν δὲ εὐθεῖα ἐπ εὐθεῖαν
σταθεῖσα τὰς ἐφεχῆς γωνίας ἴσας ἀλλήλαις ποιῇ,  ὀρθὴ
ἑκατέρα τῶν ἴσων γωνιῶν ἐστιν; ἑκατέρα ἄρα τῶν ὑπὸ ΑΖΕ, ΒΖΕ ὀρθή ἐστιν. ἡ ΓΔ ἄρα διὰ τοῦ κέντρου οὖσα τὴν ΑΒ μὴ
διὰ τοῦ κέντρου οὖσαν δίθα τέμνουσα καὶ πρὸς ὀρθὰς τέμνει.}

{\greekfont Ἀλλὰ δὴ ἡ ΓΔ τὴν ΑΒ πρὸς ὀρθὰς τεμνέτω; λέγω, ὅτι καὶ δίθα
α ὐτὴν τέμνει, τουτέστιν, ὅτι ἴση ἐστὶν ἡ ΑΖ τῇ ΖΒ.}

{\greekfont Τῶν γὰρ α ὐτῶν κατασκευασθέντων, ἐπεὶ ἴση ἐστὶν ἡ ΕΑ τῇ ΕΒ,
ἴση ἐστὶ καὶ γωνία ἡ ὑπὸ ΕΑΖ τῇ ὑπὸ ΕΒΖ. ἐστὶ δὲ καὶ
ὀρθὴ ἡ ὑπὸ ΑΖΕ ὀρθῇ τῇ ὑπὸ ΒΖΕ ἴση; δύο ἄρα τρίγωνά
ἐστι ΕΑΖ, ΕΖΒ τὰς δύο γωνίας δυσὶ γωνίαις ἴσας ἔθοντα καὶ
μίαν πλευρὰν μιᾷ πλευρᾷ ἴσην κοινὴν α ὐτῶν τὴν ΕΖ ὑποτείνουσαν ὑπὸ μίαν τῶν ἴσων γωνιῶν; καὶ τὰς λοιπὰς ἄρα πλευρὰς ταῖς
λοιπαῖς πλευραῖς ἴσας ἕχει; ἴση ἄρα ἡ ΑΖ τῇ ΖΒ.}\\~\\~\\~\\

\epsfysize=2.2in
\centerline{\epsffile{Book03/fig03g.eps}}

{\greekfont Ἐὰν ἄρα ἐν κύκλῳ εὐθεῖά τις διὰ τοῦ κέντρου εὐθεῖάν τινα μὴ διὰ τοῦ κέντρου δίθα τέμνῃ, καὶ πρὸς ὀρθὰς α ὐτὴν τέμνει; καὶ ἐὰν πρὸς ὀρθὰς α ὐτὴν τέμνῃ, καὶ δίθα α ὐτὴν τέμνει; ὅπερ ἔδει δεῖχαι.}}

\ParallelRText{
\begin{center}
{\large Proposition 3}
\end{center}

In a circle, if any straight-line through the center cuts in half any straight-line
not through the center, (then) it also cuts it at right-angles. And (conversely)
if it cuts it at right-angles, (then) it also cuts it in half.

Let $ABC$ be a circle, and, within it, let some straight-line through the center,
$CD$, 
cut in half some straight-line not through the center, $AB$, at the point $F$.
I say that ($CD$) also cuts ($AB$) at right-angles.

For let the center of the circle $ABC$ be found [Prop.~3.1], and let it
be (at point) $E$, and let $EA$ and $EB$ be joined.

And since  $AF$ is equal to $FB$, and $FE$ (is) common, two (sides  of triangle $AFE$) [are] equal to
two (sides of triangle $BFE$). And the base $EA$ (is) equal to the base $EB$. Thus, angle
$AFE$ is equal to angle $BFE$ [Prop.~1.8]. And when a straight-line
stood upon (another) straight-line makes adjacent angles (which are)
equal to one another, each of the equal angles is a right-angle [Def.~1.10].
Thus, $AFE$ and $BFE$ are each right-angles. Thus, the (straight-line) $CD$, which is
through the center
and cuts  in half the (straight-line) $AB$, which is not through the center, also cuts ($AB$) at right-angles.

And so let $CD$ cut $AB$ at right-angles. I say that it also cuts ($AB$) in half.
That is to say, that $AF$ is equal to $FB$.

For, with the same construction, since $EA$ is equal to $EB$, angle $EAF$ is also
equal to $EBF$ [Prop.~1.5]. And the right-angle $AFE$ is also equal to the
right-angle $BFE$. Thus, $EAF$ and $EFB$ are two triangles having two angles
equal to two angles, and one side equal to one side---(namely), their common (side)
$EF$, subtending one of the equal angles. Thus, they will also have the remaining sides equal to the (corresponding) remaining sides [Prop.~1.26]. Thus, $AF$ (is) equal to $FB$.

\epsfysize=2.2in
\centerline{\epsffile{Book03/fig03e.eps}}

Thus, in a circle, if any straight-line through the center cuts in half any straight-line
not through the center, (then) it also cuts it at right-angles. And (conversely)
if it cuts it at right-angles, (then) it also cuts it in half. (Which is) the very thing
it was required to show.}
\end{Parallel}

%%%%%%
% Prop 3.4
%%%%%%
\pdfbookmark[1]{Proposition 3.4}{pdf3.4}
\begin{Parallel}{}{} 
\ParallelLText{
\begin{center}
{\large \ggn{4}.}
\end{center}\vspace*{-7pt}

{\greekfont Ἐὰν ἐν κύκλῳ δύο εὐθεῖαι τέμνωσιν ἀλλήλας μὴ δὶα τοῦ
κέντρου οὖσαι, οὐ τέμνουσιν ἀλλήλας δίθα.}\\

\epsfysize=2in
\centerline{\epsffile{Book03/fig04g.eps}}

{\greekfont Ἔστω κύκλος ὁ ΑΒΓΔ, καὶ ἐν α ὐτῷ δύο εὐθεῖαι αἱ ΑΓ,
ΒΔ τεμνέτωσαν ἀλλήλας κατὰ τὸ Ε μὴ διὰ τοῦ κέντρου οὖσαι;
λέγω, ὅτι οὐ τέμνουσιν ἀλλήλας δίθα.}

{\greekfont Εἰ γὰρ δυνατόν, τεμνέτωσαν ἀλλήλας δίθα ὥστε ἴσην εἶναι
τὴν μὲν ΑΕ τῇ ΕΓ, τὴν δὲ ΒΕ τῇ ΕΔ; καὶ εἰλήφθω τὸ κέντρον
τοῦ ΑΒΓΔ κύκλου, καὶ ἔστω τὸ Ζ, καὶ ἐπεζεύθθω ἡ ΖΕ.}

{\greekfont Ἐπεὶ οὖν εὐθεῖά τις διὰ τοῦ κέντρου ἡ ΖΕ εὐθεῖάν τινα
μὴ διὰ τοῦ κέντρου τὴν ΑΓ δίθα τέμνει, καὶ πρὸς ὀρθὰς α ὐτὴν τέμνει; ὀρθὴ ἄρα ἐστὶν ἡ ὑπὸ ΖΕΑ; πάλιν, ἐπεὶ εὐθεῖά τις ἡ ΖΕ εὐθεῖάν τινα τὴν ΒΔ δίθα τέμνει, καὶ πρὸς ὀρθὰς α ὐτὴν τέμνει;
ὀρθὴ ἄρα ἡ ὑπὸ ΖΕΒ. ἐδείθθη δὲ καὶ ἡ ὑπὸ ΖΕΑ ὀρθή; ἴση ἄρα ἡ ὑπὸ ΖΕΑ τῇ ὑπὸ ΖΕΒ ἡ ἐλάττων τῇ μείζονι;
ὅπερ ἐστὶν ἀδύνατον. οὐκ ἄρα αἱ ΑΓ, ΒΔ τέμνουσιν ἀλλήλας δίθα.}

{\greekfont Ἐὰν ἄρα ἐν κύκλῳ δύο εὐθεῖαι τέμνωσιν ἀλλήλας μὴ δὶα τοῦ
κέντρου οὖσαι, οὐ τέμνουσιν ἀλλήλας δίθα; ὅπερ ἔδει δεῖχαι.}}

\ParallelRText{
\begin{center}
{\large Proposition 4}
\end{center}

In a circle, if two straight-lines, which are not through the center, cut
one another, (then) they do not cut one another in half.

\epsfysize=2in
\centerline{\epsffile{Book03/fig04e.eps}}

Let $ABCD$ be a circle, and within it, let two straight-lines, $AC$ and
$BD$, which are not through the center, cut one another at  (point) $E$. I say that they
do not cut one another in half.

For, if possible, let them cut one another in half, such that $AE$ is equal to $EC$, and
$BE$ to $ED$. And let the center of the circle $ABCD$ be found [Prop.~3.1], and let it be (at point) $F$, and let $FE$ be joined.

Therefore, since some straight-line through the center, $FE$, cuts in half
some straight-line not through the center, $AC$, it also cuts it at right-angles
[Prop.~3.3]. Thus, $FEA$ is a right-angle. Again, since some straight-line
$FE$ cuts in half some straight-line $BD$, it also cuts it at right-angles [Prop.~3.3]. Thus, $FEB$ (is) a right-angle. But $FEA$ was also shown (to be) a right-angle. Thus, $FEA$ (is) equal to $FEB$, the lesser to the greater.
The very thing is impossible. Thus, $AC$ and $BD$ do not cut one another in half.

Thus, in a circle, if two straight-lines, which are not through the center, cut
one another, (then) they do not cut one another in half. (Which is) the very
thing it was required to show.}
\end{Parallel}

%%%%%%
% Prop 3.5
%%%%%%
\pdfbookmark[1]{Proposition 3.5}{pdf3.5}
\begin{Parallel}{}{} 
\ParallelLText{
\begin{center}
{\large \ggn{5}.}
\end{center}\vspace*{-7pt}

{\greekfont Ἐὰν δύο κύκλοι τέμνωσιν ἀλλήλους, οὐκ ἔσται α ὐτῶν τὸ α ὐτὸ
κέντρον.}

\epsfysize=2in
\centerline{\epsffile{Book03/fig05g.eps}}

{\greekfont Δύο γὰρ κύκλοι οἱ ΑΒΓ, ΓΔΗ τεμνέτωσαν ἀλλήλους κατὰ τὰ Β, Γ σημεῖα. λέγω, ὅτι οὐκ ἔσται α ὐτῶν τὸ α ὐτὸ κέντρον.}

{\greekfont Εἰ γὰρ δυνατόν, ἔστω τὸ Ε, καὶ ἐπεζεύθθω ἡ ΕΓ, καὶ διήθθω
ἡ ΕΖΗ, ὡς ἔτυθεν. καὶ ἐπεὶ τὸ Ε σημεῖον κέντρον ἐστὶ τοῦ ΑΒΓ κύκλου, ἵση ἐστὶν ἡ ΕΓ τῇ ΕΖ. πάλιν, ἐπεὶ τὸ Ε σημεῖον κέντρον ἐστὶ τοῦ ΓΔΗ κύκλου, ἴση ἐστὶν ἡ ΕΓ τῇ ΕΗ; ἐδείθθη
δὲ ἡ ΕΓ καὶ τῇ ΕΖ ἴση; καὶ ἡ ΕΖ ἄρα τῇ ΕΗ ἐστιν ἴση ἡ
ἐλάσσων τῇ μείζονι; ὅπερ ἐστὶν ἀδύνατον.  οὐκ ἄρα τὸ Ε
σημεῖον κέντρον ἐστὶ τῶν ΑΒΓ, ΓΔΗ κύκλων.}

{\greekfont Ἐὰν ἄρα δύο κύκλοι τέμνωσιν ἀλλήλους, οὐκ ἔστιν α ὐτῶν τὸ
α ὐτὸ κέντρον; ὅπερ ἔδει δεῖχαι.}}

\ParallelRText{
\begin{center}
{\large Proposition 5}
\end{center}

If two circles cut one another,  (then) they will not have the same center.

\epsfysize=2in
\centerline{\epsffile{Book03/fig05e.eps}}

For let the two circles $ABC$ and $CDG$ cut one another at  points $B$ and $C$.
I say that they will not have the same center.

For, if possible, let $E$ be (the common center), and let $EC$ be joined,
and let $EFG$ be drawn through (the two circles), at random. And since point $E$
is the center of the circle $ABC$, $EC$ is equal to $EF$. Again, since point $E$ is
the center of the circle $CDG$, $EC$ is equal to $EG$. But $EC$  was also shown (to be) equal to $EF$. Thus, $EF$ is also equal to $EG$, the lesser to the
greater. The very thing is impossible. Thus, point $E$ is not the (common)
center of the circles $ABC$ and $CDG$.

Thus, if two circles cut one another, (then) they will not have the same
center. (Which is) the very thing it was required to show.}
\end{Parallel}

%%%%%%
% Prop 3.6
%%%%%%
\pdfbookmark[1]{Proposition 3.6}{pdf3.6}
\begin{Parallel}{}{} 
\ParallelLText{
\begin{center}
{\large \ggn{6}.}
\end{center}\vspace*{-7pt}

{\greekfont Ἐὰν δύο κύκλοι ἐφάπτωνται ἀλλήλων, οὐκ ἔσται α ὐτῶν τὸ α ὐτὸ κέντρον.}

{\greekfont Δύο γὰρ κύκλοι οἱ ΑΒΓ, ΓΔΕ ἐφαπτέσθωσαν ἀλλήλων κατὰ τὸ Γ
σημεῖον; λέγω, ὅτι οὐκ ἔσται α ὐτῶν τὸ α ὐτὸ κέντρον.}

{\greekfont Εἰ γὰρ δυνατόν, ἔστω τὸ Ζ, καὶ ἐπεζεύθθω ἡ ΖΓ, καὶ διήθθω, ὡς ἔτυθεν, ἡ ΖΕΒ.}

{\greekfont Ἐπεὶ οῦ̓ν τὸ Ζ σημεῖον κέντρον ἐστὶ τοῦ ΑΒΓ κύκλου, ἴση
ἐστὶν ἡ ΖΓ τῇ ΖΒ. πάλιν, ἐπεὶ τὸ Ζ σημεῖον κέντρον ἐστὶ τοῦ ΓΔΕ κύκλου, ἴση ἐστὶν ἡ ΖΓ τῇ ΖΕ. ἐδείθθη δὲ ἡ ΖΓ τῇ ΖΒ
ἴση; καὶ ἡ ΖΕ ἄρα τῇ ΖΒ ἐστιν ἴση, ἡ ἐλάττων τῇ μείζονι;
ὅπερ ἐστὶν ἀδύνατον. οὐκ ἄρα τὸ Ζ σημεῖον κέντρον ἐστὶ
τῶν ΑΒΓ, ΓΔΕ κύκλων.}

\epsfysize=2.2in
\centerline{\epsffile{Book03/fig06g.eps}}

{\greekfont Ἐὰν ἄρα δύο κύκλοι ἐφάπτωνται ἀλλήλων, οὐκ ἔσται α ὐτῶν τὸ α ὐτὸ κέντρον; ὅπερ ἔδει δεῖχαι.}}

\ParallelRText{
\begin{center}
{\large Proposition 6}
\end{center}

If  two circles touch one another, (then) they  will not have the same center.

For let the two circles $ABC$ and $CDE$ touch one another at point $C$. I say that
they will not have the same center.

For, if possible, let $F$ be (the common center), and let $FC$ be
joined, and let $FEB$ be drawn through (the two circles), at random.

Therefore, since point $F$ is the center of the circle $ABC$, $FC$ is equal to
$FB$. Again, since point $F$ is the center of the circle $CDE$, $FC$ is
equal to $FE$. But $FC$ was shown (to be) equal to $FB$. Thus, $FE$ is
also equal to $FB$, the lesser to the greater. The very thing is impossible.
Thus, point $F$ is not the (common) center of the circles
$ABC$ and $CDE$.

\epsfysize=2.2in
\centerline{\epsffile{Book03/fig06e.eps}}

Thus, if two circles touch one another, (then) they  will not have the same center. (Which is) the very thing it was required to show.}
\end{Parallel}

%%%%%%
% Prop 3.7
%%%%%%
\pdfbookmark[1]{Proposition 3.7}{pdf3.7}
\begin{Parallel}{}{} 
\ParallelLText{
\begin{center}
{\large \ggn{7}.}
\end{center}\vspace*{-7pt}

{\greekfont Ἐὰν κύκλου ἐπὶ τῆς διαμέτρου ληφθῇ τι σημεῖον, ὃ μή ἐστι κέντρον τοῦ κύκλου, ἀπὸ δὲ τοῦ σημείου πρὸς τὸν κύκλον προσπίπτωσιν εὐθεῖαί τινες, μεγίστη μὲν ἔσται, ἐφ ἧς
τὸ κέντρον, ἐλαθίστη δὲ ἡ λοιπή, τῶν δὲ ἄλλων ἀεὶ ἡ
ἔγγιον τῆς δὶα τοῦ κέντρου τῆς ἀπώτερον μείζων ἐστίν, δύο
δὲ μόνον ἴσαι ἀπὸ τοῦ σημείου προσπεσοῦνται πρὸς τὸν
κύκλον ἐφ ἑκάτερα τῆς ἐλαθίστης.}

{\greekfont Ἔστω κύκλος ὁ ΑΒΓΔ, διάμετρος δὲ α ὐτοῦ ἔστω ἡ ΑΔ, καὶ ἐπὶ
τῆς ΑΔ εἰλήφθω τι σημεῖον τὸ Ζ, ὃ μή ἐστι κέντρον τοῦ κύκλου, κέντρον δὲ τοῦ κύκλου ἔστω τὸ Ε, καὶ ἀπὸ τοῦ Ζ πρὸς τὸν ΑΒΓΔ κύκλον προσπιπτέτωσαν εὐθεῖαί τινες αἱ ΖΒ, ΖΓ, ΖΗ; λέγω, ὅτι
μεγίστη μέν ἐστιν ἡ ΖΑ, ἐλαθίστη δὲ ἡ ΖΔ, τῶν δὲ ἄλλων ἡ μὲν
ΖΒ τῆς ΖΓ μείζων, ἡ δὲ ΖΓ τῆς ΖΗ.}

{\greekfont Ἐπεζεύθθωσαν γὰρ αἱ ΒΕ, ΓΕ, ΗΕ. καὶ ἐπεὶ παντὸς τριγώνου αἱ δύο πλευραὶ τῆς λοιπῆς μείζονές εἰσιν, αἱ ἄρα ΕΒ, ΕΖ τῆς ΒΖ μείζονές
εἰσιν. ἴση δὲ ἡ ΑΕ τῇ ΒΕ [αἱ ἄρα ΒΕ, ΕΖ ἴσαι εἰσὶ τῇ ΑΖ];
μείζων ἄρα ἡ ΑΖ τῆς ΒΖ. πάλιν, ἐπεὶ ἴση ἐστὶν ἡ ΒΕ τῇ ΓΕ,
κοινὴ δὲ ἡ ΖΕ, δύο δὴ αἱ ΒΕ, ΕΖ δυσὶ ταῖς ΓΕ, ΕΖ ἴσαι εἰσίν.
ἀλλὰ καὶ γωνία ἡ ὑπὸ ΒΕΖ γωνίας τῆς ὑπὸ ΓΕΖ μείζων; βάσις
ἄρα ἡ ΒΖ βάσεως τῆς ΓΖ μείζων ἐστίν. διὰ τὰ α ὐτὰ δὴ καὶ
ἡ ΓΖ τῆς ΖΗ μείζων ἐστίν.}

{\greekfont Πάλιν, ἐπεὶ αἱ ΗΖ, ΖΕ τῆς ΕΗ μείζονές εἰσιν, ἴση δὲ ἡ ΕΗ
τῇ ΕΔ, αἱ ἄρα ΗΖ, ΖΕ τῆς ΕΔ μείζονές εἰσιν. κοινὴ ἀφῃρήσθω ἡ ΕΖ; λοιπὴ ἄρα ἡ ΗΖ λοιπῆς τῆς ΖΔ μείζων ἐστίν. μεγίστη μὲν ἄρα ἡ ΖΑ, ἐλαθίστη δὲ ἡ ΖΔ, μείζων δὲ ἡ μὲν ΖΒ τῆς ΖΓ, ἡ
δὲ ΖΓ τῆς ΖΗ.}\\~\\~\\~\\~\\~\\~\\

\epsfysize=2.5in
\centerline{\epsffile{Book03/fig07g.eps}}

{\greekfont Λέγω, ὅτι καὶ ἀπὸ τοῦ Ζ σημείου δύο μόνον ἴσαι προσπεσοῦνται πρὸς τὸν ΑΒΓΔ κύκλον ἐφ ἑκάτερα τῆς ΖΔ ἐλαθίστης. συνεστάτω γὰρ πρὸς τῇ ΕΖ εὐθείᾳ καὶ τῷ πρὸς α ὐτῇ σημείῳ τῷ Ε τῇ
ὑπὸ ΗΕΖ γωνίᾳ ἴση ἡ ὑπὸ ΖΕΘ, καὶ ἐπεζεύθθω ἡ ΖΘ. ἐπεὶ οὖν ἴση ἐστὶν ἡ ΗΕ τῇ ΕΘ, κοινὴ δὲ ἡ ΕΖ, δύο
δὴ αἱ ΗΕ, ΕΖ δυσὶ ταῖς ΘΕ, ΕΖ ἴσαι εἰσίν; καὶ γωνία ἡ ὑπὸ ΗΕΖ
γωνίᾳ τῇ ὑπὸ ΘΕΖ ἴση; βάσις ἄρα ἡ ΖΗ βάσει τῇ ΖΘ ἴση
ἐστίν. λέγω δή, ὅτι τῇ ΖΗ ἄλλη ἴση οὐ προσπεσεῖται πρὸς τὸν
κύκλον ἀπὸ τοῦ Ζ σημείου. εἰ γὰρ δυνατόν, προσπιπτέτω ἡ ΖΚ.
καὶ ἐπεὶ ἡ ΖΚ τῇ ΖΗ ἴση ἐστίν, ἀλλὰ ἡ ΖΘ τῇ ΖΗ
[ἴση ἐστίν], καὶ  ἡ  ΖΚ
ἄρα τῇ ΖΘ ἐστιν ἴση, ἡ ἔγγιον
τῆς διὰ τοῦ κέντρου τῇ ἀπώτερον ἴση; ὅπερ ἀδύνατον.
οὐκ ἄρα ἀπὸ τοῦ Ζ σημείου ἑτέρα τις προσπεσεῖται πρὸς τὸν
κύκλον ἴση τῇ ΗΖ; μία ἄρα μόνη.}

{\greekfont Ἐὰν ἄρα κύκλου ἐπὶ τῆς διαμέτρου ληφθῇ τι σημεῖον, ὃ μή ἐστι κέντρον τοῦ κύκλου, ἀπὸ δὲ τοῦ σημείου πρὸς τὸν κύκλον προσπίπτωσιν εὐθεῖαί τινες, μεγίστη μὲν ἔσται, ἐφ ἧς
τὸ κέντρον, ἐλαθίστη δὲ ἡ λοιπή, τῶν δὲ ἄλλων ἀεὶ ἡ
ἔγγιον τῆς δὶα τοῦ κέντρου τῆς ἀπώτερον μείζων ἐστίν, δύο
δὲ μόνον ἴσαι ἀπὸ τοῦ α ὐτοῦ σημείου προσπεσοῦνται πρὸς τὸν
κύκλον ἐφ ἑκάτερα τῆς ἐλαθίστης; ὅπερ ἔδει δεῖχαι.}}

\ParallelRText{
\begin{center}
{\large Proposition 7}
\end{center}

If some point, which is not the
center of the circle, is taken on the diameter of a circle,
and some straight-lines radiate from the point towards
the (circumference of the) circle, (then) the greatest (straight-line) will be that on which
the center (lies), and the least the remainder (of the same diameter). And
for the others, a (straight-line) nearer$^\dag$ to the (straight-line) through the center is always
greater than a (straight-line) further away. And only two equal (straight-lines)
will radiate from the point towards the (circumference of the) circle, (one)
on each (side) of the least (straight-line).

Let $ABCD$ be a circle, and let $AD$ be its diameter, and
let some point $F$, which is not the center of the circle, be taken on $AD$.
Let $E$ be the center of the circle. And let some straight-lines, $FB$, $FC$,
and $FG$, radiate from $F$ towards (the circumference of) circle $ABCD$.
I say that $FA$ is the greatest (straight-line), $FD$ the least, and of the others, $FB$ (is)
greater than $FC$, and $FC$ than $FG$.

For let $BE$, $CE$, and $GE$ be joined. And since for every triangle
(any) two sides are greater than the remaining (side) [Prop.~1.20],
$EB$ and $EF$ is thus greater than $BF$. And $AE$ (is) equal to $BE$ [thus,
$BE$ and $EF$ is equal to $AF$]. Thus, $AF$ (is) greater than $BF$. Again,
since $BE$ is equal to $CE$, and $FE$ (is) common, the two (straight-lines)
$BE$, $EF$ are equal to the two (straight-lines) $CE$, $EF$ (respectively). But, angle $BEF$ (is)
also greater than angle $CEF$.$^\ddag$ Thus, the base $BF$ is greater than the base $CF$.
Thus, the base $BF$ is greater than the base $CF$ [Prop.~1.24]. So, for the same (reasons), $CF$ is also greater than $FG$.

Again, since $GF$ and $FE$ are greater than $EG$ [Prop.~1.20], and $EG$
(is) equal to $ED$, $GF$ and $FE$ are thus greater than $ED$. Let $EF$ be
taken from both. Thus, the remainder $GF$ is greater than the remainder $FD$.
Thus, $FA$ (is) the greatest  (straight-line), $FD$ the least, and $FB$ (is) greater than $FC$, and $FC$ than $FG$.

\epsfysize=2.5in
\centerline{\epsffile{Book03/fig07e.eps}}

I also say that from point $F$ only two equal (straight-lines) will radiate
towards (the circumference of) circle $ABCD$, (one) on each (side) of the
least (straight-line) $FD$. For let the (angle) $FEH$, equal to angle $GEF$, be constructed on the straight-line $EF$, at the point $E$ on it [Prop.~1.23], and
let $FH$ be joined. Therefore, since $GE$ is equal to $EH$, and $EF$ (is) common, the two (straight-lines) $GE$, $EF$ are equal to the two (straight-lines)
$HE$, $EF$ (respectively). And angle $GEF$ (is) equal to angle $HEF$. Thus, the base $FG$ is
equal to the base $FH$ [Prop.~1.4]. So I say that another (straight-line)
equal to $FG$
will not radiate towards (the circumference of)
the circle from point $F$. For, if possible, let $FK$ (so) radiate. And since
$FK$ is equal to $FG$, but $FH$ [is equal] to $FG$, $FK$ is thus also equal to
$FH$, the nearer to the (straight-line) through the center equal to the
further away. The very thing (is) impossible. Thus, another (straight-line)
equal to $GF$ will not radiate from the point $F$ towards (the circumference of) the circle.
Thus, (there is) only one (such straight-line).

Thus, if some point, which is not the
center of the circle, is taken on the diameter of a circle,
and some straight-lines radiate from the point towards
the (circumference of the) circle, (then) the greatest (straight-line) will be that on which
the center (lies), and the least the remainder (of the same diameter). And
for the others, a (straight-line) nearer to the (straight-line) through the center is always
greater than a (straight-line) further away. And only two equal (straight-lines)
will radiate from the same point towards the (circumference of the) circle,  (one)
on each (side) of the least (straight-line). (Which is) the very thing it was required to show.}
\end{Parallel}
{\footnotesize \noindent$^\dag$ Presumably, in an angular sense.\\
$^\ddag$ This is not proved, except by reference to the figure.}

%%%%%%
% Prop 3.8
%%%%%%
\pdfbookmark[1]{Proposition 3.8}{pdf3.8}
\begin{Parallel}{}{} 
\ParallelLText{
\begin{center}
{\large \ggn{8}.}
\end{center}\vspace*{-7pt}

{\greekfont Ἐὰν κύκλου ληφθῇ τι σημεῖον ἐκτός, ἀπὸ δὲ τοῦ σημείου πρὸς
τὸν κύκλον διαθθῶσιν εὐθεῖαί τινες, ὧν μία μὲν διὰ τοῦ κέντρου,
αἱ δὲ λοιπαί, ὡς ἔτυθεν, τῶν μὲν πρὸς τὴν κοίλην περιφέρειαν προσπιπτουσῶν εὐθειῶν μεγίστη μέν ἐστιν ἡ διὰ τοῦ κέντρου, τῶν δὲ ἄλλων ἀεὶ ἡ ἔγγιον τῆς διὰ τοῦ κέντρου τῆς ἀπώτερον μείζων ἐστίν, τῶν δὲ πρὸς τὴν κυρτὴν περιφέρειαν προσπιπτουσῶν εὐθειῶν
ἐλαθίστη μέν ἐστιν ἡ μεταχὺ τοῦ τε σημείου καὶ τῆς διαμέτρου, τῶν δὲ ἄλλων ἀεὶ ἡ ἔγγιον τῆς ἐλαθίστης τῆς ἀπώτερόν ἐστιν
ἐλάττων, δύο δὲ μόνον ἴσαι ἀπὸ τοῦ σημείου προσπεσοῦνται
πρὸς τὸν κύκλον ἐφ ἑκάτερα τῆς ἐλαθίστης.}

{\greekfont Ἔστω κύκλος ὁ ΑΒΓ, καὶ τοῦ ΑΒΓ εἰλήφθω τι σημεῖον ἐκτὸς τὸ Δ,
καὶ ἀπ α ὐτοῦ διήθθωσαν εὐθεῖαί τινες αἱ ΔΑ, ΔΕ, ΔΖ, ΔΓ,
ἔστω δὲ ἡ ΔΑ διὰ τοῦ κέντρου. λέγω, ὅτι τῶν μὲν πρὸς
τὴν ΑΕΖΓ κοίλην περιφέρειαν προσπιπτουσῶν εὐθειῶν μεγίστη
μέν ἐστιν ἡ διὰ τοῦ κέντρου ἡ ΔΑ, μείζων δὲ ἡ μὲν ΔΕ τῆς ΔΖ ἡ δὲ ΔΖ τῆς ΔΓ, τῶν δὲ πρὸς τὴν ΘΛΚΗ κυρτὴν περιφέρειαν προσπιπτουσῶν εὐθειῶν ἐλαθίστη μέν ἐστιν ἡ ΔΗ ἡ
μεταχὺ τοῦ σημείου καὶ τῆς διαμέτρου τῆς ΑΗ, ἀεὶ δὲ ἡ
ἔγγιον τῆς ΔΗ ἐλαθίστης ἐλάττων ἐστὶ τῆς ἀπώτερον, ἡ μὲν
ΔΚ τῆς ΔΛ, ἡ δὲ ΔΛ τῆς ΔΘ.}

{\greekfont Εἰλήφθω γὰρ τὸ κέντρον τοῦ ΑΒΓ κύκλου καὶ ἔστω τὸ Μ; καὶ
ἐπεζεύθθωσαν αἱ ΜΕ, ΜΖ, ΜΓ, ΜΚ, ΜΛ, ΜΘ.}

{\greekfont Καὶ ἐπεὶ ἴση ἐστὶν ἡ ΑΜ τῇ ΕΜ, κοινὴ προσκείσθω ἡ ΜΔ;
ἡ ἄρα ΑΔ ἴση ἐστὶ ταῖς ΕΜ, ΜΔ. ἀλλ αἱ ΕΜ, ΜΔ τῆς ΕΔ μείζονές εἰσιν; καὶ ἡ ΑΔ ἄρα τῆς ΕΔ μείζων ἐστίν. πάλιν, ἐπεὶ
ἴση ἐστὶν ἡ ΜΕ τῇ ΜΖ, κοινὴ δὲ ἡ ΜΔ, αἱ ΕΜ, ΜΔ ἄρα
ταῖς ΖΜ, ΜΔ ἴσαι εἰσίν; καὶ γωνία ἡ ὑπὸ ΕΜΔ γωνίας τῆς ὑπὸ ΖΜΔ μείζων ἐστίν. βάσις ἄρα ἡ ΕΔ βάσεως τῆς ΖΔ μείζων ἐστίν.
ὁμοίως δὴ δείχομεν, ὅτι καὶ ἡ ΖΔ τῆς
ΓΔ μείζων ἐστίν; μεγίστη μὲν ἄρα ἡ ΔΑ, μείζων δὲ ἡ μὲν ΔΕ
τῆς ΔΖ, ἡ δὲ ΔΖ τῆς ΔΓ.}

{\greekfont Καὶ ἐπεὶ αἱ ΜΚ, ΚΔ τῆς ΜΔ μείζονές εἰσιν, ἴση δὲ ἡ μ η τῇ ΜΚ, λοιπὴ ἄρα ἡ ΚΔ λοιπῆς τῆς ΗΔ μείζων ἐστίν; ὥστε
ἡ ΗΔ τῆς ΚΔ ἐλάττων ἐστίν; καὶ ἐπεὶ τριγώνου τοῦ ΜΛΔ
ἐπὶ μιᾶς τῶν πλευρῶν τῆς ΜΔ δύο εὐθεῖαι ἐντὸς
συνεστάθησαν αἱ ΜΚ, ΚΔ, αἱ ἄρα ΜΚ, ΚΔ  τῶν ΜΛ, ΛΔ ἐλάττονές εἰσιν; ἴση δὲ ἡ ΜΚ  
 τῇ ΜΛ; λοιπὴ ἄρα ἡ ΔΚ λοιπῆς τῆς ΔΛ ἐλάττων ἐστίν. ὁμοίως δὴ δείχομεν, ὅτι καὶ ἡ ΔΛ τῆς ΔΘ
ἐλάττων ἐστίν; ἐλαθίστη μὲν ἄρα ἡ ΔΗ, ἐλάττων δὲ ἡ
μὲν ΔΚ τῆς ΔΛ ἡ δὲ ΔΛ τῆς ΔΘ.}\\~\\~\\~\\~\\~\\~\\~\\

\epsfysize=2.5in
\centerline{\epsffile{Book03/fig08g.eps}}

{\greekfont Λέγω, ὅτι καὶ δύο μόνον ἴσαι ἀπὸ τοῦ Δ σημείου προσπεσοῦνται
πρὸς τὸν κύκλον ἐφ ἑκάτερα τῆς ΔΗ ἐλαθίστης; συνεστάτω πρὸς
τῇ ΜΔ εὐθείᾳ καὶ τῷ πρὸς α ὐτῇ σημείῳ τῷ Μ τῇ ὑπὸ ΚΜΔ γωνίᾳ ἴση γωνία ἡ ὑπὸ ΔΜΒ, καὶ ἐπεζεύθθω ἡ ΔΒ. καὶ ἐπεὶ ἴση ἐστὶν ἡ ΜΚ τῇ ΜΒ, κοινὴ δὲ ἡ ΜΔ, δύο δὴ αἱ ΚΜ, ΜΔ
δύο ταῖς ΒΜ, ΜΔ ἴσαι εἰσὶν ἑκατέρα ἑκατέρᾳ; καὶ γωνία ἡ
ὑπὸ ΚΜΔ γωνίᾳ τῇ ὑπὸ ΒΜΔ ἴση; βάσις  ἄρα ἡ ΔΚ βάσει
τῇ ΔΒ ἴση ἐστίν. λέγω [δή], ὅτι τῇ ΔΚ εὐθείᾳ ἄλλη ἴση
οὐ προσπεσεῖται πρὸς τὸν κύκλον ἀπὸ τοῦ Δ σημείου. εἰ γὰρ
δυνατόν, προσπιπτέτω καὶ ἔστω ἡ ΔΝ. ἐπεὶ οὖν ἡ ΔΚ
τῇ ΔΝ ἐστιν ἴση, ἀλλ ἡ ΔΚ τῇ ΔΒ ἐστιν ἴση, καὶ
ἡ ΔΒ ἄρα τῇ ΔΝ ἐστιν ἴση, ἡ ἔγγιον τῆς ΔΗ ἐλαθίστης τῇ
ἀπώτερον [ἐστιν] ἴση; ὅπερ ἀδύνατον ἐδείθθη. οὐκ
ἄρα πλείους ἢ δύο ἴσαι πρὸς τὸν ΑΒΓ κύκλον ἀπὸ τοῦ Δ σημείου
ἐφ ἑκάτερα τῆς ΔΗ ἐλαθίστης προσπεσοῦνται.}

{\greekfont Ἐὰν ἄρα κύκλου ληφθῇ τι σημεῖον ἐκτός, ἀπὸ δὲ τοῦ σημείου πρὸς
τὸν κύκλον διαθθῶσιν εὐθεῖαί τινες, ὧν μία μὲν διὰ τοῦ κέντρου
αἱ δὲ λοιπαί, ὡς ἔτυθεν, τῶν μὲν πρὸς τὴν κοίλην περιφέρειαν προσπιπτουσῶν εὐθειῶν μεγίστη μέν ἐστιν ἡ διὰ τοῦ κέντου, τῶν δὲ ἄλλων ἀεὶ ἡ ἔγγιον τῆς διὰ τοῦ κέντρου τῆς ἀπώτερον μείζων ἐστίν, τῶν δὲ πρὸς τὴν κυρτὴν περιφέρειαν προσπιπτουσῶν εὐθειῶν
ἐλαθίστη μέν ἐστιν ἡ μεταχὺ τοῦ τε σημείου καὶ τῆς διαμέτρου, τῶν δὲ ἄλλων ἀεὶ ἡ ἔγγιον τῆς ἐλαθίστης τῆς ἀπώτερόν ἐστιν
ἐλάττων, δύο δὲ μόνον ἴσαι ἀπὸ τοῦ σημείου προσπεσοῦνται
πρὸς τὸν κύκλον ἐφ ἑκάτερα τῆς ἐλαθίστης; ὅπερ ἔδει δεῖχαι.}}

\ParallelRText{
\begin{center}
{\large Proposition 8}
\end{center}

If some point is taken outside a circle, and some straight-lines are
drawn from the point to the (circumference of the) circle, one of which (passes) through the
center, the remainder  (being) random, (then) for the straight-lines
radiating towards the concave (part of the) circumference, the
greatest is  that (passing) through the center. For the others, a (straight-line)
nearer$^\dag$ to the (straight-line) through the center is always greater than one
further away. For the straight-lines radiating towards the convex (part of
the) circumference, the least is that between the point and the
diameter. For the others, a (straight-line) nearer to the least
(straight-line) is always less than one further away. And only
two equal (straight-lines) will radiate from the point towards the (circumference of the) circle, (one) on each (side)
of the least (straight-line).

Let $ABC$ be a circle, and let some point $D$ be taken outside $ABC$,
and from it let some straight-lines, $DA$, $DE$, $DF$, and $DC$, be
drawn through (the circle), and let $DA$ be through the center. I say 
that for the straight-lines radiating towards the concave (part of the)
circumference, $AEFC$, the greatest is the one (passing) through the center,  (namely) $AD$,
and (that) $DE$ (is) greater than $DF$, and $DF$ than $DC$. For the straight-lines
radiating towards the convex (part of the) circumference, $HLKG$, the
least is the one between the point and the diameter $AG$, (namely) $DG$, 
and  a (straight-line) nearer to the least (straight-line) $DG$ is
always less than one farther away, (so that) $DK$ (is less) than $DL$, and $DL$ than than
$DH$.

For let the center of the circle be found [Prop.~3.1], and let it be (at point) $M$
[Prop.~3.1]. And let $ME$, $MF$, $MC$, $MK$, $ML$, and $m\kern -.7pt h$ be joined.

And since $AM$ is equal to $EM$, let $MD$ be added to both. Thus, $AD$ is equal to $EM$ and $MD$. But, $EM$ and $MD$ is greater than
$ED$ [Prop.~1.20]. Thus, $AD$ is also greater than $ED$. Again, since $ME$ is
equal to $MF$, and $MD$ (is) common, the (straight-lines) $EM$, $MD$ are thus
equal to $FM$,
$MD$. And angle $EMD$ is greater than angle $FMD$.$^\ddag$
Thus, the
base $ED$ is greater than the base $FD$ [Prop.~1.24]. 
So, similarly,
we can show that $FD$ is also greater than $CD$. Thus, $AD$ (is) the greatest
(straight-line), and $DE$ (is) greater than $DF$, and $DF$ than $DC$.

And since $MK$ and $KD$ is greater than $MD$ [Prop. 1.20], and $MG$ (is)
equal to $MK$, the remainder $KD$ is thus greater than the remainder $GD$.
So $GD$ is less than $KD$. And since in triangle $MLD$, the two internal straight-lines
$MK$ and $KD$ were constructed on one of the sides, $MD$, (then) $MK$ and $KD$
are thus less than $ML$ and $LD$  [Prop.~1.21]. And $MK$ (is) equal to $ML$.
Thus, the remainder $DK$ is less than the remainder $DL$. So, similarly, we
can show that $DL$ is also less than $DH$. Thus, $DG$ (is) the least (straight-line),
and $DK$ (is) less than $DL$, and $DL$ than $DH$.

\epsfysize=2.5in
\centerline{\epsffile{Book03/fig08e.eps}}

I also say that only two equal (straight-lines) will radiate from point $D$ towards (the
circumference of) the circle, (one) on each (side) on the least (straight-line),
$DG$. Let the  angle $DMB$, equal to angle $KMD$, be constructed
on the straight-line $MD$, at the point $M$ on it [Prop.~1.23], and  let $DB$ be
joined. And since $MK$ is equal to $MB$, and $MD$ (is) common, the
two (straight-lines) $KM$, $MD$ are equal to the two (straight-lines) $BM$, $MD$,
respectively. And angle $KMD$ (is) equal to angle $BMD$. Thus, the base
$DK$ is equal to the base $DB$ [Prop.~1.4]. [So] I say that another (straight-line)
equal to $DK$ will not radiate towards the (circumference of the) circle
from point $D$. For, if possible, let (such a straight-line) radiate, and let
it be $DN$. Therefore, since $DK$ is equal to $DN$, but $DK$ is equal to $DB$, (then) $DB$
is thus also equal to $DN$, (so that) a (straight-line) nearer to the least (straight-line) $DG$ [is]
equal to one further away. The very thing was shown (to be)  impossible. 
Thus, not more than two equal (straight-lines) will radiate towards
(the circumference of) circle $ABC$ from point $D$, (one) on each side
of the least (straight-line) $DG$.

Thus, if some point is taken outside a circle, and some straight-lines are
drawn from the point to the (circumference of the) circle, one of which (passes) through the
center, the remainder  (being) random, (then) for the straight-lines
radiating towards the concave (part of the) circumference, the
greatest is  that (passing) through the center. For the others, a (straight-line)
nearer to the (straight-line) through the center is always greater than one
further away. For the straight-lines radiating towards the convex (part of
the) circumference, the least is that between the point and the
diameter. For the others, a (straight-line) nearer to the least
(straight-line) is always less than one further away. And only
two equal (straight-lines) will radiate from the point towards the (circumference of the) circle, (one) on each (side)
of the least (straight-line).  (Which is) the very thing it was required to show.}
\end{Parallel}
{\footnotesize \noindent$^\dag$ Presumably, in an angular sense.\\
$\ddag$  This is not
proved, except by reference to the figure.} 

%%%%%%
% Prop 3.9
%%%%%%
\pdfbookmark[1]{Proposition 3.9}{pdf3.9}
\begin{Parallel}{}{} 
\ParallelLText{
\begin{center}
{\large \ggn{9}.}
\end{center}\vspace*{-7pt}

{\greekfont Ἐὰν κύκλου ληφθῇ τι σημεῖον ἐντός, ἀπο δὲ τοῦ σημείου πρὸς
τὸν κύκλον προσπίπτωσι πλείους ἢ δύο ἴσαι εὐθεῖαι, τὸ ληφθὲν σημεῖον κέντρον ἐστὶ τοῦ κύκλου.}

{\greekfont Ἔστω κύκλος ὁ ΑΒΓ, ἐντὸς δὲ α ὐτοῦ σημεῖον τὸ Δ, καὶ ἀπὸ
τοῦ Δ πρὸς τὸν ΑΒΓ κύκλον προσπιπτέτωσαν πλείους ἢ δύο ἴσαι
εὐθεῖαι αἱ ΔΑ, ΔΒ, ΔΓ; λέγω, ὅτι τὸ Δ σημεῖον κέντρον ἐστὶ τοῦ
ΑΒΓ κύκλου.}\\

\epsfysize=2.2in
\centerline{\epsffile{Book03/fig09g.eps}}

{\greekfont Ἐπεζεύθθωσαν γὰρ αἱ ΑΒ, ΒΓ καὶ τετμήσθωσαν δίθα κατὰ τὰ Ε, Ζ σημεῖα, καὶ ἐπιζευθθεῖσαι αἱ ΕΔ, ΖΔ διήθθωσαν ἐπὶ τὰ
Η, Κ, Θ, Λ σημεῖα.}

{\greekfont Ἐπεὶ οὖν ἴση ἐστὶν ἡ ΑΕ τῇ ΕΒ, κοινὴ δὲ ἡ ΕΔ,
δύο δὴ αἱ ΑΕ, ΕΔ δύο ταῖς ΒΕ, ΕΔ ἴσαι εἰσίν; καὶ βάσις ἡ
ΔΑ βάσει τῇ ΔΒ ἴση; γωνία ἄρα ἡ ὑπὸ ΑΕΔ γωνίᾳ
τῇ ὑπὸ ΒΕΔ ἴση ἐστίν; ὀρθὴ ἄρα ἑκατέρα τῶν ὑπὸ ΑΕΔ,
ΒΕΔ γωνιῶν; ἡ ΗΚ ἄρα τὴν ΑΒ τέμνει δίθα καὶ πρὸς ὀρθάς. καὶ
ἐπεί, ἐὰν ἐν κύκλῳ εὐθεῖά τις εὐθεῖάν τινα δίθα τε καὶ πρὸς
ὀρθὰς τέμνῃ, ἐπὶ τῆς τεμνούσης ἐστὶ τὸ κέντρον τοῦ κύκλου,
ἐπὶ τῆς ΗΚ ἄρα ἐστὶ τὸ κέντρον τοῦ κύκλου. διὰ τὰ
α ὐτὰ δὴ καὶ ἐπὶ τῆς ΘΛ ἐστι τὸ κέντρον τοῦ ΑΒΓ κύκλου.
καὶ οὐδὲν ἕτερον κοινὸν ἔθουσιν αἱ ΗΚ, ΘΛ εὐθεῖαι ἢ τὸ Δ
σημεῖον; τὸ Δ ἄρα σημεῖον κέντρον ἐστὶ τοῦ ΑΒΓ
κύκλου.}

{\greekfont Ἐὰν ἄρα κύκλου ληφθῇ τι σημεῖον ἐντός, ἀπὸ δὲ τοῦ σημείου πρὸς
τὸν κύκλον προσπίπτωσι πλείους ἢ δύο ἴσαι εὐθεῖαι, τὸ ληφθὲν σημεῖον κέντρον ἐστὶ τοῦ κύκλου; ὅπερ ἔδει δεῖχαι.}}

\ParallelRText{
\begin{center}
{\large Proposition 9}
\end{center}

If  some point is taken inside a circle, and more than two equal straight-lines
radiate from the point towards the (circumference of the) circle, (then) the
point taken  is the center of the circle.

Let $ABC$ be a circle, and $D$ a point inside it, and let more than two equal
straight-lines, $DA$, $DB$, and $DC$, radiate from $D$ towards (the circumference of) circle
$ABC$. I say that point $D$ is the center of circle $ABC$.

\epsfysize=2.2in
\centerline{\epsffile{Book03/fig09e.eps}}

For let $AB$ and $BC$ be joined, and (then) be cut in half at points $E$ and $F$ 
(respectively) [Prop.~1.10]. And $ED$ and $FD$ being joined, let them be
drawn through to points $G$, $K$, $H$, and $L$.

Therefore, since $AE$ is equal to $EB$, and $ED$ (is) common, the two (straight-lines) $AE$, $ED$ are equal to the two (straight-lines) $BE$, $ED$ (respectively). And the 
base $DA$ (is) equal to the base $DB$. Thus, angle $AED$ is equal to
angle $BED$ [Prop.~1.8]. Thus, angles $AED$ and $BED$ (are) each right-angles [Def.~1.10]. Thus, $GK$ cuts $AB$ in half, and at right-angles. And since,
if some straight-line in a circle cuts some (other) straight-line in half,
and at right-angles, (then) the center of the circle is on the former (straight-line)
[Prop.~3.1~corr.], the center of the circle is thus on $GK$. So, for the
same (reasons), the center of circle $ABC$ is also on $HL$. And the straight-lines
$GK$ and $HL$ have no  common (point) other than point $D$. Thus, point
$D$ is the center of circle $ABC$.

Thus, if some point is taken inside a circle, and more than two equal straight-lines
radiate from the point towards the (circumference of the) circle, (then) the
point taken  is the center of the circle. (Which is) the very thing it was
required to show.}
\end{Parallel}

%%%%%%
% Prop 3.10
%%%%%%
\pdfbookmark[1]{Proposition 3.10}{pdf3.10}
\begin{Parallel}{}{} 
\ParallelLText{
\begin{center}
{\large \ggn{10}.}
\end{center}\vspace*{-7pt}

{\greekfont Κύκλος  κύκλον οὐ τέμνει κατὰ πλείονα σημεῖα ἢ δύο.}

{\greekfont Εἰ γὰρ δυνατόν, κύκλος ὁ ΑΒΓ κύκλον τὸν ΔΕΖ τεμνέτω κατὰ πλείονα
σημεῖα ἢ δύο τὰ Β, Η, Ζ, Θ, καὶ ἐπιζευθθεῖσαι αἱ ΒΘ, ΒΗ δίθα τεμνέσθωσαν κατὰ τὰ Κ, Λ σημεῖα; καὶ ἀπὸ τῶν Κ, Λ ταῖς ΒΘ, ΒΗ
πρὸς ὀρθὰς ἀθθεῖσαι αἱ ΚΓ, ΛΜ διήθθωσαν ἐπὶ τὰ Α, Ε σημεῖα.}\\~\\

\epsfysize=2.2in
\centerline{\epsffile{Book03/fig10g.eps}}

{\greekfont Ἐπεὶ οὖν ἐν κύκλῳ τῷ ΑΒΓ εὐθεῖά τις ἡ ΑΓ εὐθεῖάν τινα
τὴν ΒΘ δίθα καὶ πρὸς ὀρθὰς τέμνει, ἐπὶ τῆς ΑΓ ἄρα ἐστὶ τὸ κέντρον τοῦ ΑΒΓ κύκλου. πάλιν, ἐπεὶ ἐν κύκλῳ τῷ α ὐτῷ
τῷ ΑΒΓ εὐθεῖά τις ἡ ΝΧ εὐθεῖάν τινα τὴν ΒΗ δίθα καὶ πρὸς
ὀρθὰς τέμνει, ἐπὶ τῆς ΝΧ ἄρα ἐστὶ τὸ κέντρον τοῦ ΑΒΓ κύκλου.
ἐδείθθη δὲ καὶ ἐπὶ τῆς ΑΓ, καὶ κατ οὐδὲν συμβάλλουσιν
αἱ ΑΓ, ΝΧ εὐθεῖαι ἢ κατὰ τὸ Ο; τὸ Ο ἄρα σημεῖον κέντρον
ἐστὶ τοῦ ΑΒΓ κύκλου. ὁμοίως δὴ δείχομεν, ὅτι καὶ τοῦ ΔΕΖ κύκλου κέντρον ἐστὶ τὸ Ο; δύο ἄρα κύκλων τεμνόντων ἀλλήλους
τῶν ΑΒΓ, ΔΕΖ τὸ α ὐτό ἐστι κέντρον τὸ Ο; 
ὅπερ ἐστὶν ἀδύνατον.}

{\greekfont Οὐκ ἄρα κύκλος  κύκλον τέμνει κατὰ πλείονα σημεῖα ἢ δύο;
ὅπερ ἔδει δεῖχαι.}}

\ParallelRText{
\begin{center}
{\large Proposition 10}
\end{center}

A circle does not cut a(nother) circle at more than two points.

For, if possible, let the circle $ABC$ cut the circle $DEF$ at more than
two points, $B$, $G$, $F$, and $H$. And $BH$ and $BG$ being joined,
let them (then) be cut in half at points $K$ and $L$ (respectively). And $KC$ and $LM$
being drawn at right-angles to $BH$ and $BG$ from $K$ and $L$ (respectively) [Prop.~1.11], let them (then) be
drawn through to points $A$ and $E$ (respectively).

\epsfysize=2.2in
\centerline{\epsffile{Book03/fig10e.eps}}

Therefore, since in circle $ABC$ some straight-line $AC$ cuts some (other)
straight-line $BH$ in half, and at right-angles, the center of circle $ABC$
is thus on $AC$ [Prop.~3.1~corr.]. Again, since in the same circle $ABC$
some straight-line $NO$ cuts some (other straight-line) $BG$ in half, and
at right-angles, the center of circle $ABC$ is thus on $NO$ [Prop.~3.1~corr.].
And it was also shown (to be) on $AC$. And the straight-lines $AC$ and $NO$
meet at no other (point) than $P$. Thus, point $P$ is the center of circle $ABC$.
So, similarly, we can show that $P$ is also the center of circle $DEF$.
Thus,  two circles cutting one another, $ABC$ and $DEF$, have the same center
$P$. The very thing is impossible [Prop.~3.5].

Thus, a circle does not cut a(nother) circle at more than two points. (Which is)
the very thing it was required to show.}
\end{Parallel}

%%%%%%
% Prop 3.11
%%%%%%
\pdfbookmark[1]{Proposition 3.11}{pdf3.11}
\begin{Parallel}{}{} 
\ParallelLText{
\begin{center}
{\large \ggn{11}.}
\end{center}\vspace*{-7pt}

{\greekfont Ἐὰν δύο κύκλοι ἐφάπτωνται ἀλλήλων ἐντός, καὶ ληφθῇ α ὐτῶν
τὰ κέντρα, ἡ ἐπὶ τὰ κέντρα α ὐτῶν ἐπιζευγνυμένη εὐθεῖα καὶ
ἐκβαλλομένη ἐπὶ τὴν συναφὴν πεσεῖται τῶν κύκλων.}

{\greekfont Δύο γὰρ κύκλοι οἱ ΑΒΓ, ΑΔΕ ἐφαπτέσθωσαν ἀλλήλων ἐντὸς κατὰ τὸ
Α σημεῖον, καὶ εἰλήφθω τοῦ μὲν ΑΒΓ κύκλου κέντρον τὸ Ζ, τοῦ δὲ
ΑΔΕ τὸ Η; λέγω, ὅτι ἡ ἀπὸ τοῦ Η ἐπὶ τὸ Ζ ἐπιζευγνυμένη
εὐθεῖα ἐκβαλλομένη ἐπὶ τὸ Α πεσεῖται.}

{\greekfont μὴ γάρ, ἀλλ εἰ δυνατόν, πιπτέτω ὡς ἡ ΖΗΘ, καὶ
ἐπεζεύθθωσαν αἱ ΑΖ, ΑΗ.}

{\greekfont Ἐπεὶ οὖν αἱ ΑΗ, ΗΖ τῆς ΖΑ, τουτέστι τῆς ΖΘ, μείζονές εἰσιν, 
κοινὴ ἀφῃρήσθω ἡ ΖΗ; λοιπὴ ἄρα ἡ ΑΗ λοιπῆς τῆς ΗΘ μείζων ἐστίν. ἴση δὲ ἡ ΑΗ τῇ ΗΔ; καὶ ἡ ΗΔ ἄρα τῆς ΗΘ μείζων ἐστὶν
ἡ ἐλάττων τῆς μείζονος; ὅπερ ἐστὶν ἀδύνατον; οὐκ ἄρα ἡ ἀπὸ
τοῦ Ζ ἐπὶ τὸ Η ἐπιζευγνυμένη εὐθεὶα ἐκτὸς πεσεῖται;
κατὰ τὸ Α ἄρα ἐπὶ τῆς συναφῆς πεσεῖται.}\\~\\

\epsfysize=2.2in
\centerline{\epsffile{Book03/fig11g.eps}}

{\greekfont Ἐὰν ἄρα δύο κύκλοι ἐφάπτωνται ἀλλήλων ἐντός, [καὶ ληφθῇ α ὐτῶν
τὰ κέντρα], ἡ ἐπὶ τὰ κέντρα α ὐτῶν ἐπιζευγνυμένη εὐθεῖα [καὶ
ἐκβαλλομένη] ἐπὶ τὴν συναφὴν πεσεῖται τῶν κύκλων; ὅπερ
ἔδει δεῖχαι.}}

\ParallelRText{
\begin{center}
{\large Proposition 11}
\end{center}

If two circles  touch one another internally, and their
centers are found, (then)  the straight-line joining their
centers, being produced, will fall upon the point of union of
the circles.

For let two circles, $ABC$ and $ADE$, touch one another internally
at point $A$, and let the center $F$ of circle $ABC$ be
found [Prop.~3.1], and (the center) $G$ of (circle) $ADE$ [Prop.~3.1].
I say that the straight-line joining $G$ to $F$, being produced, will fall on $A$.

For (if) not (then), if possible, let it fall like $FGH$ (in the figure), and
let $AF$ and $AG$ be joined.

Therefore, since $AG$ and $GF$ is greater than $FA$, that is to
say $FH$ [Prop.~1.20], let $FG$ be taken from both.
Thus, the remainder $AG$ is greater than the remainder $GH$.
And $AG$ (is) equal to $GD$. Thus, $GD$ is also greater than $GH$, the lesser
than the greater. The very thing is impossible. Thus, the straight-line
joining $F$ to $G$ will not fall outside (one circle but inside the other). Thus, it will fall upon the point of
union (of the circles) at point $A$. 

\epsfysize=2.2in
\centerline{\epsffile{Book03/fig11e.eps}}

Thus, if two circles  touch one another internally, [and their
centers are found], (then)  the straight-line joining their
centers, [being produced], will fall upon the point of union of
the circles. (Which is) the very thing it was required to show.}
\end{Parallel}

%%%%%%
% Prop 3.12
%%%%%%
\pdfbookmark[1]{Proposition 3.12}{pdf3.12}
\begin{Parallel}{}{} 
\ParallelLText{
\begin{center}
{\large \ggn{12}.}
\end{center}\vspace*{-7pt}

{\greekfont Ἐὰν δύο κύκλοι ἐφάπτωνται ἀλλήλων ἐκτός, ἡ ἐπὶ τὰ κέντρα
α ὐτῶν ἐπιζευγνυμένη διὰ τῆς ἐπαφῆς ἐλεύσεται.}\\

\epsfysize=2.5in
\centerline{\epsffile{Book03/fig12g.eps}}

{\greekfont Δύο γὰρ κύκλοι οἱ ΑΒΓ, ΑΔΕ ἐφαπτέσθωσαν ἀλλήλων ἐκτὸς
κατὰ τὸ Α σημεῖον, καὶ εἰλήφθω τοῦ μὲν ΑΒΓ κέντρον τὸ Ζ,
τοῦ δὲ ΑΔΕ τὸ Η; λέγω, ὅτι ἡ ἀπὸ τοῦ Ζ ἐπὶ τὸ Η
ἐπιζευγνυμένη εὐθεῖα διὰ τῆς κατὰ τὸ Α ἐπαφῆς ἐλεύσεται.}

{\greekfont μὴ γάρ, ἀλλ εἰ δυνατόν, ἐρθέσθω ὡς ἡ ΖΓΔΗ, καὶ ἐπεζεύθθωσαν
αἱ ΑΖ, ΑΗ.}

{\greekfont Ἐπεὶ οὖν τὸ Ζ σημεῖον κέντρον ἐστὶ τοῦ ΑΒΓ κύκλου,
ἴση ἐστὶν ἡ ΖΑ τῇ ΖΓ. πάλιν, ἐπεὶ τὸ Η σημεῖον κέντρον
ἐστὶ τοῦ ΑΔΕ κύκλου, ἴση ἐστὶν ἡ ΗΑ τῇ ΗΔ. ἐδείθθη δὲ
καὶ ἡ ΖΑ τῇ ΖΓ ἴση; αἱ ἄρα ΖΑ, ΑΗ ταῖς ΖΓ, ΗΔ ἴσαι εἰσίν; ὥστε ὅλη ἡ ΖΗ τῶν ΖΑ, ΑΗ μείζων ἐστίν; ἀλλὰ καὶ ἐλάττων;
ὅπερ ἐστὶν ἀδύνατον. οὐκ ἄρα ἡ ἀπὸ τοῦ Ζ ἐπὶ τὸ Η
ἐπιζευγνυμένη εὐθεῖα διὰ τῆς κατὰ τὸ Α ἐπαφῆς οὐκ ἐλεύσεται;
δι α ὐτῆς ἄρα.}

{\greekfont Ἐὰν ἄρα δύο κύκλοι ἐφάπτωνται ἀλλήλων ἐκτός, ἡ ἐπὶ τὰ κέντρα
α ὐτῶν ἐπιζευγνυμένη [εὐθεῖα] διὰ τῆς ἐπαφῆς ἐλεύσεται; ὅπερ
ἔδει δεῖχαι.}}

\ParallelRText{
\begin{center}
{\large Proposition 12}
\end{center}

If two circles touch one another externally, (then) the (straight-line)
joining their centers will go through the point of union.

\epsfysize=2.5in
\centerline{\epsffile{Book03/fig12e.eps}}

For let two circles, $ABC$ and $ADE$, touch one another externally at point $A$,
and let the center $F$ of $ABC$ be found [Prop.~3.1], and
(the center) $G$ of $ADE$ [Prop.~3.1]. I say that the straight-line joining
$F$ to $G$ will go through the point of union at $A$.

For (if) not (then), if possible, let it go like $FCDG$ (in the figure), and let $AF$ and
$AG$ be joined.

Therefore, since point $F$ is the center of circle $ABC$, $FA$ is equal to $FC$.
Again, since point $G$ is the center of circle $ADE$, $GA$ is equal to $GD$.
And $FA$ was also shown (to be)  equal to $FC$. Thus, the (straight-lines)
$FA$ and $AG$ are equal to the (straight-lines) $FC$ and $GD$. So the whole
of $FG$ is greater than $FA$ and $AG$. But, (it is) also less [Prop.~1.20].
The very thing is impossible. Thus, the straight-line joining $F$ to $G$
cannot not go through the point of union at $A$. Thus, (it will go) through it.

Thus, if two circles touch one another externally, (then) the [straight-line]
joining their centers will go through the point of union. (Which is)
the very thing it was required to show.}
\end{Parallel}

%%%%%%
% Prop 3.13
%%%%%%
\pdfbookmark[1]{Proposition 3.13}{pdf3.13}
\begin{Parallel}{}{} 
\ParallelLText{
\begin{center}
{\large \ggn{13}.}
\end{center}\vspace*{-7pt}

{\greekfont Κύκλος κύκλου οὐκ ἐφάπτεται κατὰ πλείονα σημεῖα
ἢ  καθ ἕν, ἐάν τε ἐντὸς ἐάν τε ἐκτὸς ἐφάπτηται.}

\epsfysize=2.2in
\centerline{\epsffile{Book03/fig13g.eps}}

{\greekfont Εἰ γὰρ δυνατόν, κύκλος ὁ ΑΒΓΔ κύκλου τοῦ ΕΒΖΔ ἐφαπτέσθω
πρότερον ἐντὸς κατὰ πλείονα σημεῖα ἢ ἓν τὰ Δ, Β.}

{\greekfont Καὶ εἰλήφθω τοῦ μὲν ΑΒΓΔ κύκλου κέντρον τὸ Η, τοῦ δὲ ΕΒΖΔ
τὸ Θ.}

{\greekfont Ἡ ἄρα ἀπὸ τοῦ Η ἐπὶ τὸ Θ ἐπιζευγνυμένη ἐπὶ τὰ Β, Δ πεσεῖται.
πιπτέτω ὡς ἡ ΒΗΘΔ. καὶ ἐπεὶ τὸ Η σημεῖον κέντρον ἐστὶ τοῦ
ΑΒΓΔ κύκλου, ἴση ἐστὶν ἡ ΒΗ τῇ ΗΔ; μείζων ἄρα ἡ ΒΗ τῆς
ΘΔ; πολλῷ ἄρα μείζων ἡ ΒΘ τῆς ΘΔ. πάλιν, ἐπεὶ τὸ Θ σημεῖον
κέντρον ἐστὶ τοῦ ΕΒΖΔ κύκλου, ἴση ἐστὶν ἡ ΒΘ τῇ ΘΔ;
ἐδείθθη δὲ α ὐτῆς καὶ πολλῷ μείζων; ὅπερ ἀδύνατον;
οὐκ ἄρα κύκλος κύκλου ἐφάπτεται ἐντὸς κατὰ πλείονα σημεῖα
ἢ ἕν.}

{\greekfont Λέγω δή, ὅτι οὐδὲ ἐκτός.}

{\greekfont Εἰ γὰρ δυνατόν, κύκλος ὁ ΑΓΚ κύκλου τοῦ ΑΒΓΔ ἐφαπτέσθω ἐκτὸς
κατὰ πλείονα σημεῖα ἢ ἓν τὰ Α, Γ, καὶ ἐπεζεύθθω ἡ ΑΓ.}

{\greekfont Ἔπεὶ οὖν κύκλων τῶν ΑΒΓΔ, ΑΓΚ εἴληπται ἐπὶ
τῆς περιφερείας ἑκατέρου δύο τυθόντα σημεῖα τὰ Α, Γ, ἡ
ἐπὶ τὰ σημεῖα ἐπιζευγνυμένη εὐθεῖα ἐντὸς ἑκατέρου πεσεῖται;
ἀλλὰ τοῦ μὲν ΑΒΓΔ ἐντὸς ἔπεσεν, τοῦ δὲ ΑΓΚ ἐκτός; ὅπερ
ἄτοπον; οὐκ ἄρα κύκλος κύκλου ἐφάπτεται ἐκτὸς κατὰ πλείονα
σημεῖα ἢ ἕν. ἐδείθθη δέ, ὅτι οὐδὲ ἐντός.}

{\greekfont Κύκλος ἄρα κύκλου οὐκ ἐφάπτεται κατὰ πλείονα σημεῖα ἢ
[καθ] ἕν, ἐάν τε ἐντὸς ἐάν τε ἐκτὸς ἐφάπτηται; ὅπερ ἔδει
δεῖχαι.}}

\ParallelRText{
\begin{center}
{\large Proposition 13}
\end{center}

A circle does not touch a(nother) circle at more than one point, whether they
touch internally or externally.

\epsfysize=2.2in
\centerline{\epsffile{Book03/fig13e.eps}}

For, if possible, let circle $ABDC$$^{\,\dag}$ touch circle $EBFD$---first of all, internally---at more than one point, $D$ and $B$.

And let the center $G$ of circle $ABDC$ be found [Prop.~3.1], and
(the center) $H$ of $EBFD$ [Prop.~3.1].

Thus, the (straight-line) joining $G$ and $H$ will fall on $B$ and $D$ [Prop.~3.11].
Let it fall like $BGHD$ (in the figure). And since point $G$ is the center of
circle $ABDC$, $BG$ is  equal to $GD$. Thus, $BG$ (is) greater than $HD$.
Thus, $BH$ (is) much greater than $HD$. Again, since point $H$ is the
center of circle $EBFD$, $BH$ is equal to $HD$. But it was also shown
(to be) much greater than it. The very thing (is) impossible.
Thus, a circle does not touch a(nother) circle internally at more than one point.

So, I say that neither (does it touch) externally (at more than one point).

For, if possible, let circle $ACK$ touch circle $ABDC$ externally at more
than one point, $A$ and $C$. And let $AC$ be joined.

Therefore, since  two points, $A$ and $C$,
be taken at random on the circumference
of each of the circles $ABDC$ and $ACK$, the straight-line joining
the points will fall inside each (circle) [Prop.~3.2]. But, it fell
inside $ABDC$, and outside $ACK$ [Def.~3.3]. The very thing (is) absurd.
Thus, a circle does not touch a(nother) circle externally  at more than one
point. And it was shown that neither (does it) internally.

Thus, a circle does not touch a(nother) circle at more than one point, whether they touch internally or externally. (Which is) the very thing it was required to
show.}
\end{Parallel}
{\footnotesize \noindent$^\dag$ The Greek text has ``$ABCD$'', which is obviously a mistake.}

%%%%%%
% Prop 3.14
%%%%%%
\pdfbookmark[1]{Proposition 3.14}{pdf3.14}
\begin{Parallel}{}{} 
\ParallelLText{
\begin{center}
{\large \ggn{14}.}
\end{center}\vspace*{-7pt}

{\greekfont Ἐν κύκλῳ αἱ ἴσαι εὐθεῖαι ἴσον ἀπέθουσιν ἀπὸ τοῦ κέντρου, καὶ
αἱ ἴσον ἀπέθουσαι ἀπὸ τοῦ κέντρου ἴσαι ἀλλήλαις εἰσίν.}\\

\epsfysize=2.2in
\centerline{\epsffile{Book03/fig14g.eps}}

{\greekfont Ἔστω κύκλος ὁ ΑΒΓΔ, καὶ ἐν α ὐτῷ ἴσαι εὐθεῖαι ἔστωσαν αἱ
ΑΒ, ΓΔ; λέγω, ὅτι αἱ ΑΒ, ΓΔ ἴσον ἀπέθουσιν ἀπὸ τοῦ
κέντρου.}

{\greekfont Εἰλήφθω γὰρ τὸ κέντον τοῦ ΑΒΓΔ κύκλου καὶ ἔστω τὸ Ε, καὶ
ἀπὸ τοῦ Ε ἐπὶ τὰς ΑΒ, ΓΔ κάθετοι ἤθθωσαν αἱ ΕΖ, ΕΗ, καὶ
ἐπεζεύθθωσαν αἱ ΑΕ, ΕΓ.}

{\greekfont Ἐπεὶ οὖν εὐθεῖά τις δὶα τοῦ κέντρου ἡ ΕΖ εὐθεῖάν τινα μὴ
διὰ τοῦ κέντρου τὴν ΑΒ πρὸς ὀρθὰς τέμνει, καὶ δίθα α ὐτὴν
τέμνει. ἴση ἄρα ἡ ΑΖ τῇ ΖΒ; διπλῆ ἄρα ἡ ΑΒ τῆς ΑΖ. διὰ τὰ
α ὐτὰ δὴ καὶ ἡ ΓΔ τῆς ΓΗ ἐστι διπλῆ; καί ἐστιν ἴση ἡ ΑΒ
τῇ ΓΔ; ἴση ἄρα καὶ ἡ ΑΖ τῇ ΓΗ. καὶ ἐπεὶ ἴση ἐστὶν ἡ ΑΕ
τῇ ΕΓ, ἴσον καὶ τὸ ἀπὸ τῆς ΑΕ τῷ ἀπὸ τῆς ΕΓ.
ἀλλὰ τῷ μὲν ἀπὸ τῆς ΑΕ ἴσα τὰ ἀπὸ τῶν ΑΖ, ΕΖ; ὀρθὴ γὰρ ἡ πρὸς τῷ Ζ γωνία; τῷ δὲ ἀπὸ τῆς ΕΓ ἴσα τὰ ἀπὸ τῶν ΕΗ, ΗΓ;
ὀρθὴ γὰρ ἡ πρὸς τῷ Η γωνία; τὰ ἄρα ἀπὸ τῶν ΑΖ, ΖΕ ἴσα
ἐστὶ τοῖς ἀπὸ τῶν ΓΗ, ΗΕ, ὧν τὸ ἀπὸ τῆς ΑΖ ἴσον ἐστὶ
τῷ ἀπὸ τῆς ΓΗ; ἴση γάρ ἐστιν ἡ ΑΖ τῇ ΓΗ; λοιπὸν ἄρα
τὸ ἀπὸ τῆς ΖΕ τῷ ἀπὸ τῆς ΕΗ ἴσον ἐστίν; ἴση ἄρα ἡ
ΕΖ τῇ ΕΗ. ἐν δὲ κύκλῳ ἴσον ἀπέθειν ἀπὸ τοῦ κέντρου
εὐθεῖαι λέγονται, ὅταν αἱ ἀπὸ τοῦ κέντρου ἐπ α ὐτὰς
κάθετοι ἀγόμεναι ἴσαι ὦσιν; αἱ ἄρα ΑΒ, ΓΔ ἴσον ἀπέθουσιν
ἀπὸ τοῦ κέντρου.}

{\greekfont Ἀλλὰ δὴ αἱ ΑΒ, ΓΔ εὐθεῖαι ἴσον ἀπεθέτωσαν ἀπὸ τοῦ κέντρου,
τουτέστιν ἴση ἔστω ἡ ΕΖ τῇ ΕΗ. λέγω, ὅτι ἴση ἐστὶ καὶ
ἡ ΑΒ τῇ ΓΔ.}

{\greekfont Τῶν γὰρ α ὐτῶν κατασκευασθέντων ὁμοίως δείχομεν, ὅτι διπλῆ
ἐστιν ἡ μὲν ΑΒ τῆς ΑΖ, ἡ δὲ ΓΔ τῆς ΓΗ; καὶ ἐπεὶ ἴση ἐστὶν ἡ ΑΕ τῇ ΓΕ, ἴσον ἐστὶ τὸ ἀπὸ τῆς 
ΑΕ τῷ ἀπὸ τῆς ΓΕ; ἀλλὰ τῷ
μὲν ἀπὸ τῆς ΑΕ ἴσα ἐστὶ τὰ ἀπὸ τῶν ΕΖ, ΖΑ, τῷ δὲ ἀπὸ
τῆς ΓΕ 
ἴσα τὰ ἀπὸ τῶν ΕΗ, ΗΓ. τὰ ἄρα ἀπὸ τῶν
ΕΖ, ΖΑ ἴσα ἐστὶ τοῖς ἀπὸ τῶν ΕΗ, ΗΓ; ὧν τὸ ἀπὸ τῆς
ΕΖ τῷ ἀπὸ τῆς ΕΗ ἐστιν ἴσον; ἴση γὰρ ἡ ΕΖ τῇ ΕΗ; λοιπὸν ἄρα
τὸ ἀπὸ τῆς ΑΖ ἴσον ἐστὶ τῷ ἀπὸ τῆς ΓΗ; ἴση ἄρα ἡ ΑΖ
τῇ ΓΗ; καί ἐστι τῆς μὲν ΑΖ διπλῆ ἡ ΑΒ, τῆς δὲ ΓΗ διπλῆ ἡ ΓΔ;
ἴση ἄρα ἡ ΑΒ τῇ ΓΔ.}

{\greekfont Ἐν κύκλῳ ἄρα αἱ ἴσαι εὐθεῖαι ἴσον ἀπέθουσιν ἀπὸ τοῦ κέντρου, καὶ
αἱ ἴσον ἀπέθουσαι ἀπὸ τοῦ κέντρου ἴσαι ἀλλήλαις εἰσίν;
ὅπερ ἔδει δεῖχαι.}}

\ParallelRText{
\begin{center}
{\large Proposition 14}
\end{center}

In  a circle, equal straight-lines are equally far from the center, and
(straight-lines) which are equally far from the center are equal to one another.

\epsfysize=2.2in
\centerline{\epsffile{Book03/fig14e.eps}}

Let $ABDC$$^{\,\dag}$ be a circle, and let $AB$ and $CD$ be equal straight-lines within it.
I say that $AB$ and $CD$ are equally far from the center.

For let the center of circle $ABDC$ be found [Prop.~3.1], and
let it be (at) $E$. And let $EF$ and $EG$ be drawn from (point) $E$, perpendicular
to $AB$ and $CD$ (respectively) [Prop.~1.12]. And let $AE$ and $EC$ be joined.

Therefore, since some straight-line, $EF$, through the center (of the circle),
cuts some (other) straight-line, $AB$, not through the center, at right-angles,
it also cuts it in half [Prop.~3.3]. Thus, $AF$ (is) equal to $FB$. Thus, 
$AB$ (is) double $AF$. So, for the same (reasons), $CD$ is also double $CG$.
And $AB$ is equal to $CD$. Thus, $AF$ (is) also equal to $CG$. And since
$AE$ is equal to $EC$, the (square) on $AE$ (is) also equal to the (square)
on $EC$. But, the (sum of the squares) on $AF$
and $EF$ (is) equal to the (square) on $AE$. For the angle at $F$ (is) a right-angle [Prop.~1.47]. And the
(sum of the squares) on $EG$ and $GC$ (is) equal to the (square) on $EC$. For the angle at $G$ (is) a right-angle [Prop.~1.47]. Thus, the (sum of the squares)
on $AF$ and $FE$ is equal to the (sum of the squares) on $CG$ and $GE$, 
of which
the (square) on $AF$ is equal to the (square) on $CG$. For $AF$ is equal to
$CG$. Thus, the remaining (square) on $FE$ is equal to the (remaining square) on  $EG$. Thus, $EF$ (is) equal to $EG$. And straight-lines in a circle are said to be equally far from the center when perpendicular (straight-lines)
which are drawn to them from the center  are equal [Def.~3.4]. Thus, $AB$ and
$CD$ are equally far from the center.

So, let the straight-lines $AB$ and $CD$ be equally far from the center. That is to say, let $EF$ be equal to $EG$. I say that $AB$ is also equal to $CD$.

For, with the same construction,  we can, similarly, show that
$AB$ is double $AF$, and $CD$ (double) $CG$. And since $AE$ is equal to $CE$, the (square) on $AE$ is equal to the (square) on $CE$. But, the (sum of the squares)
on $EF$ and $FA$ is equal to the (square) on $AE$ [Prop.~1.47]. And the (sum of
the squares) on $EG$ and $GC$ (is) equal to the (square) on $CE$ [Prop.~1.47]. Thus, the (sum of the squares) on $EF$ and $FA$ is equal to the (sum of the squares) on $EG$ and $GC$, of which the (square) on $EF$ is equal to the (square)
on $EG$. For $EF$ (is) equal to $EG$. Thus, the remaining (square) on  $AF$ is equal to the (remaining square) on  $CG$. Thus, $AF$ (is) equal to $CG$. And $AB$ is double $AF$, and
$CD$ double $CG$. Thus, $AB$ (is) equal to $CD$.

Thus, in a circle, equal straight-lines are equally far from the center, and
(straight-lines) which are equally far from the center are equal to one another.
(Which is) the very thing it was required to show.}
\end{Parallel}
{\footnotesize \noindent$^\dag$ The Greek text has ``$ABCD$'', which is obviously a
mistake.}

%%%%%%
% Prop 3.15
%%%%%%
\pdfbookmark[1]{Proposition 3.15}{pdf3.15}
\begin{Parallel}{}{} 
\ParallelLText{
\begin{center}
{\large \ggn{15}.}
\end{center}\vspace*{-7pt}

{\greekfont Ἐν κύκλῳ μεγίστη μὲν ἡ διάμετρος, τῶν δὲ ἄλλων ἀεὶ ἡ ἔγγιον τοῦ κέντρου τῆς ἀπώτερον μείζων ἐστίν.}

{\greekfont Ἔστω κύκλος ὁ ΑΒΓΔ, διάμετρος δὲ α ὐτοῦ ἔστω ἡ ΑΔ, κέντρον δὲ τὸ Ε, καὶ ἔγγιον μὲν τῆς ΑΔ διαμέτρου ἔστω ἡ ΒΓ, ἀπώτερον δὲ ἡ ΖΗ; λέγω, ὅτι μεγίστη μέν ἐστιν ἡ ΑΔ, μείζων δὲ ἡ ΒΓ τῆς ΖΗ.}

{\greekfont Ἤθθωσαν γὰρ ἀπὸ τοῦ Ε κέντρου ἐπὶ τὰς ΒΓ, ΖΗ κάθετοι αἱ ΕΘ, ΕΚ.
καὶ ἐπεὶ ἔγγιον μὲν τοῦ κέντρου ἐστὶν ἡ ΒΓ, ἀπώτερον δὲ
ἡ ΖΗ, μείζων ἄρα ἡ ΕΚ τῆς ΕΘ. κείσθω τῇ ΕΘ ἴση ἡ ΕΛ,
καὶ διὰ τοῦ Λ τῇ ΕΚ πρὸς ὀρθὰς ἀθθεῖσα ἡ ΛΜ διήθθω ἐπὶ
τὸ Ν, καὶ ἐπεζεύθθωσαν αἱ ΜΕ, ΕΝ, ΖΕ, ΕΗ.}

{\greekfont Καὶ ἐπεὶ ἴση ἐστὶν ἡ ΕΘ τῇ ΕΛ, ἴση ἐστὶ καὶ ἡ ΒΓ τῇ ΜΝ. πάλιν, ἐπεὶ ἴση ἐστὶν ἡ μὲν ΑΕ τῇ ΕΜ, ἡ δὲ ΕΔ τῇ ΕΝ, ἡ ἄρα
ΑΔ ταῖς ΜΕ, ΕΝ ἴση ἐστίν. ἀλλ αἱ μὲν ΜΕ, ΕΝ τῆς ΜΝ μείζονές
εἰσιν [καὶ ἡ ΑΔ τῆς ΜΝ μείζων ἐστίν], ἴση δὲ ἡ ΜΝ τῇ ΒΓ;
ἡ ΑΔ ἄρα τῆς ΒΓ μείζων ἐστίν. καὶ ἐπεὶ δύο αἱ ΜΕ, ΕΝ δύο
ταῖς ΖΕ, ΕΗ ἴσαι εἰσίν, καὶ γωνία ἡ ὑπὸ ΜΕΝ γωνίας τῆς ὑπὸ
ΖΕΗ μείζων [ἐστίν], βάσις ἄρα ἡ ΜΝ βάσεως τῆς ΖΗ μείζων
ἐστίν. ἀλλὰ ἡ ΜΝ τῇ ΒΓ ἐδείθθη ἴση [καὶ ἡ ΒΓ τῆς ΖΗ
μείζων ἐστίν]. μεγίστη μὲν ἄρα ἡ ΑΔ διάμετρος, μείζων
δὲ ἡ ΒΓ τῆς ΖΗ.}~\\~\\~\\~\\

\epsfysize=2.2in
\centerline{\epsffile{Book03/fig15g.eps}}

{\greekfont Ἐν κύκλῳ ἄρα μεγίστη μὲνέστιν ἡ διάμετρος, τῶν δὲ ἄλλων ἀεὶ ἡ ἔγγιον τοῦ κέντρου τῆς ἀπώτερον μείζων ἐστίν; ὅπερ
ἔδει δεῖχαι.}}

\ParallelRText{
\begin{center}
{\large Proposition 15}
\end{center}

In a circle, a diameter (is) the greatest (straight-line), and for 
the others, a (straight-line) nearer to the center is always greater than one further away.

Let $ABCD$ be a circle, and let $AD$ be its diameter, and $E$ (its) center.
And let $BC$ be nearer to the diameter $AD$,$^\dag$ and $FG$ further away. I say that
$AD$ is the greatest (straight-line), and $BC$ (is) greater than $FG$.

For let $EH$ and $EK$ be drawn from the center $E$, at right-angles
to $BC$ and $FG$ (respectively) [Prop.~1.12]. And since $BC$ is nearer to the center, and $FG$ further away, $EK$ (is) thus greater than $EH$ [Def.~3.5]. Let $EL$ be made equal to
$EH$ [Prop.~1.3]. And $LM$ being drawn through $L$, at right-angles to $EK$ [Prop.~1.11], 
let it be drawn through to $N$. And let $ME$, $EN$, $FE$, and $EG$
be joined.

And since $EH$ is equal to $EL$, $BC$ is also equal to $MN$ [Prop.~3.14]. Again,
since $AE$ is equal to $EM$, and $ED$ to $EN$, $AD$ is thus equal to $ME$ and $EN$.
But, $ME$ and $EN$ is greater than $MN$ [Prop.~1.20] [also $AD$ is greater than $MN$], and $MN$ (is) equal to $BC$. Thus, $AD$ is greater than $BC$. And since
the two (straight-lines) $ME$, $EN$ are equal to the two (straight-lines) $FE$, $EG$ (respectively),
and angle $MEN$ [is] greater than angle $FEG$,$^\ddag$ the base $MN$ is thus
greater than the base $FG$ [Prop.~1.24]. But, $MN$ was shown (to be) equal to $BC$ [(so) $BC$ is also greater than $FG$]. Thus, the diameter $AD$ (is) the
greatest (straight-line), and $BC$ (is) greater than $FG$.

\epsfysize=2.2in
\centerline{\epsffile{Book03/fig15e.eps}}

Thus, in a circle, a diameter (is) the greatest (straight-line), and for 
the others, a (straight-line) nearer to the center is always greater than one further away. (Which is) the very thing it was required to show.}
\end{Parallel}
{\footnotesize \noindent$^\dag$ Euclid should have said ``to the center'', rather than "to the diameter $AD$", since $BC$, $AD$ and $FG$ are not necessarily parallel.\\
$^\ddag$  This is not proved,
except by reference to the figure.}

%%%%%%
% Prop 3.16
%%%%%%
\pdfbookmark[1]{Proposition 3.16}{pdf3.16}
\begin{Parallel}{}{} 
\ParallelLText{
\begin{center}
{\large \ggn{16}.}
\end{center}\vspace*{-7pt}

{\greekfont Ἡ τῇ διαμέτρῳ τοῦ κύκλου πρὸς ὀρθὰς ἀπ ἄκρας ἀγομένη
ἐκτὸς πεσεῖται τοῦ κύκλου, καὶ εἰς τὸν μεταχὺ τόπον τῆς τε
εὐθείας καὶ τῆς περιφερείας ἑτέρα εὐθεῖα οὐ παρεμπεσεῖται,
καὶ ἡ μὲν τοῦ ἡμικυκλίου γωνία ἁπάσης γωνίας ὀχείας
εὐθυγράμμου μείζων ἐστίν, ἡ δὲ λοιπὴ ἐλάττων.}

{\greekfont Ἔστω κύκλος ὁ ΑΒΓ περὶ κέντρον τὸ Δ καὶ διάμετρον τὴν
ΑΒ; λέγω, ὅτι ἡ ἀπὸ τοῦ Α τῇ ΑΒ πρὸς ὀρθὰς ἀπ ἄκρας ἀγομένη ἐκτὸς πεσεῖται τοῦ κύκλου.}

{\greekfont μὴ γάρ, ἀλλ εἰ δυνατόν, πιπτέτω ἐντὸς ὡς ἡ ΓΑ, καὶ ἐπεζεύθθω ἡ ΔΓ.}

{\greekfont Ἐπεὶ ἴση ἐστὶν ἡ ΔΑ τῇ ΔΓ, ἴση ἐστὶ καὶ γωνία ἡ ὑπὸ ΔΑΓ γωνίᾳ τῇ ὑπὸ ΑΓΔ. ὀρθὴ δὲ ἡ ὑπὸ ΔΑΓ; ὀρθὴ ἄρα καὶ ἡ
ὑπὸ ΑΓΔ; τριγώνου δὴ τοῦ ΑΓΔ αἱ δύο γωνίαι αἱ ὑπὸ ΔΑΓ, ΑΓΔ
δύο ὀρθαῖς ἴσαι εἰσίν; ὅπερ ἐστὶν ἀδύνατον. οὐκ ἄρα ἡ ἀπὸ
τοῦ Α σημείου τῇ ΒΑ πρὸς ὀρθὰς ἀγομένη ἐντὸς πεσεῖται
τοῦ κύκλου. ὁμοίως δὴ δεῖχομεν, ὅτι οὐδ ἐπὶ τῆς περιφερείας;
ἐκτὸς ἄρα.}

{\greekfont Πιπτέτω ὡς ἡ ΑΕ; λέγω δή, ὅτι εἰς τὸν μεταχὺ τόπον τῆς τε ΑΕ
εὐθείας καὶ τῆς ΓΘΑ περιφερείας ἑτέρα εὐθεῖα οὐ παρεμπεσεῖται.}

{\greekfont Εἰ γὰρ δυνατόν, παρεμπιπτέτω ὡς ἡ ΖΑ, καὶ ἤθθω
ἀπὸ τοῦ Δ σημείου ἐπὶ τὴν ΖΑ κάθετος ἡ ΔΗ. καὶ ἐπεὶ ὀρθή
ἐστιν ἡ ὑπὸ ΑΗΔ, ἐλάττων δὲ ὀρθῆς ἡ ὑπὸ ΔΑΗ, μείζων ἄρα
ἡ ΑΔ τῆς ΔΗ. ἴση δὲ ἡ ΔΑ τῇ ΔΘ; μείζων ἄρα ἡ ΔΘ τῆς
ΔΗ, ἡ ἐλάττων τῆς μείζονος; ὅπερ ἐστὶν ἀδύνατον. οὐκ ἄρα
εἰς τὸν μεταχὺ τόπον τῆς τε εὐθείας καὶ τῆς περιφερείας ἑτέρα
εὐθεῖα παρεμπεσεῖται.}\\~\\~\\~\\

\epsfysize=2in
\centerline{\epsffile{Book03/fig16g.eps}}

{\greekfont Λέγω, ὅτι καὶ ἡ μὲν τοῦ ἡμικυκλίου γωνία ἡ περιεθομένη ὑπό
τε τῆς ΒΑ εὐθείας καὶ τῆς ΓΘΑ περιφερείας ἁπάσης γωνίας
ὀχείας εὐθυγράμμου μείζων ἐστίν, ἡ δὲ λοιπὴ ἡ περιεθομένη
ὑπό τε τῆς ΓΘΑ περιφερείας καὶ τῆς ΑΕ εὐθείας ἁπάσης γωνίας
ὀχείας εὐθυγράμμου ἐλάττων ἐστίν.}

{\greekfont Εἰ γὰρ ἐστί τις γωνία εὐθύγραμμος μείζων μὲν τῆς περιεθομένης
ὑπό τε τῆς ΒΑ εὐθείας καὶ τῆς ΓΘΑ περιφερείας, ἐλάττων δὲ τῆς
περιεθομένης ὑπό τε τῆς ΓΘΑ περιφερείας καὶ τὴς ΑΕ εὐθείας, εἰς
τὸν μεταχὺ τόπον τῆς τε ΓΘΑ περιφερείας καὶ τῆς ΑΕ εὐθείας
εὐθεῖα παρεμπεσεῖται, ἥτις ποιήσει μείζονα
μὲν τῆς περιεθομένης ὑπὸ τε τῆς ΒΑ εὐθείας καὶ τῆς  ΓΘΑ
περιφερείας ὑπὸ εὐθειῶν περιεθομένην, ἐλάττονα δὲ τῆς περιεθομένης ὑπό τε τῆς ΓΘΑ περιφερείας καὶ τῆς ΑΕ εὐθείας.
οὐ παρεμπίπτει δέ; οὐκ ἄρα τῆς περιεθομένης γωνίας ὑπό
τε τῆς ΒΑ εὐθείας καὶ τῆς ΓΘΑ περιφερείας ἔσται μείζων ὀχεῖα
ὑπὸ εὐθειῶν περιεθομένη, οὐδὲ μὴν ἐλάττων τῆς περιεθομένης
ὑπό τε τῆς ΓΘΑ περιφερείας καὶ τῆς ΑΕ εὐθείας.}\\~\\

\begin{center}
{\large {\greekfont Πόρισμα}.}
\end{center}\vspace*{-7pt}

{\greekfont Ἐκ δὴ τούτου φανερόν, ὅτι ἡ τῇ διαμέτρῳ τοῦ κύκλου πρὸς
ὀρθὰς ἀπ ἄκρας ἀγομένη ἐφάπτεται τοῦ κύκλου [καὶ ὅτι εὐθεῖα κύκλου καθ ἓν μόνον ἐφάπτεται σημεῖον, ἐπειδήπερ καὶ ἡ κατὰ δύο
α ὐτῷ συμβάλλουσα ἐντὸς α ὐτοῦ πίπτουσα ἐδείθθη]; ὅπερ
ἔδει δεῖχαι.}}

\ParallelRText{
\begin{center}
{\large Proposition 16}
\end{center}

A (straight-line) drawn at right-angles to the diameter of a circle, from its
end, will fall outside the circle. And another straight-line cannot be inserted
into the space between the (aforementioned) straight-line and the circumference. And the
  angle of the semi-circle is greater than any acute rectilinear angle whatsoever, 
and the remaining (angle is) less (than any acute rectilinear angle).

Let $ABC$ be a circle around the center $D$ and the diameter $AB$. I say that
the (straight-line) drawn from $A$, at right-angles to $AB$ [Prop~1.11], from its end, will fall
outside the circle.

For (if) not (then), if possible, let it fall inside, like $CA$ (in the figure), and let
$DC$ be joined.

Since $DA$ is equal to $DC$, angle $DAC$ is also equal to angle $ACD$ [Prop.~1.5].
And $DAC$ (is) a right-angle. Thus, $ACD$ (is) also a right-angle. So, in triangle
$ACD$, the two angles $DAC$ and $ACD$ are equal to two right-angles. 
The very thing is
impossible [Prop.~1.17]. Thus, the (straight-line) drawn from point $A$,
at right-angles to $BA$, will not fall inside the circle. So, similarly, we can show 
that neither (will it fall) on the circumference. Thus, (it will fall) outside
(the circle).

Let it fall like $AE$ (in the figure). So, I say that another straight-line cannot
be inserted into the space between the straight-line $AE$ and the
circumference $CHA$.

For, if possible, let it be inserted like $FA$ (in the figure), and let $DG$ 
be drawn from point $D$, perpendicular to $FA$ [Prop.~1.12]. And since $AGD$ is a right-angle, and
 $DAG$ (is) less than a right-angle, $AD$ (is) thus greater than $DG$  [Prop.~1.19]. And $DA$ (is) equal to $DH$. Thus, $DH$ (is)
greater than $DG$, the lesser than the greater. The very thing is impossible.
Thus, another straight-line cannot be inserted into the space between the
straight-line ($AE$) and the circumference.

\epsfysize=2in
\centerline{\epsffile{Book03/fig16e.eps}}

And I also say that the  semi-circular angle  contained by the straight-line
$BA$ and the circumference $CHA$ is greater than any acute rectilinear angle whatsoever,
and the remaining (angle) contained by the circumference $CHA$ and the straight-line
$AE$ is less than any acute rectilinear angle whatsoever.

For if any rectilinear angle is greater than the (angle) contained by the
straight-line $BA$ and the circumference $CHA$, or  less than the (angle)
contained by the circumference $CHA$ and the straight-line $AE$, (then)
a straight-line can be inserted into the space between the circumference 
$CHA$ and the straight-line $AE$---anything which will make (an angle)
contained by straight-lines greater than the angle contained by the straight-line $BA$ and the circumference $CHA$, or less than the (angle) contained 
by the circumference $CHA$ and the straight-line $AE$. But (such a straight-line)
cannot be inserted. Thus, an acute (angle) contained by straight-lines
cannot be greater than the angle contained by the straight-line $BA$ and
the circumference $CHA$, neither (can it be) less than the (angle) contained by
the circumference $CHA$ and the straight-line $AE$.\\

\begin{center}
{\large Corollary}
\end{center}\vspace*{-7pt}

So, from this, (it is) manifest that a (straight-line) drawn at right-angles
to the diameter of a circle, from its extremity, touches the circle [and that the
straight-line  touches the circle at a single point, inasmuch as it was also
shown that a (straight-line) meeting (the circle) at two (points) falls inside it [Prop.~3.2]\,]. (Which is) the very thing it was required to show. }
\end{Parallel}

%%%%%%
% Prop 3.17
%%%%%%
\pdfbookmark[1]{Proposition 3.17}{pdf3.17}
\begin{Parallel}{}{} 
\ParallelLText{
\begin{center}
{\large \ggn{17}.}
\end{center}\vspace*{-7pt}

{\greekfont Ἀπὸ τοῦ δοθέντος σημείου τοῦ δοθέντος κύκλου ἐφαπτομένην
εὐθεῖαν γραμμὴν ἀγαγεῖν.}

{\greekfont Ἔστω τὸ μὲν δοθὲν σημεῖον τὸ Α, ὁ δὲ δοθεὶς κύκλος ὁ ΒΓΔ; δεῖ δὴ ἀπὸ τοῦ Α σημείου τοῦ ΒΓΔ κύκλου ἐφαπτομένην
εὐθεῖαν γραμμὴν ἀγαγεῖν.}

\epsfysize=2.2in
\centerline{\epsffile{Book03/fig17g.eps}}

{\greekfont Εἰλήφθω γὰρ τὸ κέντρον τοῦ κύκλου τὸ Ε, καὶ ἐπεζεύθθω ἡ ΑΕ,
καὶ κέντρῳ μὲν τῷ Ε διαστήματι δὲ τῷ ΕΑ κύκλος γεγράφθω ὁ ΑΖΗ,
καὶ ἀπὸ τοῦ  Δ τῇ ΕΑ πρὸς ὀρθὰς ἤθθω ἡ ΔΖ, καὶ
ἐπεζεύθθωσαν αἱ ΕΖ, ΑΒ; λέγω, ὅτι ἀπὸ τοῦ Α σημείου τοῦ
ΒΓΔ κύκλου ἐφαπτομένη ἦκται ἡ ΑΒ.}

{\greekfont Ἐπεὶ γὰρ τὸ Ε κέντρον ἐστὶ τῶν ΒΓΔ, ΑΖΗ κύκλων, ἴση
ἄρα ἐστὶν ἡ μὲν ΕΑ τῇ ΕΖ, ἡ δὲ ΕΔ τῇ ΕΒ; δύο δὴ αἱ ΑΕ, ΕΒ δύο ταῖς ΖΕ, ΕΔ ἴσαι εἰσίν; καὶ γωνίαν κοινὴν περιέθουσι τὴν
πρὸς τῷ Ε; βάσις ἄρα ἡ ΔΖ βάσει τῇ ΑΒ ἴση ἐστίν, καὶ τὸ ΔΕΖ τρίγωνον τῷ ΕΒΑ τριγώνῳ ἴσον ἐστίν, καὶ αἱ λοιπαὶ
γωνίαι ταῖς λοιπαῖς γωνίαις; ἴση ἄρα ἡ ὑπὸ ΕΔΖ τῇ ὑπὸ ΕΒΑ.
ὀρθὴ δὲ ἡ ὑπὸ ΕΔΖ; ὀρθὴ ἄρα καὶ ἡ ὑπὸ ΕΒΑ. καί ἐστιν
ἡ ΕΒ ἐκ τοῦ κέντρου; ἡ δὲ τῇ διαμέτρῳ τοῦ κύκλου πρὸς
ὀρθὰς ἀπ ἄκρας ἀγομένη ἐφάπτεται τοῦ κύκλου; ἡ ΑΒ ἄρα
ἐφάπτεται τοῦ ΒΓΔ κύκλου.}

{\greekfont Ἀπὸ τοῦ  ἄρα δοθέντος σημείου τοῦ Α τοῦ δοθέντος κύκλου τοῦ ΒΓΔ ἐφαπτομένη
εὐθεῖα γραμμὴ ἦκται ἡ ΑΒ; ὅπερ ἔδει ποιῆσαι.}}

\ParallelRText{
\begin{center}
{\large Proposition 17}
\end{center}

To draw a straight-line touching a given circle from a given point.

Let $A$ be the given point, and $BCD$ the given circle. So it is required to draw a
straight-line touching  circle $BCD$ from point $A$.

\epsfysize=2.2in
\centerline{\epsffile{Book03/fig17e.eps}}

For let the center $E$ of the circle be found [Prop.~3.1], and let $AE$ be joined. And let (the circle) $AFG$ be drawn with center
$E$ and radius $EA$. And let $DF$ be drawn from from (point) $D$,
at right-angles to $EA$ [Prop.~1.11]. And let $EF$ and $AB$ 
be  joined. I say that the (straight-line) $AB$ has been drawn from  point $A$
touching circle $BCD$.

For since $E$ is the center of circles $BCD$ and $AFG$, $EA$ is thus equal to $EF$,
and $ED$ to $EB$. So the two (straight-lines) $AE$, $EB$ are equal to
the two (straight-lines) $FE$, $ED$ (respectively). And they contain a common angle at $E$.
Thus, the base $DF$ is equal to the base $AB$, and triangle $DEF$ is equal to
triangle $EBA$, and the remaining angles (are equal) to the (corresponding)
remaining angles [Prop.~1.4]. Thus, (angle) $EDF$ (is) equal to $EBA$. And
$EDF$ (is) a right-angle. Thus, $EBA$ (is) also a right-angle. And $EB$ is a radius.
And a (straight-line) drawn at right-angles to the diameter of a circle, from
its extremity,  touches the circle [Prop.~3.16~corr.]. Thus, $AB$ touches circle $BCD$.

Thus, the straight-line $AB$ has been drawn touching the given circle $BCD$ from the given point $A$. (Which is) the very thing it was required to do.}
\end{Parallel}

%%%%%%
% Prop 3.18
%%%%%%
\pdfbookmark[1]{Proposition 3.18}{pdf3.18}
\begin{Parallel}{}{} 
\ParallelLText{
\begin{center}
{\large \ggn{18}.}
\end{center}\vspace*{-7pt}

{\greekfont Ἐὰν κύκλου ἐφάπτηταί τις εὐθεῖα, ἀπὸ δὲ τοῦ κέντρου ἐπὶ τὴν ἁφὴν ἐπιζευθθῇ τις εὐθεῖα, ἡ ἐπιζευθθεῖσα κάθετος ἔσται
ἐπὶ τὴν ἐφαπτομένην.}

{\greekfont Κύκλου γὰρ τοῦ ΑΒΓ ἐφαπτέσθω τις εὐθεῖα ἡ ΔΕ κατὰ τὸ Γ
σημεῖον, καὶ εἰλήφθω τὸ κέντρον τοῦ ΑΒΓ κύκλου τὸ
Ζ, καὶ ἀπὸ τοῦ Ζ ἐπὶ τὸ Γ ἐπεζεύθθω ἡ ΖΓ; λέγω, ὅτι ἡ ΖΓ
κάθετός ἐστιν ἐπὶ τὴν ΔΕ.}

{\greekfont Εἰ γὰρ μή, ἤθθω ἀπὸ τοῦ Ζ ἐπὶ τὴν ΔΕ κάθετος ἡ ΖΗ.}

{\greekfont Ἐπεὶ οὖν ἡ ὑπὸ ΖΗΓ γωνία ὀρθή ἐστιν, ὀχεῖα ἄρα
ἐστὶν ἡ ὑπὸ ΖΓΗ; ὑπὸ δὲ τὴν μείζονα γωνίαν ἡ μείζων
πλευρὰ ὑποτείνει; μείζων ἄρα ἡ ΖΓ τῆς ΖΗ; ἴση δὲ ἡ ΖΓ τῇ ΖΒ;
μείζων ἄρα καὶ ἡ ΖΒ τῆς ΖΗ ἡ ἐλάττων τῆς μείζονος; ὅπερ
ἐστὶν ἀδύνατον. οὐκ ἄρα ἡ ΖΗ κάθετός ἐστιν ἐπὶ τὴν ΔΕ.
ὁμοίως δὴ δεῖχομεν, ὅτι οὐδ ἄλλη τις πλὴν τῆς ΖΓ;
ἡ ΖΓ ἄρα κάθετός ἐστιν ἐπὶ τὴν ΔΕ.}\\~\\~\\

\epsfysize=2.2in
\centerline{\epsffile{Book03/fig18g.eps}}

{\greekfont Ἐὰν ἄρα κύκλου ἐφάπτηταί τις εὐθεῖα, ἀπὸ δὲ τοῦ κέντρου ἐπὶ τὴν ἁφὴν ἐπιζευθθῇ τις εὐθεῖα, ἡ ἐπιζευθθεῖσα κάθετος ἔσται
ἐπὶ τὴν ἐφαπτομένην; ὅπερ ἔδει δεῖχαι.}}

\ParallelRText{
\begin{center}
{\large Proposition 18}
\end{center}

If some straight-line touches a circle, and some (other) straight-line is
joined from the center (of the circle) to the point of contact, (then) the (straight-line)
so joined will be perpendicular to the tangent.

For let some straight-line $DE$  touch the circle $ABC$  at  point $C$, and
let the center $F$ of circle $ABC$ be found [Prop.~3.1], 
and let $FC$ be joined from $F$ to $C$. I say that $FC$ is perpendicular
to $DE$.

For if not, let $FG$ be drawn from $F$, perpendicular to $DE$ [Prop.~1.12].

Therefore, since angle $FGC$ is a right-angle, (angle) $FCG$ is thus acute [Prop.~1.17]. And the greater angle is subtended by the greater side [Prop.~1.19]. Thus, $FC$ (is) greater than $FG$. And $FC$ (is) equal to
$FB$. Thus, $FB$ (is)  also greater than $FG$, the lesser than the greater.
The very thing is impossible. Thus, $FG$ is not perpendicular to
$DE$. So, similarly, we can show that neither (is) any other
(straight-line) except $FC$. Thus, $FC$ is perpendicular  to $DE$.

\epsfysize=2.2in
\centerline{\epsffile{Book03/fig18e.eps}}

Thus, if some straight-line  touches a circle, and some (other) straight-line is
joined from the center (of the circle) to the point of contact, (then) the (straight-line)
so joined will be perpendicular to the tangent. (Which is)
the very thing it was required to show.}
\end{Parallel}

%%%%%%
% Prop 3.19
%%%%%%
\pdfbookmark[1]{Proposition 3.19}{pdf3.19}
\begin{Parallel}{}{} 
\ParallelLText{
\begin{center}
{\large \ggn{19}.}
\end{center}\vspace*{-7pt}

{\greekfont Ἐὰν κύκλου ἐφάπτηταί τις εὐθεῖα, ἀπὸ δὲ τῆς ἁφῆς τῇ
ἐφαπτομένῃ πρὸς ὀρθὰς [γωνίας] εὐθεῖα γραμμὴ ἀθθῇ,
ἐπὶ τῆς ἀθθείσης ἔσται τὸ κέντρον τοῦ κύκλου.}\\

\epsfysize=2.2in
\centerline{\epsffile{Book03/fig19g.eps}}

{\greekfont Κύκλου γὰρ τοῦ ΑΒΓ ἐφαπτέσθω τις εὐθεῖα ἡ ΔΕ κατὰ τὸ Γ
σημεῖον, καὶ ἀπὸ τοῦ Γ τῇ ΔΕ πρὸς ὀρθὰς ἤθθω ἡ ΓΑ;
λέγω, ὅτι ἐπὶ τῆς ΑΓ ἐστι τὸ κέντρον τοῦ κύκλου.}

{\greekfont μὴ γάρ, ἀλλ εἰ δυνατόν, ἔστω τὸ Ζ, καὶ ἐπεζεύθθω ἡ ΓΖ.}

{\greekfont Ἐπεὶ [οὖν] κύκλου τοῦ ΑΒΓ ἐφάπτεταί τις εὐθεῖα ἡ ΔΕ,
ἀπὸ δὲ τοῦ κέντρου ἐπὶ τὴν ἁφὴν ἐπέζευκται ἡ ΖΓ, ἡ ΖΓ
ἄρα κάθετός ἐστιν ἐπὶ τὴν ΔΕ; ὀρθὴ ἄρα ἐστὶν ἡ ὑπὸ ΖΓΕ. ἐστὶ
δὲ καὶ ἡ ὑπὸ ΑΓΕ ὀρθή; ἴση ἄρα ἐστὶν ἡ ὑπὸ ΖΓΕ
τῇ ὑπὸ ΑΓΕ ἡ ἐλάττων τῇ μείζονι; ὅπερ ἐστὶν ἀδύνατον.
οὐκ ἄρα τὸ Ζ κέντρον ἐστὶ τοῦ ΑΒΓ κύκλου. ὁμοίως δὴ δείχομεν,
ὅτι οὐδ ἄλλο τι πλὴν ἐπὶ τῆς ΑΓ.}

{\greekfont Ἐὰν ἄρα κύκλου ἐφάπτηταί τις εὐθεῖα, ἀπὸ δὲ τῆς ἁφῆς τῇ
ἐφαπτομένῃ πρὸς ὀρθὰς εὐθεῖα γραμμὴ ἀθθῇ,
ἐπὶ τῆς ἀθθείσης ἔσται τὸ κέντρον τοῦ κύκλου; ὅπερ ἔδει
δεῖχαι.}}

\ParallelRText{
\begin{center}
{\large Proposition 19}
\end{center}

If some straight-line  touches a circle, and a straight-line is drawn from the
point of contact, at
right-[angles] to the tangent, (then) the center (of the circle) will be on the (straight-line) so drawn.

\epsfysize=2.2in
\centerline{\epsffile{Book03/fig19e.eps}}

For let some straight-line $DE$  touch the circle $ABC$  at point $C$. And let $CA$
be drawn from $C$, at right-angles to $DE$ [Prop.~1.11]. I say that the center of the circle is on $AC$.

For (if) not, if possible, let $F$ be (the center of the circle), and let
$CF$ be joined.

\mbox{[}Therefore], since some straight-line $DE$ touches the circle $ABC$, and
$FC$ has been joined from the center to the point of contact, $FC$ is thus
perpendicular to $DE$ [Prop.~3.18]. Thus, $FCE$ is a right-angle.
And $ACE$ is also a right-angle. Thus, $FCE$ is equal to $ACE$, the lesser to the
greater. The very thing is impossible. Thus, $F$ is not the center of circle $ABC$. 
So, similarly, we can show that neither is any (point) other (than one) on $AC$.

Thus, if some straight-line touches a circle, and a straight-line is drawn from the
point of contact, at
right-angles to the tangent, (then) the center (of the circle) will be on the (straight-line) so drawn. (Which is) the very thing it was required to show.}
\end{Parallel}

%%%%%%
% Prop 3.20
%%%%%%
\pdfbookmark[1]{Proposition 3.20}{pdf3.20}
\begin{Parallel}{}{} 
\ParallelLText{
\begin{center}
{\large \ggn{20}.}
\end{center}\vspace*{-7pt}

{\greekfont Ἐν κύκλῳ ἡ πρὸς τῷ κέντρῳ γωνία διπλασίων ἐστὶ τῆς πρὸς
τῇ περιφερείᾳ, ὅταν τὴν α ὐτὴν περιφέρειαν βάσιν ἔθωσιν αἱ
γωνίαι.}

{\greekfont Ἔστω κύκλος ὁ ΑΒΓ, καὶ πρὸς μὲν τῷ κέντρῳ α ὐτοῦ γωνία
ἔστω ἡ ὑπὸ ΒΕΓ, πρὸς δὲ τῇ περιφερείᾳ ἡ ὑπὸ ΒΑΓ, ἐθέτωσαν
δὲ τὴν α ὐτὴν περιφέρειαν βάσιν τὴν ΒΓ; λέγω, ὅτι διπλασίων ἐστὶν
ἡ ὑπὸ ΒΕΓ γωνία τῆς ὑπὸ ΒΑΓ.}

{\greekfont Ἐπιζευθθεῖσα γὰρ ἡ ΑΕ διήθθω ἐπὶ τὸ Ζ.}

{\greekfont Ἐπεὶ οὖν ἴση ἐστὶν ἡ ΕΑ τῇ ΕΒ, ἴση καὶ γωνία
ἡ ὑπὸ ΕΑΒ τῇ ὑπὸ ΕΒΑ; αἱ ἄρα ὑπὸ ΕΑΒ, ΕΒΑ γωνίαι
τῆς ὑπὸ ΕΑΒ διπλασίους εἰσίν. ἴση δὲ ἡ ὑπὸ ΒΕΖ ταῖς
ὑπὸ ΕΑΒ, ΕΒΑ; καὶ ἡ ὑπὸ ΒΕΖ ἄρα τῆς ὑπὸ ΕΑΒ ἐστι διπλῆ. διὰ τὰ α ὐτὰ δὴ καὶ ἡ ὑπὸ ΖΕΓ τῆς ὑπὸ ΕΑΓ ἐστι διπλῆ. ὅλη
ἄρα ἡ ὑπὸ ΒΕΓ ὅλης τῆς ὑπὸ ΒΑΓ ἐστι διπλῆ.}

\epsfysize=2.2in
\centerline{\epsffile{Book03/fig20g.eps}}

{\greekfont Κεκλάσθω δὴ πάλιν, καὶ ἔστω ἑτέρα γωνία ἡ ὑπὸ ΒΔΓ, καὶ
ἐπιζευθθεῖσα ἡ ΔΕ ἐκβεβλήσθω ἐπὶ τὸ Η. ὁμοίως
δὴ δείχομεν, ὅτι διπλῆ ἐστιν ἡ ὑπὸ ΗΕΓ γωνία τῆς ὑπὸ
ΕΔΓ, ὧν ἡ ὑπὸ ΗΕΒ διπλῆ ἐστι τῆς ὑπὸ ΕΔΒ; λοιπὴ
ἄρα ἡ ὑπὸ ΒΕΓ διπλῆ ἐστι τῆς ὑπὸ ΒΔΓ.}

{\greekfont Ἐν κύκλῳ ἄρα ἡ πρὸς τῷ κέντρῳ γωνία διπλασίων ἐστὶ τῆς πρὸς
τῇ περιφερείᾳ, ὅταν τὴν α ὐτὴν περιφέρειαν βάσιν ἔθωσιν [αἱ
γωνίαι]; ὅπερ ἔδει δεῖχαι.}}

\ParallelRText{
\begin{center}
{\large Proposition 20}
\end{center}

In a circle, the angle at the center is double that at the circumference, when
the angles have the same circumference base.

Let $ABC$ be a circle, and let $BEC$ be an angle at its center, and $BAC$
(one) at (its) circumference. And let them have the same circumference base
$BC$. I say that angle $BEC$ is double (angle) $BAC$.

For being joined, let $AE$ be drawn through to $F$.

Therefore, since $EA$ is equal to $EB$, angle $EAB$ (is) also equal to
$EBA$ [Prop.~1.5]. Thus, angle $EAB$ and $EBA$ is double (angle) $EAB$.
And $BEF$ (is) equal to $EAB$ and $EBA$ [Prop.~1.32]. Thus, $BEF$ is also
double $EAB$. So, for the same (reasons), $FEC$ is also double $EAC$. Thus,
the whole (angle) $BEC$ is double the whole (angle) $BAC$.

\epsfysize=2.2in
\centerline{\epsffile{Book03/fig20e.eps}}

So let another  (straight-line)  be inflected, and let there be another
angle, $BDC$. And $DE$ being joined, let it be produced to $G$.
So, similarly, we can show that angle $GEC$ is double $EDC$, of which
$GEB$ is double $EDB$. Thus, the remaining (angle) $BEC$ is double
the (remaining angle) $BDC$.

Thus, in a circle, the angle at the center is double that at the circumference, when
[the angles] have the same circumference base. (Which is) the very thing it
was required to show.\\}
\end{Parallel}

%%%%%%
% Prop 3.21
%%%%%%
\pdfbookmark[1]{Proposition 3.21}{pdf3.21}
\begin{Parallel}{}{} 
\ParallelLText{
\begin{center}
{\large \ggn{21}.}
\end{center}\vspace*{-7pt}

{\greekfont Ἐν κύκλῳ αἱ ἐν τῷ α ὐτῷ τμήματι γωνίαι ἴσαι ἀλλήλαις
εἰσίν.}

\epsfysize=2.2in
\centerline{\epsffile{Book03/fig21g.eps}}

{\greekfont Ἔστω κύκλος ὁ ΑΒΓΔ, καὶ ἐν τῷ α ὐτῷ τμήματι τῷ ΒΑΕΔ γωνίαι ἔστωσαν αἱ ὑπὸ ΒΑΔ, ΒΕΔ; λέγω, ὅτι αἱ ὑπὸ ΒΑΔ,
ΒΕΔ γωνίαι ἴσαι ἀλλήλαις εἰσίν.}

{\greekfont Εἰλήφθω γὰρ τοῦ ΑΒΓΔ κύκλου τὸ κέντρον, καὶ ἔστω τὸ Ζ, καὶ
ἐπεζεύθθωσαν αἱ ΒΖ, ΖΔ.}

{\greekfont Καὶ ἐπεὶ ἡ μὲν ὑπὸ ΒΖΔ γωνία πρὸς τῷ κέντρῳ ἐστίν,
ἡ δὲ ὑπὸ ΒΑΔ πρὸς τῇ περιφερείᾳ, καὶ ἔθουσι τὴν α ὐτὴν
περιφέρειαν βάσιν τὴν ΒΓΔ, ἡ ἄρα ὑπὸ ΒΖΔ γωνία
διπλασίων ἐστὶ τῆς ὑπὸ ΒΑΔ. διὰ τὰ α ὐτὰ δὴ ἡ ὑπὸ ΒΖΔ καὶ τῆς ὑπὸ ΒΕΔ ἐστι διπλσίων; ἴση ἄρα ἡ ὑπὸ ΒΑΔ τῇ ὑπὸ
ΒΕΔ.}

{\greekfont Ἐν κύκλῳ ἄρα αἱ ἐν τῷ α ὐτῷ τμήματι γωνίαι ἴσαι ἀλλήλαις
εἰσίν; ὅπερ ἔδει δεῖχαι.}}

\ParallelRText{
\begin{center}
{\large Proposition 21}
\end{center}

In a circle, angles in the same segment are equal to one another.

\epsfysize=2.2in
\centerline{\epsffile{Book03/fig21e.eps}}

Let $ABCD$ be a circle, and let $BAD$ and $BED$ be angles in the same segment
$BAED$. I say that angles $BAD$ and $BED$ are equal to one another.

For let the center of circle $ABCD$ be found [Prop.~3.1],
and let it be (at point) $F$. And let $BF$ and $FD$ be joined.

And since angle $BFD$ is at the center, and $BAD$ at the
circumference, and they have the same circumference base $BCD$, angle
$BFD$ is thus double $BAD$ [Prop.~3.20]. So, for the
same (reasons), $BFD$ is also double  $BED$. Thus, $BAD$ (is) equal to $BED$.

Thus, in a circle, angles in the same segment are equal to one another.
(Which is) the very thing it was required to show.}
\end{Parallel}

%%%%%%
% Prop 3.22
%%%%%%
\pdfbookmark[1]{Proposition 3.22}{pdf3.22}
\begin{Parallel}{}{} 
\ParallelLText{
\begin{center}
{\large \ggn{22}.}
\end{center}\vspace*{-7pt}

{\greekfont Τῶν ἐν τοῖς κύκλοις τετραπλεύρων αἱ ἀπεναντίον γωνίαι δυσὶν ὀρθαῖς ἴσαι εἰσίν.}

\epsfysize=2.2in
\centerline{\epsffile{Book03/fig22g.eps}}

{\greekfont Ἔστω κύκλος ὁ ΑΒΓΔ, καὶ ἐν α ὐτῷ τετράπλευρον ἔστω τὸ ΑΒΓΔ;
λέγω, ὅτι αἱ ἀπεναντίον γωνίαι δυσὶν ὀρθαῖς ἴσαι εἰσίν.}

{\greekfont Ἐπεζεύθθωσαν αἱ ΑΓ, ΒΔ.}

{\greekfont Ἐπεὶ οὖν παντὸς τριγώνου αἱ τρεῖς γωνίαι δυσὶν ὀρθαῖς ἴσαι
εἰσίν, τοῦ ΑΒΓ ἄρα τριγώνου αἱ τρεῖς γωνίαι αἱ ὑπὸ ΓΑΒ, ΑΒΓ, ΒΓΑ δυσὶν ὀρθαῖς ἴσαι εἰσίν. ἴση δὲ ἡ μὲν ὑπὸ ΓΑΒ τῇ ὑπὸ
ΒΔΓ; ἐν γὰρ τῷ α ὐτῷ τμήματί εἰσι τῷ ΒΑΔΓ; ἡ δὲ ὑπὸ ΑΓΒ
τῇ ὑπὸ ΑΔΒ; ἐν γὰρ τῷ α ὐτῷ τμήματί εἰσι τῷ ΑΔΓΒ;
ὅλη ἄρα ἡ ὑπὸ ΑΔΓ ταῖς ὑπὸ ΒΑΓ, ΑΓΒ ἴση ἐστίν.
κοινὴ προσκείσθω ἡ ὑπὸ ΑΒΓ; αἱ ἄρα ὑπὸ ΑΒΓ, ΒΑΓ, ΑΓΒ ταῖς
ὑπὸ ΑΒΓ, ΑΔΓ ἴσαι εἰσίν. ἀλλ αἱ ὑπὸ ΑΒΓ, ΒΑΓ, ΑΓΒ δυσὶν ὀρθαῖς ἴσαι εἰσίν. καὶ αἱ ὑπὸ ΑΒΓ, ΑΔΓ ἄρα δυσὶν ὀρθαῖς
ἴσαι εἰσίν. ὁμοίως δὴ δείχομεν, ὅτι καὶ αἱ ὑπὸ ΒΑΔ, ΔΓΒ γωνίαι δυσὶν ὀρθαῖς ἴσαι εἰσίν.}

{\greekfont Τῶν ἄρα ἐν τοῖς κύκλοις τετραπλεύρων αἱ ἀπεναντίον γωνίαι δυσὶν ὀρθαῖς ἴσαι εἰσίν; ὅπερ ἔδει δεῖχαι.}}

\ParallelRText{
\begin{center}
{\large Proposition 22}
\end{center}

For quadrilaterals within circles, the (sum of the) opposite angles is equal to two right-angles.

\epsfysize=2.2in
\centerline{\epsffile{Book03/fig22e.eps}}

Let $ABCD$ be a circle, and let $ABCD$ be a quadrilateral within it. I say that
the (sum of the) opposite angles is equal to two right-angles.

Let $AC$ and $BD$ be joined.

Therefore, since the three angles of any triangle are equal to two
right-angles [Prop.~1.32], the
three angles $CAB$, $ABC$, and $BCA$ of triangle $ABC$ are thus equal to two right-angles. And $CAB$ (is) equal to $BDC$. For they are in the
same segment $BADC$ [Prop.~3.21]. And $ACB$ (is equal) to $ADB$.
For they are in the same segment $ADCB$ [Prop.~3.21].
Thus, the whole of $ADC$ is equal to $BAC$ and $ACB$. Let $ABC$ be added to both. Thus, $ABC$, $BAC$, and
$ACB$ are equal to $ABC$ and $ADC$. But, $ABC$, $BAC$, and $ACB$ are equal to
two right-angles. Thus, $ABC$ and $ADC$ are also equal to two right-angles.
Similarly, we can show that angles $BAD$ and $DCB$ are also equal to
two right-angles.

Thus, for quadrilaterals within circles, the (sum of the) opposite angles is equal to two right-angles. (Which is) the very thing it was required to show.}
\end{Parallel}

%%%%%%
% Prop 3.23
%%%%%%
\pdfbookmark[1]{Proposition 3.23}{pdf3.23}
\begin{Parallel}{}{} 
\ParallelLText{
\begin{center}
{\large \ggn{23}.}
\end{center}\vspace*{-7pt}

{\greekfont Ἐπὶ τῆς α ὐτῆς εὐθείας δύο τμήματα κύκλων ὅμοια καὶ ἄνισα
οὐ συσταθήσεται ἐπὶ τὰ α ὐτὰ μέρη.}

{\greekfont Εἰ γὰρ δυνατόν, ἐπὶ τῆς α ὐτῆς εὐθείας τῆς ΑΒ δύο τμήματα
κύκλων ὅμοια καὶ ἄνισα συνεστάτω ἐπὶ τὰ α ὐτὰ μέρη τὰ ΑΓΒ, ΑΔΒ,
καὶ διήθθω ἡ ΑΓΔ, καὶ ἐπεζεύθθωσαν αἱ ΓΒ, ΔΒ.}

\epsfysize=1.25in
\centerline{\epsffile{Book03/fig23g.eps}}

{\greekfont Ἐπεὶ οὖν ὅμοιόν ἐστι τὸ ΑΓΒ τμῆμα τῷ ΑΔΒ τμήματι,
ὅμοια δὲ τμήματα κύκλων ἐστὶ τὰ δεθόμενα γωνίας ἴσας, ἴση
ἄρα ἐστὶν ἡ ὑπὸ ΑΓΒ γωνία τῇ ὑπὸ ΑΔΒ ἡ ἐκτὸς τῇ ἐντός;
ὅπερ ἐστὶν ἀδύνατον.}

{\greekfont Οὐκ ἄρα ἐπὶ τῆς α ὐτῆς εὐθείας δύο τμήματα κύκλων ὅμοια
καὶ ἄνισα συσταθήσεται ἐπὶ τὰ α ὐτὰ μέρη; ὅπερ ἔδει δεῖχαι.}}

\ParallelRText{
\begin{center}
{\large Proposition 23}
\end{center}

Two similar and unequal segments of circles cannot be constructed on the same side of the same straight-line.

For, if possible, let the two similar and unequal segments of circles, $ACB$ and
$ADB$, be constructed on the same side of the same straight-line $AB$. And let $ACD$ be drawn through (the segments), and let $CB$ and $DB$ be joined.

\epsfysize=1.25in
\centerline{\epsffile{Book03/fig23e.eps}}

Therefore, since segment $ACB$ is similar to segment $ADB$, and similar segments of circles are those accepting equal angles [Def.~3.11], 
angle $ACB$ is thus equal to $ADB$, the external to the internal. The very thing
is impossible [Prop.~1.16].

Thus,  two similar and unequal segments of circles cannot be constructed on the same side of the same straight-line.}
\end{Parallel}

%%%%%%
% Prop 3.24
%%%%%%
\pdfbookmark[1]{Proposition 3.24}{pdf3.24}
\begin{Parallel}{}{} 
\ParallelLText{
\begin{center}
{\large\ggn{24}.}
\end{center}\vspace*{-7pt}

{\greekfont Τὰ ἐπὶ ἴσων εὐθειῶν ὅμοια τμήματα κύλων ἴσα ἀλλήλοις ἐστίν.}

{\greekfont Ἔστωσαν γὰρ ἐπὶ ἴσων εὐθειῶν τῶν ΑΒ, ΓΔ ὅμοια
τμήματα κύκλων τὰ ΑΕΒ, ΓΖΔ; λέγω, ὅτι ἴσον ἐστὶ τὸ ΑΕΒ
τμῆμα τῷ ΓΖΔ τμήματι.}

{\greekfont Ἐφαρμοζομένου γὰρ τοῦ ΑΕΒ τμήματος ἐπὶ τὸ ΓΖΔ καὶ τιθεμένου τοῦ μὲν Α σημείου ἐπὶ τὸ Γ τῆς δὲ ΑΒ εὐθείας ἐπὶ τὴν
ΓΔ, ἐφαρμόσει καὶ τὸ Β σημεῖον ἐπὶ τὸ Δ σημεῖον διὰ τὸ ἴσην
εἶναι τὴν ΑΒ τῇ ΓΔ; τῆς δὲ ΑΒ ἐπὶ τὴν ΓΔ ἐφαρμοσάσης ἐφαρμόσει καὶ τὸ ΑΕΒ τμῆμα ἐπὶ τὸ ΓΖΔ. εἰ γὰρ ἡ ΑΒ εὐθεῖα
ἐπὶ τὴν ΓΔ ἐφαρμόσει, τὸ δὲ ΑΕΒ τμῆμα ἐπὶ τὸ ΓΖΔ μὴ ἐφαρμόσει, ἤτοι ἐντὸς α ὐτοῦ πεσεῖται ἢ ἐκτὸς ἢ παραλλάχει,
ὡς τὸ ΓΗΔ, καὶ κύκλος κύκλον τέμνει κατὰ πλείονα σημεῖα ἢ
δύο; ὅπερ ἐστίν ἀδύνατον. οὐκ ἄρα ἐφαρμοζομένης τῆς ΑΒ
εὐθείας ἐπὶ τὴν ΓΔ οὐκ ἐφαρμόσει καὶ τὸ ΑΕΒ τμῆμα ἐπὶ
τὸ ΓΖΔ; ἐφαρμόσει ἄρα, καὶ ἴσον α ὐτῷ ἔσται.}\\~\\

\epsfysize=2.4in
\centerline{\epsffile{Book03/fig24g.eps}}

{\greekfont Τὰ ἄρα ἐπὶ ἴσων εὐθειῶν ὅμοια τμήματα κύκλων
ἴσα ἀλλήλοις ἐστίν; ὅπερ ἔδει δεῖχαι.}}

\ParallelRText{
\begin{center}
{\large Proposition 24}
\end{center}

Similar segments of circles on equal straight-lines are equal to one another.

For let $AEB$ and $CFD$ be similar segments of circles on the equal straight-lines $AB$ and $CD$ (respectively). I say that segment $AEB$ is equal to segment
$CFD$.

For if the segment $AEB$ is applied to the segment $CFD$,  and point $A$
is placed on (point) $C$, and the straight-line $AB$ on $CD$, (then) point $B$ will also coincide with point $D$, on account of $AB$ being equal to $CD$. And
if $AB$ coincides with $CD$, (then) the segment $AEB$ will also coincide with $CFD$.
For if the straight-line $AB$ coincides with $CD$, and the segment $AEB$ does not
coincide with $CFD$, (then) it will surely either fall inside it, outside (it),$^\dag$ or
it will miss like $CGD$ (in the figure), and a circle (will) cut (another) circle
at more than two points. The very thing is impossible [Prop.~3.10].
Thus, if the straight-line $AB$ is applied to $CD$, the segment $AEB$ cannot
not also coincide with $CFD$. Thus, it will coincide, and will be equal to it [C.N.~4].\\

\epsfysize=2.4in
\centerline{\epsffile{Book03/fig24e.eps}}

Thus, similar segments of circles on equal straight-lines are equal to one another. (Which is) the very thing it was required to show.}
\end{Parallel}
{\footnotesize \noindent$^\dag$ Both this possibility, and the previous one, are precluded by Prop.~3.23.}

%%%%%%
% Prop 3.25
%%%%%%
\pdfbookmark[1]{Proposition 3.25}{pdf3.25}
\begin{Parallel}{}{} 
\ParallelLText{
\begin{center}
{\large \ggn{25}.}
\end{center}\vspace*{-7pt}

{\greekfont Κύκλου τμήματος δοθέντος προσαναγράψαι τὸν κύκλον, οὗπέρ ἐστι τμῆμα.}

{\greekfont Ἔστω τὸ δοθὲν τμῆμα κύκλου τὸ ΑΒΓ; δεῖ δὴ τοῦ ΑΒΓ τμήματος
προσαναγράψαι τὸν κύκλον, οὖπέρ ἐστι τμῆμα.}

{\greekfont Τετμήσθω γὰρ ἡ ΑΓ δίθα κατὰ τὸ Δ, καὶ ἤθθω ἀπὸ τοῦ Δ
σημείου τῇ ΑΓ πρὸς ὀρθὰς ἡ ΔΒ, καὶ ἐπεζεύθθω ἡ ΑΒ; ἡ ὑπὸ
ΑΒΔ γωνία ἄρα τῆς ὑπὸ ΒΑΔ ἤτοι μείζων ἐστὶν ἢ ἴση ἢ
ἐλάττων.}

{\greekfont Ἔστω πρότερον μείζων, καὶ συνεστάτω πρὸς τῇ ΒΑ εὐθείᾳ καὶ
τῷ πρὸς α ὐτῇ σημείῳ τῷ Α τῇ ὑπὸ ΑΒΔ γωνίᾳ ἴση ἡ ὑπὸ
ΒΑΕ, καὶ διήθθω ἡ ΔΒ ἐπὶ τὸ Ε, καὶ ἐπεζεύθθω ἡ ΕΓ. ἐπεὶ οὖν
ἴση ἐστὶν ἡ ὑπὸ ΑΒΕ γωνία τῇ ὑπὸ ΒΑΕ, ἴση ἄρα ἐστὶ
καὶ ἡ ΕΒ εὐθεῖα τῇ ΕΑ. καὶ ἐπεὶ ἴση ἐστὶν ἡ ΑΔ τῇ
ΔΓ, κοινὴ δὲ ἡ ΔΕ, δύο δὴ αἱ ΑΔ, ΔΕ δύο ταῖς ΓΔ, ΔΕ ἴσαι
εἰσὶν ἑκατέρα ἑκατέρᾳ; καὶ γωνία ἡ ὑπὸ ΑΔΕ γωνίᾳ τῇ ὑπὸ
ΓΔΕ ἐστιν ἴση; ὀρθὴ γὰρ ἑκατέρα; βάσις ἄρα ἡ ΑΕ βάσει τῇ ΓΕ
ἐστιν ἴση. ἀλλὰ ἡ ΑΕ τῇ ΒΕ ἐδείθθη ἴση; καὶ ἡ ΒΕ ἄρα
τῇ ΓΕ ἐστιν ἴση; αἱ τρεῖς ἄρα αἱ ΑΕ, ΕΒ, ΕΓ ἴσαι
ἀλλήλαις εἰσίν; ὁ ἄρα κέντρῷ τῷ Ε διαστήματι δὲ ἑνὶ τῶν ΑΕ, ΕΒ, ΕΓ κύκλος γραφόμενος ἥχει καὶ διὰ τῶν λοιπῶν σημείων
καὶ ἔσται προσαναγεγραμμένος. κύκλου ἄρα τμήματος δοθέντος
προσαναγέγραπται ὁ κύκλος. καὶ δῆλον, ὡς τὸ ΑΒΓ τμῆμα ἔλαττόν
ἐστιν ἡμικυκλίου διὰ τὸ τὸ Ε κέντρον ἐκτὸς α ὐτοῦ τυγθάνειν.}\\~\\~\\

\epsfysize=1.3in
\centerline{\epsffile{Book03/fig25g.eps}}

{\greekfont Ὁμοίως [δὲ] κἂν ᾖ ἡ ὑπὸ ΑΒΔ γωνία ἴση τῇ ὑπὸ
ΒΑΔ, τῆς ΑΔ ἴσης γενομένης ἑκατέρᾳ τῶν ΒΔ, ΔΓ αἱ
τρεῖς αἱ ΔΑ, ΔΒ, ΔΓ ἴσαι ἀλλήλαις ἔσονται, καὶ ἔσται τὸ Δ κέντρον
τοῦ προσαναπεπληρωμένου κύκλου, καὶ δηλαδὴ ἔσται τὸ ΑΒΓ
ἡμικύκλιον.}

{\greekfont Ἐὰν δὲ ἡ ὑπὸ ΑΒΔ ἐλάττων ᾖ τῆς ὑπὸ ΒΑΔ,
καὶ συστησώμεθα πρὸς τῇ ΒΑ εὐθείᾳ καὶ τῷ πρὸς
α ὐτῇ σημείῳ τῷ Α τῇ ὑπὸ ΑΒΔ γωνίᾳ ἴσην, ἐντὸς
τοῦ ΑΒΓ τμήματος πεσεῖται τὸ κέντρον ἐπὶ τῆς ΔΒ, καὶ
ἔσται δηλαδὴ τὸ ΑΒΓ τμῆμα μεῖζον ἡμικυκλίου.}

{\greekfont Κύκλου ἄρα τμήματος δοθέντος προσαναγέγραπται ὁ κύκλος;
ὅπερ ἔδει ποιῆσαι.}}

\ParallelRText{
\begin{center}
{\large Proposition 25}
\end{center}

For a given segment of a circle, to complete the circle, the very one of which it is a segment.

Let $ABC$ be the given segment of a circle. So it is required
to complete the circle for segment $ABC$, the very one of which it is a segment.

For let $AC$ be cut in half at (point) $D$ [Prop.~1.10], and let $DB$ be drawn
from point $D$, at right-angles to $AC$ [Prop.~1.11]. And let $AB$ be joined.
Thus, angle $ABD$ is surely either greater than,  equal to, or less than (angle)
$BAD$.

First of all, let it be greater. And let (angle) $BAE$, equal to
angle $ABD$, be constructed on the straight-line $BA$, at the point $A$ on it [Prop.~1.23].
And let $DB$ be drawn through to $E$, and let $EC$ be joined.
Therefore, since angle $ABE$ is equal to $BAE$, the straight-line $EB$ is thus also
equal to $EA$ [Prop.~1.6]. And since $AD$ is equal to $DC$, and $DE$ (is)
common, the two (straight-lines) $AD$, $DE$ are equal to the two (straight-lines)
$CD$, $DE$, respectively. And angle $ADE$ is equal to angle $CDE$. For each (is)
a right-angle. Thus, the base $AE$ is equal to the base $CE$ [Prop.~1.4].
But, $AE$ was shown (to be) equal to $BE$. Thus, $BE$ is also equal
to $CE$. Thus, the three (straight-lines) $AE$, $EB$, and $EC$ are equal to
one another. Thus, if a circle is drawn with center $E$, and radius one of $AE$, $EB$, or $EC$,  it will also go through the remaining points (of the segment), and the (associated circle) will be
 completed [Prop.~3.9]. Thus, a circle has been completed from the given
 segment of a circle. And (it is) clear that the segment $ABC$ is less
 than a semi-circle, because the center $E$ happens to lie outside it.

\epsfysize=1.3in
\centerline{\epsffile{Book03/fig25e.eps}}
 
\mbox{[}And], similarly, even if angle $ABD$ is equal to $BAD$,  (since) $AD$ becomes equal to
 each of $BD$ [Prop.~1.6] and $DC$, the three (straight-lines) $DA$, $DB$, and $DC$
 will  be equal to one another. And point $D$ will be the center of the completed
 circle. And $ABC$ will manifestly be a semi-circle.
 
 And if $ABD$ is less than $BAD$, and we construct (angle $BAE$), equal to angle
 $ABD$,  on the straight-line $BA$, at the point $A$ on it [Prop.~1.23], then
 the center will fall on $DB$,  inside the segment $ABC$. And segment
 $ABC$ will manifestly be greater than a semi-circle.
 
 Thus, a circle has been completed from the given segment of a circle. (Which is) the very thing it was required to do.}
\end{Parallel}

%%%%%%
% Prop 3.26
%%%%%%
\pdfbookmark[1]{Proposition 3.26}{pdf3.26}
\begin{Parallel}{}{} 
\ParallelLText{
\begin{center}
{\large \ggn{26}.}
\end{center}\vspace*{-7pt}

{\greekfont Ἐν τοῖς ἴσοις κύκλοις αἱ ἴσαι γωνίαι ἐπὶ ἴσων περιφερειῶν βεβήκασιν, ἐάν τε πρὸς τοῖς κέντροις ἐάν τε πρὸς ταῖς περιφερείαις
ὦσι βεβηκυῖαι.}

{\greekfont Ἔστωσαν ἴσοι κύκλοι οἱ ΑΒΓ, ΔΕΖ καὶ ἐν α ὐτοῖς ἴσαι γωνίαι ἔστωσαν πρὸς μὲν τοῖς κέντροις αἱ ὑπὸ ΒΗΓ, ΕΘΖ, πρὸς δὲ ταῖς
περιφερείαις αἱ ὑπὸ ΒΑΓ, ΕΔΖ; λέγω, ὅτι ἴση ἐστὶν ἡ ΒΚΓ
περιφέρεια τῇ ΕΛΖ περιφερείᾳ.}

{\greekfont Ἐπεζεύθθωσαν γὰρ αἱ ΒΓ, ΕΖ.}

{\greekfont Καὶ ἐπεὶ ἴσοι εἰσὶν οἱ ΑΒΓ, ΔΕΖ κύκλοι, ἴσαι
εἰσὶν αἱ ἐκ τῶν κέντρων; δύο δὴ αἱ ΒΗ, ΗΓ
δύο ταῖς ΕΘ, ΘΖ ἴσαι; καὶ γωνία ἡ πρὸς τῷ Η γωνίᾳ τῇ πρὸς
τῷ Θ ἴση; βάσις ἄρα ἡ ΒΓ βάσει τῇ ΕΖ ἐστιν ἴση. καὶ ἐπεὶ
ἴση ἐστὶν ἡ πρὸς τῷ Α γωνία τῇ πρὸς τῷ Δ, ὅμοιον ἄρα ἐστὶ τὸ ΒΑΓ τμῆμα τῷ ΕΔΖ τμήματι; καί εἰσιν ἐπὶ ἴσων εὐθειῶν [τῶν ΒΓ, ΕΖ]; τὰ δὲ ἐπὶ ἴσων εὐθειῶν ὅμοια
τμήματα κύκλων ἴσα ἀλλήλοις ἐστίν; ἴσον ἄρα τὸ ΒΑΓ τμῆμα τῷ
ΕΔΖ. ἔστι δὲ καὶ ὅλος ὁ ΑΒΓ κύκλος ὅλῳ τῷ ΔΕΖ κύκλῳ
ἴσος; λοιπὴ ἄρα ἡ ΒΚΓ περιφέρεια τῇ ΕΛΖ περιφερείᾳ ἐστὶν ἴση.}\\~\\~\\

\epsfysize=1.5in
\centerline{\epsffile{Book03/fig26g.eps}}

{\greekfont Ἐν ἄρα τοῖς ἴσοις κύκλοις αἱ ἴσαι γωνίαι ἐπὶ ἴσων
περιφερειῶν βεβήκασιν, ἐάν τε πρὸς τοῖς κέντροις ἐάν
τε πρὸς ταῖς περιφερείας ὦσι βεβηκυῖαι; ὅπερ ἔδει
δεῖχαι.}}

\ParallelRText{
\begin{center}
{\large Proposition 26}
\end{center}

In equal circles, equal angles stand upon equal circumferences  whether they
are standing at the center or at the circumference.

Let $ABC$ and $DEF$ be equal circles, and within them let $BGC$ and $EHF$ be equal angles at the center, and $BAC$ and $EDF$ (equal angles) at the circumference. I say that circumference $BKC$ is equal to circumference $ELF$.

For let $BC$ and $EF$ be joined.

And since circles $ABC$ and $DEF$ are equal, their radii are equal.
So the two (straight-lines) $BG$, $GC$ (are) equal to the two (straight-lines)
$EH$, $HF$ (respectively). And the angle at $G$ (is) equal to the angle at $H$.
Thus, the base $BC$ is equal to the base $EF$ [Prop.~1.4].
And since the angle at $A$ is equal to the (angle) at $D$, the segment $BAC$ is
thus similar to the segment $EDF$ [Def.~3.11]. And they are on
equal straight-lines [$BC$ and $EF$]. And similar segments of circles on equal
straight-lines are equal to one another [Prop.~3.24]. Thus, segment
$BAC$ is equal to (segment) $EDF$. And the whole circle $ABC$ is also
equal to the whole circle $DEF$. Thus, the remaining circumference $BKC$
is equal to the (remaining) circumference $ELF$.

\epsfysize=1.5in
\centerline{\epsffile{Book03/fig26e.eps}}

Thus,   in equal circles, equal angles stand upon equal circumferences, whether they
are standing  at the center or at the circumference. (Which is)
the very thing which it was required to show.}
\end{Parallel}

%%%%%%
% Prop 3.27
%%%%%%
\pdfbookmark[1]{Proposition 3.27}{pdf3.27}
\begin{Parallel}{}{} 
\ParallelLText{
\begin{center}
{\large\ggn{27}.}
\end{center}\vspace*{-7pt}

{\greekfont Ἐν τοῖς ἴσοις κύκλοις αἱ ἐπὶ ἴσων περιφερειῶν βεβηκυῖαι γωνίαι
ἴσαι ἀλλήλαις εἰσίν, ἐάν τε πρὸς τοῖς κέντροις ἐάν τε πρὸς
ταῖς περιφερείαις ὦσι βεβηκυῖαι.}

\epsfysize=1.5in
\centerline{\epsffile{Book03/fig27g.eps}}

{\greekfont Ἐν γὰρ ἴσοις κύκλοις τοῖς ΑΒΓ, ΔΕΖ ἐπὶ ἴσων περιφερειῶν
τῶν ΒΓ, ΕΖ πρὸς μὲν τοῖς Η, Θ κέντροις γωνίαι βεβηκέτωσαν αἱ
ὑπὸ ΒΗΓ, ΕΘΖ, πρὸς δὲ ταῖς περιφερείαις αἱ ὑπὸ ΒΑΓ, ΕΔΖ;
λέγω, ὅτι ἡ μὲν ὑπὸ ΒΗΓ γωνία τῇ ὑπὸ ΕΘΖ ἐστιν ἴση,
ἡ δὲ ὑπὸ ΒΑΓ τῇ ὑπὸ ΕΔΖ ἐστιν ἴση.}

{\greekfont Εἰ γὰρ ἄνισός ἐστιν ἡ ὑπὸ ΒΗΓ τῇ ὑπὸ ΕΘΖ,
μία α ὐτῶν μείζων ἐστίν. ἔστω μείζων ἡ ὑπὸ ΒΗΓ, καὶ
συνεστάτω πρὸς τῇ ΒΗ εὐθείᾳ καὶ τῷ πρὸς α ὐτῇ σημείῳ
τῷ Η τῇ ὑπὸ ΕΘΖ γωνίᾳ ἴση ἡ ὑπὸ ΒΗΚ; αἱ δὲ ἴσαι
γωνίαι ἐπὶ ἴσων περιφερειῶν βεβήκασιν, ὅταν πρὸς τοῖς κέντροις
ὦσιν; ἴση ἄρα ἡ ΒΚ περιφέρεια τῇ ΕΖ περιφερείᾳ. ἀλλὰ ἡ ΕΖ τῇ
ΒΓ ἐστιν ἴση; καὶ ἡ ΒΚ ἄρα τῇ ΒΓ ἐστιν ἴση ἡ ἐλάττων τῇ
μείζονι; ὅπερ ἐστὶν ἀδύνατον. οὐκ ἄρα ἄνισός ἐστιν ἡ ὑπὸ
ΒΗΓ γωνία τῇ ὑπὸ ΕΘΖ; ἴση ἄρα. καί ἐστι τῆς μὲν ὑπὸ ΒΗΓ
ἡμίσεια ἡ πρὸς τῷ Α, τῆς δὲ ὑπὸ  ΕΘΖ ἡμίσεια
ἡ πρὸς τῷ Δ; ἴση ἄρα καὶ ἡ πρὸς τῷ Α γωνία τῇ πρὸς
τῷ Δ.}

{\greekfont Ἐν ἄρα τοῖς ἴσοις κύκλοις αἱ ἐπὶ ἴσων περιφερειῶν βεβηκυῖαι
γωνίαι ἴσαι ἀλλήλαις εἰσίν, ἐάν τε πρὸς τοῖς κέντροις ἐάν
τε πρὸς ταῖς περιφερείαις ὦσι βεβηκυῖαι; ὅπερ ἔδει δεῖχαι.}}

\ParallelRText{
\begin{center}
{\large Proposition 27}
\end{center}

In equal circles, angles standing upon equal circumferences  are equal to one
another, whether they are standing at the center or at the circumference.

\epsfysize=1.5in
\centerline{\epsffile{Book03/fig27e.eps}}

For let the angles $BGC$ and $EHF$ at the centers $G$ and $H$, and the (angles) $BAC$ and $EDF$
at the circumferences, stand upon the equal circumferences $BC$ and $EF$, in the equal circles $ABC$ and $DEF$ (respectively). I say that angle $BGC$ is
equal to (angle) $EHF$, and $BAC$ is equal to $EDF$.

For if $BGC$ is unequal to $EHF$, one of them is greater. Let $BGC$ be greater, 
and let the (angle) $BGK$, equal to  angle $EHF$, be constructed on the
straight-line $BG$, at the point $G$ on it [Prop.~1.23]. But equal angles (in equal circles) stand upon equal circumferences, when they are at the centers [Prop.~3.26]. Thus,
circumference $BK$ (is) equal to circumference $EF$. But, $EF$ is equal to $BC$.
Thus, $BK$ is also equal to $BC$, the lesser to the greater. The very thing is
impossible. Thus, angle $BGC$ is not unequal to $EHF$. Thus, (it is) equal.
And the (angle) at $A$ is half $BGC$, and the (angle) at $D$ half $EHF$ [Prop.~3.20]. Thus, the angle at $A$ (is) also equal to the (angle) at $D$.

Thus, in equal circles, angles standing upon equal circumferences  are equal to one
another, whether they are standing at the center or at the circumference.
(Which is) the very thing it was required to show.\\}
\end{Parallel}

%%%%%%
% Prop 3.28
%%%%%%
\pdfbookmark[1]{Proposition 3.28}{pdf3.28}
\begin{Parallel}{}{} 
\ParallelLText{
\begin{center}
{\large\ggn{28}.}
\end{center}\vspace*{-7pt}

{\greekfont Ἐν τοῖς ἴσοις κύκλοις αἱ ἴσαι εὐθεῖαι ἴσας περιφερείας ἀφαιροῦσι
τὴν μὲν μείζονα  τῇ μείζονι τὴν δὲ ἐλάττονα τῇ ἐλάττονι.}

{\greekfont Ἔστωσαν ἴσοι κύκλοι οἱ ΑΒΓ, ΔΕΖ, καὶ ἐν τοῖς κύκλοις ἴσαι
εὐθεῖαι ἔστωσαν αἱ ΑΒ, ΔΕ τὰς μὲν ΑΓΒ, ΑΖΕ περιφερείας μείζονας
ἀφαιροῦσαι τὰς δὲ ΑΗΒ, ΔΘΕ ἐλάττονας; λέγω, ὅτι ἡ μὲν ΑΓΒ
μείζων περιφέρεια ἴση ἐστὶ τῇ ΔΖΕ μείζονι περιφερείᾳ
ἡ δὲ ΑΗΒ ἐλάττων περιφέρεια τῇ ΔΘΕ.}\\

\epsfysize=1.6in
\centerline{\epsffile{Book03/fig28g.eps}}

{\greekfont Εἰλήφθω γὰρ τὰ κέντρα τῶν κύκλων τὰ Κ, Λ, καὶ ἐπεζεύθθωσαν αἱ
ΑΚ, ΚΒ, ΔΛ, ΛΕ.}

{\greekfont Καὶ ἐπεὶ ἴσοι κύκλοι εἰσίν, ἴσαι εἰσὶ καὶ αἱ ἐκ τῶν κέντρων;
δύο δὴ αἱ ΑΚ, ΚΒ δυσὶ ταῖς ΔΛ, ΛΕ
ἴσαι εἰσίν; καὶ βάσις ἡ ΑΒ βάσει τῇ ΔΕ ἴση; γωνία ἄρα ἡ ὑπὸ
ΑΚΒ γωνίᾳ τῇ ὑπὸ ΔΛΕ ἴση ἐστίν. αἱ δὲ ἴσαι γωνίαι ἐπὶ
ἴσων περιφερειῶν βεβήκασιν, ὅταν πρὸς τοῖς κέντροις ὦσιν; ἴση
ἄρα ἡ ΑΗΒ περιφέρεια τῇ ΔΘΕ. ἐστὶ δὲ καὶ ὅλος ὁ ΑΒΓ
κύκλος ὅλῳ τῷ ΔΕΖ κύκλῳ ἴσος; καὶ λοιπὴ ἄρα ἡ ΑΓΒ 
περιφέρεια λοιπῇ τῇ ΔΖΕ περιφερείᾳ ἴση ἐστίν.}

{\greekfont Ἐν ἄρα τοῖς ἴσοις κύκλοις αἱ ἴσαι εὐθεῖαι ἴσας περιφερείας ἀφαιροῦσι
τὴν μὲν μείζονα τῇ μείζονι τὴν δὲ ἐλάττονα τῇ ἐλάττονι; ὅπερ
ἔδει δεῖχαι.}}

\ParallelRText{
\begin{center}
{\large Proposition 28}
\end{center}

 In equal circles, equal straight-lines cut off equal circumferences, 
the greater (circumference being equal) to the greater, and the lesser to the lesser.

Let $ABC$ and $DEF$ be equal circles, and let $AB$ and $DE$ be equal straight-lines
in these circles, cutting off the greater circumferences $ACB$ and $DFE$, 
and the lesser (circumferences) $AGB$ and $DHE$ (respectively). I say that
the greater circumference $ACB$ is equal to the greater circumference $DFE$,
and the lesser circumference $AGB$ to (the lesser) $DHE$.

\epsfysize=1.6in
\centerline{\epsffile{Book03/fig28e.eps}}

For let the centers of the circles, $K$ and $L$, be found [Prop.~3.1],
and let $AK$, $KB$, $DL$, and $LE$ be joined.

And since ($ABC$ and $DEF$) are equal circles, their radii are also equal [Def.~3.1]. So the two
(straight-lines) $AK$, $KB$ are equal to the two (straight-lines) $DL$, $LE$ (respectively).
And the base $AB$ (is) equal to the base $DE$. Thus, angle $AKB$ is equal
to angle $DLE$ [Prop.~1.8]. And equal angles stand upon equal circumferences,
when they are at the centers [Prop.~3.26]. Thus, circumference $AGB$
(is) equal to $DHE$. And the whole circle $ABC$ is also equal to the whole
circle $DEF$. Thus, the remaining circumference $ACB$
is also equal to the remaining circumference $DFE$.

Thus, in equal circles, equal straight-lines cut off equal circumferences, 
the greater (circumference being equal) to the greater, and the lesser to the lesser. (Which is) the very thing it was required to show.}
\end{Parallel}

%%%%%%
% Prop 3.29
%%%%%%
\pdfbookmark[1]{Proposition 3.29}{pdf3.29}
\begin{Parallel}{}{} 
\ParallelLText{
\begin{center}
{\large \ggn{29}.}
\end{center}\vspace*{-7pt}

{\greekfont Ἐν τοῖς ἴσοις κύκλοις τὰς ἴσας περιφερείας ἴσαι εὐθεῖαι ὑποτείνουσιν.}

{\greekfont Ἔστωσαν ἴσοι κύκλοι οἱ ΑΒΓ, ΔΕΖ, καὶ ἐν α ὐτοῖς ἴσαι περιφέρειαι
ἀπειλήφθωσαν αἱ ΒΗΓ, ΕΘΖ, καὶ ἐπεζεύθθωσαν αἱ ΒΓ, ΕΖ εὐθεῖαι;
λέγω, ὅτι ἴση ἐστὶν ἡ ΒΓ τῇ ΕΖ.}

{\greekfont Εἰλήφθω γὰρ τὰ κέντρα τῶν κύκλων, καὶ ἔστω τὰ Κ, Λ, καὶ ἐπεζεύθθωσαν αἱ ΒΚ, ΚΓ, ΕΛ, ΛΖ.}

{\greekfont Καὶ ἐπεὶ ἴση ἐστὶν ἡ ΒΗΓ περιφέρεια τῇ ΕΘΖ περιφερείᾳ, ἴση
ἐστὶ καὶ γωνία ἡ ὑπὸ ΒΚΓ τῇ ὑπὸ ΕΛΖ. καὶ ἐπεὶ ἴσοι
εἰσὶν οἱ ΑΒΓ, ΔΕΖ κύκλοι, ἴσαι εἰσὶ καὶ αἱ ἐκ τῶν κέντρων;
δύο δὴ αἱ ΒΚ, ΚΓ δυσὶ ταῖς ΕΛ, ΛΖ ἴσαι εἰσίν; καὶ γωνίας
ἴσας περιέθουσιν; βάσις ἄρα ἡ ΒΓ βάσει τῇ ΕΖ ἴση ἐστίν;}\\~\\~\\~\\

\epsfysize=1.7in
\centerline{\epsffile{Book03/fig29g.eps}}

{\greekfont Ἐν ἄρα τοῖς ἴσοις κύκλοις τὰς ἴσας περιφερείας
ἴσαι εὐθεῖαι ὑποτείνουσιν; ὅπερ ἔδει δεῖχαι.}}

\ParallelRText{
\begin{center}
{\large Proposition 29}
\end{center}

In equal circles, equal straight-lines subtend equal circumferences.

Let $ABC$ and $DEF$ be equal circles, and within them let the equal circumferences $BGC$ and $EHF$ be cut off. And let the straight-lines $BC$ and $EF$ be joined. I say that $BC$ is equal to $EF$.

For let the centers of the circles be found [Prop.~3.1], and
let them be (at) $K$ and $L$. And let $BK$, $KC$, $EL$, and $LF$ be joined.

And since the circumference $BGC$ is equal to the circumference $EHF$,
the angle $BKC$ is also equal to (angle) $ELF$ [Prop.~3.27]. And since the circles
$ABC$ and $DEF$ are equal, their radii are also equal [Def.~3.1]. 
So the two (straight-lines) $BK$, $KC$ are equal to the two (straight-lines)
$EL$, $LF$ (respectively). And they contain equal angles. Thus, the
base $BC$ is equal to the base $EF$ [Prop.~1.4].

\epsfysize=1.7in
\centerline{\epsffile{Book03/fig29e.eps}}

Thus,  in equal circles, equal straight-lines subtend equal circumferences.
(Which is) the very thing it was required to show.}
\end{Parallel}

%%%%%%
% Prop 3.30
%%%%%%
\pdfbookmark[1]{Proposition 3.30}{pdf3.30}
\begin{Parallel}{}{} 
\ParallelLText{
\begin{center}
{\large \ggn{30}.}
\end{center}\vspace*{-7pt}

{\greekfont Τὴν δοθεῖσαν περιφέρειαν δίθα τεμεῖν.}

{\greekfont Ἔστω ἡ δοθεῖσα περιφέρεια ἡ ΑΔΒ; δεῖ δὴ τὴν ΑΔΒ περιφέρειαν
δίθα τεμεῖν.}

{\greekfont Ἐπεζεύθθω ἡ ΑΒ, καὶ τετμήσθω δίθα κατὰ τὸ Γ, καὶ ἀπὸ τοῦ
Γ σημείου τῇ ΑΒ εὐθείᾳ πρὸς ὀρθὰς ἤθθω ἡ ΓΔ, καὶ ἐπεζεύθθωσαν αἱ ΑΔ, ΔΒ.}

\epsfysize=1.3in
\centerline{\epsffile{Book03/fig30g.eps}}

{\greekfont Καὶ ἐπεὶ ἴση ἐστὶν ἡ ΑΓ τῇ ΓΒ, κοινὴ δὲ ἡ ΓΔ, δύο δὴ
αἱ ΑΓ, ΓΔ δυσὶ ταῖς ΒΓ, ΓΔ ἴσαι εἰσίν; καὶ γωνία ἡ ὑπὸ
ΑΓΔ γωνίᾳ τῇ ὑπὸ ΒΓΔ ἴση; ὀρθὴ γὰρ ἑκατέρα; βάσις
ἄρα ἡ ΑΔ βάσει τῇ ΔΒ ἴση ἐστίν. αἱ δὲ ἴσαι εὐθεῖαι ἴσας περιφερείας ἀφαιροῦσι τὴν μὲν μείζονα τῇ μείζονι τὴν δὲ
ἐλάττονα τῇ ἐλάττονι; κάι ἐστιν ἑκατέρα τῶν ΑΔ, ΔΒ περιφερειῶν
ἐλάττων ἡμικυκλίου; ἴση ἄρα ἡ ΑΔ περιφέρεια τῇ ΔΒ περιφερείᾳ.}

{\greekfont Ἡ ἄρα δοθεῖσα περιφέρεια δίθα τέτμ ηται κατὰ τὸ Δ σημεῖον; ὅπερ
ἔδει ποιῆσαι.}}

\ParallelRText{
\begin{center}
{\large Proposition 30}
\end{center}

To cut a given circumference in half.

Let $ADB$ be the given circumference. So it is required to cut circumference
$ADB$ in half.

Let $AB$ be joined, and let it be cut in half at (point) $C$ [Prop.~1.10]. And let $CD$ be drawn from point $C$, at right-angles
to $AB$ [Prop.~1.11]. And let $AD$, and $DB$ be joined.

\epsfysize=1.3in
\centerline{\epsffile{Book03/fig30e.eps}}

And since $AC$ is equal to $CB$, and $CD$ (is) common, the two
(straight-lines) $AC$, $CD$ are equal to the two (straight-lines) $BC$, $CD$ (respectively).
And angle $ACD$ (is) equal to angle $BCD$. For (they are) each right-angles.
Thus, the base $AD$ is equal to the base $DB$ [Prop.~1.4]. And equal
straight-lines cut off equal circumferences, the greater (circumference being
equal) to the greater, and the lesser to the lesser [Prop.~1.28].
And the circumferences $AD$ and $DB$ are each less than a semi-circle.
Thus, circumference $AD$ (is) equal to circumference $DB$.

Thus, the given circumference has been cut in half at point $D$. (Which is)
the very thing it was required to do. }
\end{Parallel}

%%%%%%
% Prop 3.31
%%%%%%
\pdfbookmark[1]{Proposition 3.31}{pdf3.31}
\begin{Parallel}{}{} 
\ParallelLText{
\begin{center}
{\large \ggn{31}.}
\end{center}\vspace*{-7pt}

{\greekfont Ἐν κύκλῳ ἡ μὲν ἐν τῷ ἡμικυκλίῳ γωνία ὀρθή ἐστιν, ἡ δὲ
ἐν τῷ μείζονι τμήματι ἐλάττων ὀρθῆς, ἡ δὲ ἐν τῷ ἐλάττονι
τμήματι μείζων ὀρθῆς; καὶ ἔπι ἡ μὲν τοῦ μείζονος τμήματος γωνία μείζων ἐστὶν ὀρθῆς, ἡ δὲ τοῦ ἐλάττονος τμήματος γωνία
ἐλάττων ὀρθῆς.}\\

\epsfysize=2.2in
\centerline{\epsffile{Book03/fig31g.eps}}

{\greekfont Ἔστω κύκλος ὁ ΑΒΓΔ, διάμετρος δὲ α ὐτοῦ ἔστω ἡ ΒΓ,
κέντρον δὲ τὸ Ε, καὶ ἐπεζεύθθωσαν αἱ ΒΑ, ΑΓ, ΑΔ, ΔΓ; λέγω,
ὅτι ἡ μὲν ἐν τῷ ΒΑΓ ἡμικυκλίῳ γωνία ἡ ὑπὸ ΒΑΓ ὀρθή
ἐστιν, ἡ δὲ ἐν τῷ ΑΒΓ μείζονι τοῦ ἡμικυκλίου τμήματι
γωνία ἡ ὑπὸ ΑΒΓ ἐλάττων ἐστὶν ὀρθῆς, ἡ δὲ ἐν τῷ ΑΔΓ
ἐλάττονι τοῦ ἡμικυκλίου τμήματι γωνία ἡ ὑπὸ ΑΔΓ μείζων
ἐστὶν ὀρθῆς.}

{\greekfont Ἐπεζεύθθω ἡ ΑΕ, καὶ διήθθω ἡ ΒΑ ἐπὶ τὸ Ζ.}

{\greekfont Καὶ ἐπεὶ ἴση ἐστὶν ἡ ΒΕ τῇ ΕΑ, ἴση ἐστὶ καὶ γωνία ἡ ὑπὸ
ΑΒΕ τῇ ὑπὸ ΒΑΕ. πάλιν, ἐπεὶ ἴση ἐστὶν ἡ ΓΕ τῇ ΕΑ, ἴση
ἐστὶ καὶ ἡ ὑπὸ ΑΓΕ τῇ ὑπὸ ΓΑΕ; ὅλη ἄρα ἡ ὑπὸ ΒΑΓ
δυσὶ ταῖς ὑπὸ ΑΒΓ, ΑΓΒ ἴση ἐστίν. ἐστὶ δὲ καὶ ἡ ὑπὸ 
ΖΑΓ ἐκτὸς τοῦ ΑΒΓ τριγώνου δυσὶ ταῖς ὑπὸ ΑΒΓ, ΑΓΒ γωνίαις
ἴση; ἴση ἄρα καὶ ἡ ὑπὸ ΒΑΓ γωνία τῇ ὑπὸ ΖΑΓ;
ὀρθὴ ἄρα ἑκατέρα; ἡ ἄρα ἐν τῷ ΒΑΓ ἡμικυκλίῳ 
γωνία ἡ ὑπὸ ΒΑΓ ὀρθή ἐστιν.}

{\greekfont Καὶ ἐπεὶ τοῦ ΑΒΓ τρίγωνου δύο γωνίαι αἱ ὑπὸ ΑΒΓ, ΒΑΓ δύο
ὀρθῶν ἐλάττονές εἰσιν, ὀρθὴ δὲ ἡ ὑπὸ ΒΑΓ, ἐλάττων ἄρα
ὀρθῆς ἐστιν ἡ ὑπὸ ΑΒΓ γωνία; καί ἐστιν ἐν τῷ ΑΒΓ μείζονι
τοῦ ἡμικυκλίου τμήματι.}

{\greekfont Καὶ ἐπεὶ ἐν κύκλῳ τετράπλευρόν ἐστι τὸ ΑΒΓΔ, τῶν δὲ ἐν τοῖς
κύκλοις τετραπλεύρων αἱ ἀπεναντίον γωνίαι δυσὶν ὀρθαῖς ἴσαι
εἰσίν [αἱ ἄρα ὑπὸ ΑΒΓ, ΑΔΓ γωνίαι δυσὶν ὀρθαῖς ἴσας
εἰσίν], καί ἐστιν ἡ ὑπὸ ΑΒΓ ἐλάττων ὀρθῆς; λοιπὴ ἄρα ἡ ὑπὸ
ΑΔΓ γωνία μείζων ὀρθῆς ἐστιν; καί ἐστιν ἐν τῷ ΑΔΓ
ἐλάττονι τοῦ ἡμικυκλίου τμήματι.}

{\greekfont Λέγω, ὅτι καὶ ἡ μὲν τοῦ μείζονος τμήματος γωνία ἡ περιεθομένη
ὑπό [τε] τῆς ΑΒΓ περιφερείας καὶ τῆς ΑΓ εὐθείας μείζων ἐστὶν
ὀρθῆς, ἡ δὲ τοῦ ἐλάττονος τμήματος γωνία ἡ περιεθομένη
ὑπό [τε] τῆς ΑΔ[Γ] περιφερείας καὶ τῆς ΑΓ εὐθείας ἐλάττων ἐστὶν
ὀρθῆς. καί ἐστιν α ὐτόθεν φανερόν. ἐπεὶ γὰρ ἡ ὑπὸ τῶν ΒΑ,
ΑΓ εὐθειῶν ὀρθή ἐστιν, ἡ ἄρα ὑπὸ τῆς ΑΒΓ περιφερείας καὶ τῆς
ΑΓ εὐθείας περιεθομένη μείζων ἐστὶν ὀρθῆς. πάλιν, ἐπεὶ
ἡ ὑπὸ τῶν ΑΓ, ΑΖ εὐθειῶν ὀρθή ἐστιν, ἡ ἄρα ὑπὸ
τῆς ΓΑ εὐθείας καὶ τῆς ΑΔ[Γ] περιφερείας περιεθομένη ἐλάττων
ἐστὶν ὀρθῆς.}

{\greekfont Ἐν κύκλῳ ἄρα ἡ μὲν ἐν τῷ ἡμικυκλίῳ γωνία ὀρθή ἐστιν, ἡ δὲ
ἐν τῷ μείζονι τμήματι ἐλάττων ὀρθῆς, ἡ δὲ ἐν τῷ ἐλάττονι
[τμήματι] μείζων ὀρθῆς; καὶ ἔπι ἡ μὲν τοῦ μείζονος τμήματος [γωνία] μείζων [ἐστὶν] ὀρθῆς, ἡ δὲ τοῦ ἐλάττονος τμήματος [γωνία] ἐλάττων ὀρθῆς; ὅπερ ἔδει δεῖχαι.}}

\ParallelRText{
\begin{center}
{\large Proposition 31}
\end{center}

In a circle, the angle in a semi-circle is a right-angle, and that in a
greater segment (is) less than a right-angle, and that in a
lesser segment (is) greater than a right-angle. And, further, the angle of
a segment greater (than a semi-circle) is greater than a right-angle, and
the angle of a segment less (than a semi-circle) is less than a right-angle.

\epsfysize=2.2in
\centerline{\epsffile{Book03/fig31e.eps}}

Let $ABCD$  be a circle, and let $BC$ be its diameter, and $E$ its center. And let $BA$, $AC$, $AD$, and $DC$ be joined. I say that the angle $BAC$ in the semi-circle $BAC$ is a right-angle, and the angle $ABC$ in the segment $ABC$, (which is) greater than a semi-circle, is less
than a right-angle, and the angle $ADC$ in the segment $ADC$, (which is) less than a semi-circle,
is greater than a right-angle.

Let $AE$ be joined, and let $BA$ be drawn through to $F$.

And since $BE$ is equal to $EA$, angle $ABE$ is also equal to $BAE$ [Prop.~1.5]. Again, since $CE$ is equal to $EA$, $ACE$ is also equal to  $CAE$ [Prop.~1.5].
Thus, the whole (angle) $BAC$ is equal to the two (angles) $ABC$ and $ACB$.
And $FAC$, (which is) external to triangle $ABC$, is also equal to the two angles
$ABC$ and $ACB$ [Prop.~1.32]. Thus, angle $BAC$ (is) also equal to
$FAC$. Thus, (they are) each right-angles. [Def.~1.10]. Thus, the angle
$BAC$ in the semi-circle $BAC$ is a right-angle.

And since the two angles $ABC$ and $BAC$ of triangle $ABC$ are less than 
two right-angles [Prop.~1.17], and $BAC$ is a right-angle, angle $ABC$ is
thus less than a right-angle. And it is in segment $ABC$, (which is) greater 
than
a semi-circle.

And since $ABCD$ is a quadrilateral within a circle, and for quadrilaterals within circles
the (sum of the) opposite angles is equal to two right-angles 
[Prop.~3.22] [angles $ABC$ and $ADC$ are thus equal to two right-angles],
and (angle) $ABC$ is  less than a right-angle. The remaining angle $ADC$
is thus greater than a right-angle. And it is in segment $ADC$, (which is)
less than a semi-circle.

I also say that the angle of the greater segment, (namely) that contained by
the circumference $ABC$ and the straight-line $AC$, is greater than a right-angle.
And the angle of the lesser segment, (namely) that contained by the circumference
$AD[C]$ and the straight-line $AC$, is less than a right-angle. And
this is immediately apparent. For since the (angle contained by) the
two straight-lines $BA$ and $AC$ is a right-angle, the (angle) contained
by the circumference $ABC$ and the straight-line $AC$ is thus greater than a right-angle. Again, since the (angle contained by) the straight-lines
$AC$ and $AF$ is a right-angle, the (angle) contained by the circumference
$AD[C]$ and the straight-line $CA$ is thus less than a right-angle.

Thus, in a circle, the angle in a semi-circle is a right-angle, and that in a
greater segment (is) less than a right-angle, and that in a
lesser [segment] (is) greater than a right-angle. And, further, the [angle] of
a segment greater (than a semi-circle) [is] greater than a right-angle, and
the [angle] of a segment less (than a semi-circle) is  less than a right-angle.
(Which is) the very thing it was required to show.}
\end{Parallel}

%%%%%%
% Prop 3.32
%%%%%%
\pdfbookmark[1]{Proposition 3.32}{pdf3.32}
\begin{Parallel}{}{} 
\ParallelLText{
\begin{center}
{\large \ggn{32}.}
\end{center}\vspace*{-7pt}

{\greekfont Ἐὰν κύκλου ἐφάπτηταί τις εὐθεῖα, ἀπὸ δὲ τῆς ἁφῆς εἰς τὸν κύκλον διαθθῇ τις εὐθεῖα τέμνουσα τὸν κύκλον, ἃς ποιεῖ γωνίας
πρὸς τῇ ἐφαπτομένῃ, ἴσαι ἔσονται ταῖς ἐν τοῖς ἐναλλὰχ
τοῦ κύκλου τμήμασι γωνίαις.}\\

\epsfysize=2.1in
\centerline{\epsffile{Book03/fig32e.eps}}

{\greekfont Κύκλου γὰρ τοῦ ΑΒΓΔ ἐφαπτέσθω τις εὐθεῖα ἡ ΕΖ κατὰ τὸ Β σημεῖον, καὶ ἀπὸ τοῦ Β σημείου διήθθω τις εὐθεῖα εἰς τὸν
ΑΒΓΔ κύκλον τέμνουσα α ὐτὸν ἡ ΒΔ. λέγω, ὅτι ἃς ποιεῖ γωνίας
ἡ ΒΔ μετὰ τῆς ΕΖ  ἐφαπτομένης, ἴσας ἔσονται ταῖς ἐν τοῖς
ἐναλλὰχ τμήμασι τοῦ κύκλου γωνίαις, τουτέστιν, ὅτι ἡ μὲν
ὑπὸ ΖΒΔ γωνία ἴση ἐστὶ τῇ ἐν τῷ ΒΑΔ τμήματι συνισταμένῃ γωνίᾳ, ἡ δὲ ὑπὸ ΕΒΔ γωνία ἴση ἐστὶ τῇ ἐν τῷ ΔΓΒ τμήματι συνισταμένῃ γωνίᾳ.}

{\greekfont Ἤθθω γὰρ ἀπὸ τοῦ Β τῇ ΕΖ πρὸς ὀρθὰς ἡ ΒΑ, καὶ εἰλήφθω ἐπὶ
τῆς ΒΔ περιφερείας τυθὸν σημεῖον τὸ Γ, καὶ ἐπεζεύθθωσαν αἱ ΑΔ, ΔΓ, ΓΒ.}

{\greekfont Καὶ ἐπεὶ κύκλου τοῦ ΑΒΓΔ ἐφάπτεταί τις εὐθεῖα ἡ ΕΖ κατὰ
τὸ Β, καὶ ἀπὸ τῆς ἁφῆς ἦκται τῇ ἐφαπτομένῃ πρὸς ὀρθὰς ἡ ΒΑ, ἐπὶ τῆς ΒΑ ἄρα τὸ κέντρον ἐστὶ τοῦ ΑΒΓΔ κύκλου. ἡ ΒΑ
ἄρα διάμετός ἐστι τοῦ ΑΒΓΔ κύκλου; ἡ ἄρα ὑπὸ ΑΔΒ γωνία
ἐν ἡμικυκλίῳ οὖσα ὀρθή ἐστιν. λοιπαὶ ἄρα αἱ ὑπὸ ΒΑΔ, ΑΒΔ
μιᾷ ὀρθῇ ἴσαι εἰσίν. ἐστὶ δὲ καὶ ἡ ὑπὸ ΑΒΖ ὀρθή; ἡ ἄρα ὑπὸ ΑΒΖ ἴση ἐστὶ ταῖς ὑπὸ ΒΑΔ, ΑΒΔ. κοινὴ ἀφῃρήσθω ἡ ὑπὸ
ΑΒΔ; λοιπὴ ἄρα ἡ ὑπὸ ΔΒΖ γωνία ἴση ἐστὶ τῇ ἐν τῷ ἐναλλὰχ τμήματι τοῦ κύκλου γωνίᾳ τῇ ὑπὸ ΒΑΔ. καὶ ἐπεὶ ἐν κύκλῳ
τετράπλευρόν ἐστι τὸ ΑΒΓΔ, αἱ ἀπεναντίον α ὐτοῦ γωνίαι δυσὶν
ὀρθαῖς ἴσαι  εἰσίν. εἰσὶ δὲ καὶ αἱ ὑπὸ ΔΒΖ, ΔΒΕ δυσὶν ὀρθαῖς
ἴσαι;
 αἱ ἄρα ὑπὸ ΔΒΖ, ΔΒΕ ταῖς ὑπὸ ΒΑΔ, ΒΓΔ ἴσαι
εἰσίν, ῶ̔ν ἡ ὑπὸ ΒΑΔ τῇ ὑπὸ ΔΒΖ ἐδείθθη ἴση;
λοιπὴ ἄρα ἡ ὑπὸ ΔΒΕ τῇ ἐν τῷ ἐναλλὰχ τοῦ κύκλου τμήματι
τῷ ΔΓΒ τῇ ὑπὸ ΔΓΒ γωνίᾳ ἐστὶν ἴση.}

{\greekfont Ἐὰν ἄρα κύκλου ἐφάπτηταί τις εὐθεῖα, ἀπὸ δὲ τῆς ἁφῆς εἰς τὸν κύκλον διαθθῇ τις εὐθεῖα τέμνουσα τὸν κύκλον, ἃς ποιεῖ γωνίας
πρὸς τῇ ἐφαπτομένῃ, ἴσαι ἔσονται ταῖς ἐν τοῖς ἐναλλὰχ
τοῦ κύκλου τμήμασι γωνίαις; ὅπερ ἔδει δεῖχαι.}}

\ParallelRText{
\begin{center}
{\large Proposition 32}
\end{center}

If some straight-line touches a circle, and some (other) straight-line is drawn  
across, from the point of contact into the circle,  cutting the circle (in two), (then) those angles  the (straight-line) makes
with the tangent will be equal to the angles in
the alternate segments of the circle.

\epsfysize=2.1in
\centerline{\epsffile{Book03/fig32e.eps}}

For let some straight-line $EF$ touch the circle $ABCD$ at the point $B$, and let
some (other) straight-line $BD$ be drawn from point $B$ into the
circle $ABCD$, cutting it (in two). I say that the angles $BD$ makes with the tangent 
$EF$  will be equal to the angles in the alternate segments of the circle.
That is to say, that angle $FBD$ is equal to the angle constructed
in segment $BAD$, and angle $EBD$ is equal to the angle constructed in segment $DCB$.

For let $BA$ be drawn from $B$, at right-angles to $EF$ [Prop.~1.11].
And let the point $C$ be taken at random on the circumference $BD$. And
let $AD$, $DC$, and $CB$ be joined.

And since some straight-line $EF$ touches the circle $ABCD$ at point $B$, and
$BA$ has been drawn from the point of contact, at right-angles to the
tangent, the center of circle $ABCD$ is thus on $BA$ [Prop.~3.19].
Thus, $BA$ is a diameter of circle $ABCD$.
Thus, angle $ADB$, being in a semi-circle, is a right-angle [Prop.~3.31].
Thus, the remaining angles (of triangle $ADB$) $BAD$ and $ABD$ are equal to one right-angle [Prop.~1.32]. And $ABF$ is also a right-angle. Thus,
$ABF$ is equal to $BAD$ and $ABD$. Let $ABD$ be subtracted from both.
Thus, the remaining angle $DBF$ is equal to the angle $BAD$ in the alternate
segment of the circle. And since $ABCD$ is a quadrilateral in a circle,
(the sum of) its opposite angles is equal to two right-angles [Prop.~3.22]. And $DBF$ and $DBE$ is also equal to two right-angles [Prop.~1.13]. Thus, $DBF$ and $DBE$ is equal to $BAD$ and $BCD$, of which
$BAD$ was shown (to be) equal to $DBF$. Thus, the remaining (angle) $DBE$ is
equal to the angle $DCB$ in the alternate segment $DCB$ of the circle.

Thus, if some straight-line touches a circle, and some (other) straight-line is drawn  across, from the point of contact into the circle,  cutting the circle (in two), (then) those angles  the (straight-line) makes
with the tangent will be equal to the angles in
the alternate segments of the circle. (Which is) the very thing it was required to
show.}
\end{Parallel}

%%%%%%
% Prop 3.33
%%%%%%
\pdfbookmark[1]{Proposition 3.33}{pdf3.33}
\begin{Parallel}{}{} 
\ParallelLText{
\begin{center}
{\large \ggn{33}.}
\end{center}\vspace*{-7pt}

{\greekfont Ἐπὶ τῆς δοθείσης εὐθείας γράψαι τμῆμα κύκλου δεθόμε-νον γωνίαν ἴσην τῇ δοθείσῃ γωνίᾳ εὐθυγράμμῳ.}

\epsfysize=1.35in
\centerline{\epsffile{Book03/fig33g.eps}}

{\greekfont Ἔστω ἡ δοθεῖσα εὐθεῖα ἡ ΑΒ, ἡ δὲ δοθεῖσα γωνία εὐθύγραμμος ἡ
πρὸς τῷ Γ; δεῖ δὴ ἐπὶ τῆς δοθείσης εὐθείας τῆς ΑΒ γράψαι τμῆμα
κύκλου δεθόμενον γωνίαν ἴσην τῇ πρὸς τῷ Γ.}

{\greekfont Ἡ δὴ πρὸς τῷ Γ [γωνία] ἤτοι ὀχεῖά ἐστιν ἢ ὀρθὴ ἢ ἀμβλεῖα;
ἔστω πρότερον ὀχεῖα, καὶ ὡς ἐπὶ τῆς πρώτης καταγραφῆς συνεστάτω
πρὸς τῇ ΑΒ εὐθείᾳ καὶ τῷ Α σημείῳ τῇ πρὸς τῷ Γ γωνίᾳ
ἴση ἡ ὑπὸ ΒΑΔ; ὀχεῖα ἄρα ἐστὶ καὶ ἡ ὑπὸ ΒΑΔ. ἤθθω
τῇ ΔΑ πρὸς ὀρθὰς ἡ ΑΕ, καὶ τετμήσθω ἡ ΑΒ δίθα κατὰ τὸ Ζ, καὶ
ἤθθω ἀπὸ τοῦ Ζ σημείου τῇ ΑΒ πρὸς ὀρθὰς ἡ ΖΗ, καὶ
ἐπεζεύθθω ἡ ΗΒ.}

{\greekfont Καὶ ἐπεὶ ἴση ἐστὶν ἡ ΑΖ τῇ ΖΒ, κοινὴ δὲ ἡ ΖΗ, δύο δὴ αἱ
ΑΖ, ΖΗ δύο ταῖς ΒΖ, ΖΗ ἴσαι εἰσίν; καὶ γωνία ἡ ὑπὸ ΑΖΗ [γωνίᾳ]
τῇ ὑπὸ ΒΖΗ ἴση; βάσις ἄρα ἡ ΑΗ βάσει τῇ ΒΗ ἴση ἐστίν. ὁ ἄρα κέντρῳ μὲν τῷ Η διαστήματι δὲ τῷ ΗΑ κύκλος γραφόμενος ἥχει καὶ διὰ τοῦ Β. γεγράφθω καὶ ἔστω ὁ ΑΒΕ, καὶ ἐπεζεύθθω ἡ ΕΒ.
ἐπεὶ οὖν ἀπ ἄκρας τῆς ΑΕ διαμέτρου ἀπὸ τοῦ Α τῇ ΑΕ πρὸς ὀρθάς
ἐστιν ἡ ΑΔ, ἡ ΑΔ ἄρα ἐφάπτεται τοῦ ΑΒΕ κύκλου; ἐπεὶ οὖν
κύκλου τοῦ ΑΒΕ ἐφάπτεταί τις εὐθεῖα ἡ ΑΔ, καὶ ἀπὸ τῆς κατὰ
τὸ Α ἁφῆς  εἰς τὸν ΑΒΕ κύκλον διῆκταί τις εὐθεῖα ἡ ΑΒ, ἡ
ἄρα ὑπὸ ΔΑΒ γωνία ἴση ἐστὶ τῇ ἐν τῷ ἐναλλὰχ τοῦ 
κύκλου τμήματι γωνίᾳ τῇ ὑπὸ ΑΕΒ. ἀλλ ἡ ὑπὸ ΔΑΒ τῇ
πρὸς τῷ Γ ἐστιν ἴση; καὶ ἡ πρὸς τῷ Γ ἄρα γωνία ἴση
ἐστὶ τῇ ὑπὸ ΑΕΒ.}

{\greekfont Ἐπὶ τῆς δοθείσης ἄρα εὐθείας τῆς ΑΒ τμῆμα κύκλου γέγραπται
τὸ ΑΕΒ δεθόμενον γωνίαν τὴν ὑπὸ ΑΕΒ ἴσην τῇ δοθείσῃ
τῇ πρὸς τῷ Γ.}

{\greekfont Ἀλλὰ δὴ ὀρθὴ ἔστω ἡ πρὸς τῷ Γ; καὶ δέον πάλιν ἔστω ἐπὶ
τῆς ΑΒ γράψαι τμῆμα κύκλου δεθόμενον γωνίαν ἴσην τῇ πρὸς
τῷ Γ ὀρθῇ [γωνίᾳ]. συνεστάτω [πάλιν] τῇ πρὸς τῷ Γ ὀρθῇ γωνίᾳ ἴση ἡ ὑπὸ ΒΑΔ, ὡς ἔχει ἐπὶ τῆς δευτέρας καταγραφῆς, καὶ τετμήσθω ἡ ΑΒ δίθα κατὰ τὸ Ζ, καὶ κέντρῳ τῷ Ζ, διαστήματι δὲ
ὁποτέρῳ τῶν ΖΑ, ΖΒ, κύκλος γεγράφθω ὁ ΑΕΒ.}

{\greekfont Ἐφάπτεται ἄρα ἡ ΑΔ εὐθεῖα τοῦ ΑΒΕ κύκλου διὰ τὸ ὀρθὴν εἶναι
τὴν πρὸς τῷ Α γωνίαν. καὶ ἴση ἐστὶν ἡ ὑπὸ ΒΑΔ γωνία
τῇ ἐν τῷ ΑΕΒ τμήματι; ὀρθὴ γὰρ καὶ α ὐτὴ ἐν ἡμικυκλίῳ
οὖσα. ἀλλὰ καὶ ἡ ὑπὸ ΒΑΔ τῇ πρὸς τῷ Γ ἴση ἐστίν. καὶ
ἡ ἐν τῷ ΑΕΒ ἄρα ἴση ἐστὶ  τῇ πρὸς τῷ Γ.}

{\greekfont Γέγραπται ἄρα πάλιν ἐπὶ τῆς ΑΒ τμῆμα κύκλου τὸ ΑΕΒ δεθόμενον
γωνίαν ἴσην τῇ πρὸς τῷ Γ.}

{\greekfont Ἀλλὰ δὴ ἡ πρὸς τῷ Γ ἀμβλεῖα ἔστω; καὶ συνεστάτω α ὐτῇ
ἴση πρὸς τῇ ΑΒ εὐθείᾳ καὶ τῷ Α σημείῳ ἡ ὑπὸ ΒΑΔ, ὡς
ἔχει ἐπὶ τῆς τρίτης καταγραφῆς, καὶ τῇ ΑΔ πρὸς ὀρθὰς ἤθθω ἡ ΑΕ,
καὶ τετμήσθω πάλιν ἡ ΑΒ δίθα κατὰ τὸ Ζ, καὶ τῇ ΑΒ πρὸς ὀρθὰς
ἤθθω ἡ ΖΗ, καὶ ἐπεζεύθθω ἡ ΗΒ.}

{\greekfont Καὶ ἐπεὶ πάλιν ἴση ἐστὶν ἡ ΑΖ τῇ ΖΒ, καὶ κοινὴ ἡ ΖΗ, δύο δὴ αἱ
ΑΖ, ΖΗ δύο ταῖς ΒΖ, ΖΗ ἴσαι εἰσίν; καὶ γωνία ἡ ὑπὸ
ΑΖΗ γωνίᾳ τῇ ὑπὸ ΒΖΗ ἴση; βάσις ἄρα ἡ ΑΗ βάσει τῇ ΒΗ
ἴση ἐστίν; ὁ ἄρα κέντρῳ μὲν τῷ Η διαστήματι δὲ τῷ ΗΑ κύκλος
γραφόμενος ἥχει καὶ διὰ τοῦ Β. ἐρθέσθω ὡς ὁ ΑΕΒ. καὶ ἐπεὶ
τῇ ΑΕ διαμέτρῳ ἀπ ἄκρας πρὸς ὀρθάς ἐστιν ἡ ΑΔ, ἡ ΑΔ ἄρα
ἐφάπτεται τοῦ ΑΕΒ κύκλου. καὶ ἀπὸ τῆς κατὰ τὸ Α ἐπαφῆς διῆκται
ἡ ΑΒ; ἡ ἄρα ὑπὸ ΒΑΔ γωνία ἴση ἐστὶ τῇ ἐν τῷ ἐναλλὰχ τοῦ
κύκλου τμήματι  τῷ ΑΘΒ συνισταμένῃ γωνίᾳ. ἀλλ ἡ ὑπὸ ΒΑΔ γωνία τῇ πρὸς τῷ Γ ἴση ἐστίν. καὶ ἡ ἐν τῷ ΑΘΒ ἄρα τμήματι
γωνία ἴση  ἐστὶ τῇ πρὸς τῷ Γ.}

{\greekfont Ἐπὶ τῆς ἄρα δοθείσης εὐθείας τῆς ΑΒ γέγραπται τμῆμα κύκλου τὸ
ΑΘΒ δεθόμενον γωνίαν ἴσην τῇ πρὸς τῷ Γ; ὅπερ ἔδει ποιῆσαι.}}

\ParallelRText{
\begin{center}
{\large Proposition 33}
\end{center}

To draw a segment of a circle, accepting an angle equal to a given rectilinear angle, on a given straight-line.

\epsfysize=1.53in
\centerline{\epsffile{Book03/fig33e.eps}}

Let $AB$ be the given straight-line, and $C$ the given rectilinear angle. So
it is required to draw a segment of a circle, accepting an angle equal to $C$,
on the given straight-line $AB$.

So the [angle] $C$ is surely either acute, a right-angle, or obtuse. First of all,
let it be acute. And, as in the first diagram (from the left), let (angle) $BAD$,
equal to angle $C$,  be
constructed  on the straight-line $AB$,  at the point A (on it)
[Prop.~1.23]. Thus, $BAD$ is also acute. Let $AE$ be drawn, at
right-angles to $DA$ [Prop.~1.11]. And let $AB$ be cut in half
at $F$ [Prop.~1.10]. And let $FG$ be drawn from point $F$, at right-angles to $AB$
[Prop.~1.11]. And let $GB$ be joined.

And since $AF$ is equal to $FB$, and $FG$ (is) common, the two (straight-lines)
$AF$, $FG$ are equal to the two (straight-lines) $BF$, $FG$ (respectively). And
angle $AFG$ (is) equal to [angle] $BFG$. Thus, the base $AG$ is equal to the
base $BG$ [Prop.~1.4]. Thus, the circle drawn with center $G$, and radius $GA$, will also go through $B$ (as well as $A$). Let it be drawn, and let it be (denoted) $ABE$. And let $EB$ be joined. Therefore, since $AD$ is at the extremity
of diameter $AE$,  (namely, point) $A$,  at right-angles to $AE$, the (straight-line)
$AD$ thus touches the circle $ABE$ [Prop.~3.16~corr.]. Therefore,
since some straight-line $AD$ touches the circle $ABE$, and some (other)
straight-line $AB$ has been drawn across from the point of contact $A$ into circle
$ABE$, angle $DAB$ is thus equal to the angle $AEB$ in the alternate segment of the
circle [Prop.~3.32]. But, $DAB$ is equal to $C$. Thus,
angle $C$ is also equal to $AEB$.

Thus, a segment $AEB$ of a circle, accepting the angle $AEB$ (which is) equal to the
given (angle) $C$, has been drawn on the given straight-line $AB$.

And so let $C$ be a right-angle. And let it again be necessary to draw  a segment
of a circle on $AB$, accepting an angle equal to the right-[angle] $C$. Let the (angle)
$BAD$ [again] be constructed, equal to the right-angle $C$ [Prop.~1.23], as in the second diagram (from the left). And let $AB$ be cut in half at $F$ [Prop.~1.10]. And let the circle $AEB$ be
drawn with center $F$, and radius either  $FA$ or $FB$.

Thus, the straight-line $AD$ touches the circle $ABE$, on account of the angle at $A$ being a right-angle [Prop.~3.16 corr.]. And angle $BAD$ is equal
to the angle in segment $AEB$. For (the latter angle), being in a semi-circle, is also a right-angle [Prop.~3.31]. But,  $BAD$ is also equal to $C$. Thus, the (angle) in (segment) $AEB$ is also equal to $C$.

Thus, a segment $AEB$ of a circle, accepting an angle equal to $C$, has again been
drawn on $AB$.

And so let (angle) $C$ be obtuse. And let (angle) $BAD$, equal
to ($C$), be constructed  on the straight-line $AB$, at the point $A$ 
(on it) [Prop.~1.23], as in the third diagram (from the left). And let $AE$ be
drawn, at right-angles to $AD$ [Prop.~1.11]. And let $AB$ again be 
cut in half at $F$ [Prop.~1.10]. And let $FG$ be drawn,
at right-angles to $AB$ [Prop.~1.10]. And let $GB$ 
be joined.

And again, since $AF$ is equal to $FB$, and $FG$ (is) common, the two (straight-lines) $AF$, $FG$ are equal to the two (straight-lines) $BF$, $FG$ (respectively).
And angle $AFG$ (is) equal to angle $BFG$. Thus, the base $AG$ is equal to
the base $BG$ [Prop.~1.4]. Thus, a circle of center $G$, and radius $GA$,
being drawn, will also go through $B$ (as well as $A$). Let it go like $AEB$ (in the third diagram
from the left). And since $AD$ is at right-angles to the diameter $AE$, at its extremity, $AD$ thus touches circle $AEB$ [Prop.~3.16~corr.]. And $AB$ has been drawn across (the circle) from the point of contact $A$. Thus, angle $BAD$ is equal to the angle constructed in the alternate segment $AHB$ of the circle 
[Prop.~3.32]. But, angle $BAD$ is equal to $C$. Thus, the angle in segment
$AHB$ is also equal to $C$.

Thus, a segment $AHB$ of a circle, accepting an angle equal to $C$, has been
drawn on the given straight-line $AB$. (Which is) the very thing
it was required to do.}
\end{Parallel}

%%%%%%
% Prop 3.34
%%%%%%
\pdfbookmark[1]{Proposition 3.34}{pdf3.34}
\begin{Parallel}{}{} 
\ParallelLText{
\begin{center}
{\large\ggn{34}.}
\end{center}\vspace*{-7pt}

{\greekfont Ἀπὸ τοῦ δοθέντος κύκλου τμῆμα ἀφελεῖν δεθόμενον γωνίαν
ἴσην τῇ δοθείσῃ γωνίᾳ εὐθυγράμμῳ.}

{\greekfont Ἔστω ὁ δοθεὶς κύκλος ὁ ΑΒΓ, ἡ δὲ δοθεῖσα γωνία εὐθύγραμμος
ἡ πρὸς  τῷ Δ; δεῖ δὴ ἀπὸ τοῦ ΑΒΓ κύκλου τμῆμα ἀφελεῖν
δεθόμενον γωνίαν ἴσην τῇ δοθείσῃ γωνίᾳ εὐθυγράμμῳ τῇ πρὸς
τῷ Δ.}

\epsfysize=1.8in
\centerline{\epsffile{Book03/fig34g.eps}}

{\greekfont Ἤθθω τοῦ ΑΒΓ ἐφαπτομένη ἡ ΕΖ κατὰ τὸ Β σημεῖον, καὶ συνεστάτω
πρὸς τῇ ΖΒ εὐθείᾳ καὶ τῷ πρὸς α ὐτῇ σημείῳ τῷ Β τῇ πρὸς
τῷ Δ γωνίᾳ ἴση ἡ ὑπὸ ΖΒΓ.}

{\greekfont Ἐπεὶ οὖν κύκλου τοῦ ΑΒΓ ἐφάπτεταί τις εὐθεῖα ἡ ΕΖ, καὶ ἀπὸ
τῆς κατὰ τὸ Β ἐπαφῆς διῆκται ἡ ΒΓ, ἡ ὑπὸ ΖΒΓ ἄρα γωνία
ἴση ἐστὶ τῇ ἐν τῷ ΒΑΓ ἐναλλὰχ τμήματι συνισταμένῃ γωνίᾳ.
ἀλλ ἡ ὑπὸ ΖΒΓ τῇ πρὸς τῷ Δ ἐστιν ἴση; καὶ ἡ ἐν τῷ ΒΑΓ
ἄρα τμήματι ἴση ἐστὶ τῇ πρὸς τῷ Δ [γωνίᾳ].}

{\greekfont Ἀπὸ τοῦ δοθέντος ἄρα κύκλου τοῦ ΑΒΓ τμῆμα ἀφῄρηται τὸ ΒΑΓ
δεθόμενον γωνίαν ἴσην τῇ δοθείσῃ γωνίᾳ εὐθυγράμ-μῳ τῇ
πρὸς τῷ Δ; ὅπερ ἔδει ποιῆσαι.}}

\ParallelRText{
\begin{center}
{\large Proposition 34}
\end{center}

To cut off a segment, accepting an angle equal to a given rectilinear angle,
from a given circle.

Let $ABC$ be the given circle, and $D$ the given rectilinear angle. So it is
required to cut off a segment, accepting an angle equal to the given
rectilinear angle $D$, from the given circle $ABC$.

\epsfysize=1.8in
\centerline{\epsffile{Book03/fig34e.eps}}

Let $EF$ be drawn touching $ABC$ at point $B$.$^\dag$ And let (angle)
$FBC$, equal to angle $D$, be constructed  on the
straight-line $FB$, at the point $B$ on it [Prop.~1.23].

Therefore, since some straight-line $EF$ touches the circle $ABC$, and
$BC$ has been drawn across (the circle) from the point of contact $B$, angle
$FBC$ is thus equal to the angle constructed in the alternate segment
$BAC$ [Prop.~1.32]. But, $FBC$ is equal to $D$. Thus, the (angle)
in the segment $BAC$ is also equal to [angle] $D$.

Thus, the segment $BAC$, accepting an angle equal to the given
rectilinear angle $D$, has been cut off from the given circle $ABC$. (Which is)
the very thing it was required to do.}
\end{Parallel}
{\footnotesize \noindent$^\dag$ Presumably,
by finding the center of $ABC$ [Prop.~3.1], drawing a straight-line between
the center and point $B$, and then drawing $EF$ through point $B$, at right-angles
to the aforementioned straight-line [Prop.~1.11].}

%%%%%%
% Prop 3.35
%%%%%%
\pdfbookmark[1]{Proposition 3.35}{pdf3.35}
\begin{Parallel}{}{} 
\ParallelLText{
\begin{center}
{\large \ggn{35}.}
\end{center}\vspace*{-7pt}

{\greekfont Ἐὰν ἐν κύκλῳ δύο εὐθεῖαι τέμνωσιν ἀλλήλας, τὸ ὑπὸ τῶν
τῆς μιᾶς τμ ημάτων περιεχόμενον ὀρθογώνιον ἴσον ἐστὶ
τῷ ὑπὸ τῶν τῆς ἑτέρας τμ ημάτων περιεθομένῳ ὀρθογωνίῳ.}

{\greekfont Ἐν γὰρ κύκλῳ τῷ ΑΒΓΔ δύο εὐθεῖαι αἱ ΑΓ, ΒΔ τεμνέτωσαν ἀλλήλας κατὰ τὸ Ε σημεῖον; λέγω, ὅτι τὸ ὑπὸ τῶν ΑΕ, ΕΓ
περιεχόμενον ὀρθογώνιον ἴσον ἐστὶ τῷ ὑπὸ τῶν ΔΕ, ΕΒ
περιεθομένῳ ὀρθογωνίῳ.}

{\greekfont Εἰ μὲν οὖν αἱ ΑΓ, ΒΔ διὰ τοῦ κέντρου εἰσὶν ὥστε τὸ Ε
κέντρον εἶναι τοῦ ΑΒΓΔ κύκλου, φανερόν, ὅτι
ἴσων οὐσῶν τῶν ΑΕ, ΕΓ, ΔΕ, ΕΒ καὶ τὸ ὑπὸ τῶν ΑΕ,
ΕΓ περιεχόμενον ὀρθογώνιον ἴσον ἐστὶ τῷ ὑπὸ τῶν
ΔΕ, ΕΒ περιεθομένῳ ὀρθογωνίῳ.}

{\greekfont μὴ ἔστωσαν δὴ αἱ ΑΓ, ΔΒ διὰ τοῦ κέντρου, καὶ εἰλήφθω τὸ
 κέντρον τοῦ ΑΒΓΔ, καὶ ἔστω τὸ Ζ,  καὶ ἀπὸ τοῦ Ζ ἐπὶ
τὰς ΑΓ, ΔΒ εὐθείας κάθετοι ἤθθωσαν αἱ ΖΗ, ΖΘ, καὶ ἐπεζεύθθωσαν
αἱ ΖΒ, ΖΓ, ΖΕ.}

{\greekfont Καὶ ἐπεὶ εὐθεῖά τις διὰ τοῦ κέντρου ἡ ΗΖ εὐθεῖάν τινα μὴ διὰ
τοῦ κέντρου τὴν ΑΓ πρὸς ὀρθὰς τέμνει, καὶ δίθα α ὐτὴν τέμνει;
ἴση ἄρα ἡ ΑΗ τῇ ΗΓ. ἐπεὶ οὖν εὐθεῖα ἡ ΑΓ τέτμ ηται εἰς μὲν
ἴσα κατὰ τὸ Η, εἰς δὲ ἄνισα κατὰ τὸ Ε, τὸ ἄρα ὑπὸ τῶν
ΑΕ, ΕΓ περιεχόμενον ὀρθογώνιον μετὰ τοῦ ἀπὸ τῆς ΕΗ τετραγώνου
ἴσον ἐστὶ τῷ ἀπὸ τῆς ΗΓ; [κοινὸν] προσκείσθω τὸ ἀπὸ τῆς ΗΖ;
τὸ ἄρα ὑπὸ τῶν ΑΕ, ΕΓ μετὰ τῶν ἀπὸ τῶν ΗΕ, ΗΖ ἴσον ἐστὶ
τοῖς ἀπὸ τῶν ΓΗ, ΗΖ. ἀλλὰ τοῖς μὲν ἀπὸ τῶν ΕΗ, ΗΖ ἴσον
ἐστὶ τὸ ἀπὸ τῆς ΖΕ, τοὶς δὲ ἀπὸ τῶν ΓΗ, ΗΖ ἴσον
ἐστὶ τὸ ἀπὸ τῆς ΖΓ; τὸ ἄρα ὑπὸ τῶν ΑΕ, ΕΓ  μετὰ τοῦ ἀπὸ τῆς ΖΕ ἴσον  ἐστὶ τῷ ἀπὸ τῆς ΖΓ. ἴση δὲ ἡ ΖΓ τῇ ΖΒ; τὸ
ἄρα ὑπὸ τῶν ΑΕ, ΕΓ μετὰ τοῦ ἀπὸ τῆς ΕΖ ἴσον ἐστὶ
τῷ ἀπὸ τῆς ΖΒ. διὰ τὰ α ὐτὰ δὴ καὶ τὸ 
 ὑπὸ  τῶν ΔΕ, ΕΒ μετὰ τοῦ ἀπὸ τῆς
ΖΕ ἰσον ἐστὶ τῷ ἀπὸ τῆς ΖΒ. ἐδείθθη δὲ καὶ τὸ ὑπὸ
τῶν ΑΕ, ΕΓ μετὰ τοῦ ἀπὸ τῆς ΖΕ ἴσον τῷ ἀπὸ τῆς
ΖΒ; τὸ ἄρα ὑπὸ τῶν ΑΕ, ΕΓ μετὰ τοῦ ἀπὸ τῆς ΖΕ ἴσον
ἐστὶ τῷ ὑπὸ τῶν
ΔΕ, ΕΒ μετὰ τοῦ ἀπὸ τῆς ΖΕ. κοινὸν ἀφῇρήσθω
τὸ ἀπὸ τῆς ΖΕ; λοιπὸν ἄρα τὸ ὑπὸ
τῶν ΑΕ, ΕΓ περιεχόμενον ὀρθογώνιον ἴσον ἐστὶ τῷ ὑπὸ
τῶν ΔΕ, ΕΒ περιεθομένῳ ὀρθογωνίῳ.}\\~\\~\\~\\~\\~\\~\\

\epsfysize=1.5in
\centerline{\epsffile{Book03/fig35g.eps}}

{\greekfont Ἐὰν ἄρα ἐν κύκλῳ εὐθεῖαι δύο τέμνωσιν ἀλλήλας, τὸ ὑπὸ τῶν
τῆς μιᾶς τμ ημάτων περιεχόμενον ὀρθογώνιον ἴσον ἐστὶ
τῷ ὑπὸ τῶν τῆς ἑτέρας τμ ημάτων περιεθομένῳ ὀρθογωνίῳ;
ὅπερ ἔδει δεῖχαι.}}

\ParallelRText{
\begin{center}
{\large Proposition 35}
\end{center}

If two straight-lines in a circle cut one another, (then) the rectangle
contained by the pieces of one is equal to the rectangle contained
by the pieces of the other.

For let the two straight-lines $AC$ and $BD$, in the circle $ABCD$,
cut one another at  point $E$. I say that the rectangle contained by
$AE$ and $EC$ is equal to the rectangle contained  by $DE$ and $EB$.

In fact, if $AC$ and $BD$ are through the center (as in the first diagram from the left), so that $E$ is the center
of circle $ABCD$, (then it is) clear that, $AE$, $EC$, $DE$, and $EB$ being equal,
the rectangle contained by $AE$ and $EC$ is also equal to the rectangle
contained by $DE$ and $EB$.

So let $AC$ and $DB$ not be though the center (as in the second diagram from the left), and let the center of $ABCD$ 
be found [Prop.~3.1], and let it be (at) $F$. And let $FG$ and $FH$
be drawn from $F$, perpendicular to the straight-lines $AC$ and $DB$
(respectively) [Prop.~1.12]. And let $FB$, $FC$, and $FE$ 
be joined.

And since some straight-line, $GF$, through the center, cuts at right-angles  some (other) straight-line, $AC$, not through the center,  (then) it also cuts
it in half [Prop.~3.3]. Thus, $AG$ (is) equal to $GC$. Therefore,
since the straight-line $AC$ is cut equally at $G$, and unequally at $E$,
the rectangle contained by $AE$ and $EC$ plus the square on $EG$ is thus
equal to the (square) on $GC$ [Prop.~2.5]. Let the (square) on
$GF$ be added [to both]. Thus, the (rectangle contained) by 
$AE$ and $EC$ plus the (sum of the squares) on $GE$ and $GF$ is equal to
the (sum of the squares) on $CG$ and $GF$. But, the (square) on $FE$ is equal to the (sum of the squares)
on $EG$ and $GF$   [Prop.~1.47],
and the (square)
on $FC$  is equal to the (sum of the squares) on $CG$ and $GF$ [Prop.~1.47]. Thus, the (rectangle contained) by
$AE$ and $EC$ plus the (square) on $FE$ is equal to the (square) on $FC$.
And $FC$ (is) equal to $FB$. Thus, the (rectangle contained) by
$AE$ and $EC$ plus the (square) on $FE$ is equal to the (square) on $FB$.
So, for the same (reasons), the (rectangle contained) by
$DE$ and $EB$ plus the (square) on $FE$ is equal to the (square) on $FB$.
And the (rectangle contained) by
$AE$ and $EC$ plus the (square) on $FE$ was also shown (to be) equal to the (square) on $FB$. Thus,  the (rectangle contained) by
$AE$ and $EC$ plus the (square) on $FE$ is equal to the (rectangle contained) by
$DE$ and $EB$ plus the (square) on $FE$. Let the (square) on $FE$ be
taken from both. Thus, the remaining rectangle contained by $AE$ and
$EC$ is equal to the rectangle contained by $DE$ and $EB$.

\epsfysize=1.5in
\centerline{\epsffile{Book03/fig35e.eps}}

Thus, if two straight-lines in a circle cut one another, (then) the rectangle
contained by the pieces of one is equal to the rectangle contained
by the pieces of the other. (Which is) the very thing it was required to
show.}
\end{Parallel}

%%%%%%
% Prop 3.36
%%%%%%
\pdfbookmark[1]{Proposition 3.36}{pdf3.36}
\begin{Parallel}{}{} 
\ParallelLText{
\begin{center}
{\large \ggn{36}.}
\end{center}\vspace*{-7pt}

{\greekfont Ἐὰν κύκλου ληφθῇ τι σημεῖον ἐκτός, καὶ ἀπ α ὐτοῦ πρὸς τὸν
κύκλον προσπίπτωσι δύο εὐθεῖαι, καὶ ἡ μὲν α ὐτῶν τέμνῃ τὸν
κύκλον, ἡ δὲ ἐφάπτηται, ἔσται τὸ ὑπὸ ὅλης τῆς τεμνούσης
καὶ τῆς ἐκτὸς ἀπολαμβανομένης μεταχὺ τοῦ τε σημείου καὶ
τῆς κυρτῆς περιφερείας ἴσον τῷ ἀπὸ τῆς ἐφαπτομένης τετραγώνῳ.}

{\greekfont Κύκλου γὰρ τοῦ ΑΒΓ εἰλήφθω τι σημεῖον ἐκτὸς τὸ Δ, καὶ ἀπὸ
τοῦ Δ πρὸς τὸν ΑΒΓ κύκλον προσπιπτέτωσαν δύο εὐθεῖαι αἱ ΔΓ[Α],
ΔΒ; καὶ ἡ μὲν ΔΓΑ τεμνέτω τὸν ΑΒΓ κύκλον, ἡ δὲ ΒΔ ἐφαπτέσθω;
λέγω, ὅτι τὸ ὑπὸ τῶν ΑΔ, ΔΓ περιεχόμενον ὀρθογώνιον
ἴσον ἐστὶ τῷ ἀπὸ τῆς ΔΒ τετραγώνῳ.}\\

\epsfysize=1.8in
\centerline{\epsffile{Book03/fig36g.eps}}

{\greekfont Ἡ ἄρα [Δ]ΓΑ ἤτοι διὰ τοῦ κέντρου ἐστὶν ἢ οὔ. ἔστω πρότερον
διὰ τοῦ κέντρου, καὶ ἔστω τὸ Ζ κέντρον τοῦ ΑΒΓ κύκλου, καὶ ἐπεζεύθθω ἡ ΖΒ; ὀρθὴ ἄρα ἐστὶν ἡ ὑπὸ ΖΒΔ. καὶ ἐπεὶ εὐθεῖα
ἡ ΑΓ δίθα τέτμ ηται κατὰ τὸ Ζ, πρόσκειται δὲ α ὐτῇ ἡ ΓΔ, τὸ
ἄρα ὑπὸ τῶν ΑΔ, ΔΓ μετὰ τοῦ ἀπὸ τῆς ΖΓ ἴσον
ἐστὶ τῷ ἀπὸ τῆς ΖΔ. ἴση δὲ ἡ ΖΓ τῇ ΖΒ; τὸ ἄρα ὑπὸ
τῶν ΑΔ, ΔΓ μετὰ τοῦ ἀπὸ τῆς ΖΒ ἴσον ἐστὶ τῷ ἀπὸ τῆς ΖΔ.
 τῷ δὲ ἀπὸ τῆς ΖΔ ἴσα
 ἐστὶ τὰ ἀπὸ τῶν ΖΒ, ΒΔ; τὸ ἄρα ὑπὸ τῶν ΑΔ, ΔΓ μετὰ
τοῦ ἀπὸ τῆς ΖΒ ἴσον ἐστὶ τοῖς ἀπὸ τῶν ΖΒ, ΒΔ. κοινὸν
ἀφῃρήσθω τὸ ἀπὸ  	τῆς ΖΒ; λοιπὸν ἄρα τὸ ὑπὸ τῶν ΑΔ, ΔΓ
ἴσον ἐστὶ τῷ ἀπὸ τῆς ΔΒ ἐφαπτομένης.}

{\greekfont Ἀλλὰ δὴ ἡ ΔΓΑ μὴ ἔστω διὰ τοῦ κέντρου τοῦ ΑΒΓ κύκλου,
καὶ εἰλήφθω τὸ κέντρον τὸ Ε, καὶ 
ἀπὸ τοῦ Ε ἐπὶ τὴν ΑΓ κάθετος ἤθθω ἡ ΕΖ, καὶ ἐπεζεύθθωσαν
αἱ ΕΒ, ΕΓ, ΕΔ; ὀρθὴ ἄρα ἐστὶν ἡ ὑπὸ ΕΒΔ. καὶ ἐπεὶ εὐθεῖά
τις διὰ τοῦ κέντρου ἡ ΕΖ εὐθεῖάν τινα μὴ διὰ τοῦ κέντρου τὴν
ΑΓ πρὸς ὀρθὰς τέμνει, καὶ δίθα α ὐτὴν τέμνει; ἡ ΑΖ ἄρα τῇ
ΖΓ ἐστιν ἴση. καὶ ἐπεὶ εὐθεῖα ἡ ΑΓ τέτμ ηται δίθα κατὰ τὸ Ζ
σημεῖον, πρόσκειται δὲ α ὐτῇ ἡ ΓΔ, τὸ ἄρα ὑπὸ τῶν ΑΔ, ΔΓ μετὰ
τοῦ ἀπὸ τῆς ΖΓ ἴσον ἐστὶ τῷ ἀπὸ τῆς ΖΔ. κοινὸν προσκείσθω τὸ
ἀπὸ τῆς ΖΕ; τὸ ἄρα ὑπὸ τῶν ΑΔ, ΔΓ μετὰ τῶν ἀπὸ
τῶν ΓΖ, ΖΕ ἴσον ἐστὶ τοῖς ἀπὸ τῶν ΖΔ, ΖΕ. τοῖς δὲ ἀπὸ
τῶν ΓΖ, ΖΕ ἴσον ἐστὶ τὸ ἀπὸ τῆς ΕΓ; ὀρθὴ γὰρ [ἐστιν]
 ἡ ὑπὸ ΕΖΓ [γωνία]; τοῖς δὲ ἀπὸ τῶν ΔΖ, ΖΕ ἴσον ἐστὶ
 τὸ ἀπὸ τῆς ΕΔ; τὸ ἄρα ὑπὸ τῶν ΑΔ, ΔΓ μετὰ τοῦ ἀπὸ
 τῆς ΕΓ ἴσον ἐστὶ τῷ ἀπὸ τῆς ΕΔ. ἴση
 δὲ ἡ ΕΓ τῂ ΕΒ; 
 τὸ ἄρα ὑπὸ τῶν ΑΔ, ΔΓ μετὰ τοῦ ἀπὸ τῆς ΕΒ ἴσον
 ἐστὶ τῷ ἄπὸ τῆς ΕΔ.
  τῷ δὲ ἀπὸ τῆς ΕΔ
 ἴσα ἐστὶ τὰ ἀπὸ τῶν ΕΒ, ΒΔ; ὀρθὴ γὰρ ἡ ὑπὸ ΕΒΔ γωνία;
 τὸ ἄρα ὑπὸ τῶν ΑΔ, ΔΓ μετὰ τοῦ ἀπὸ 
τῆς ΕΒ ἴσον 
 ἐστὶ τοῖς ἀπὸ τῶν ΕΒ, ΒΔ. κοινὸν ἀφῃρήσθω τὸ ἀπὸ τῆς
 ΕΒ; λοιπὸν ἄρα τὸ ὑπὸ τῶν ΑΔ, ΔΓ ἴσον ἐστὶ τῷ ἀπὸ τῆς
 ΔΒ.}
 
 {\greekfont Ἐὰν ἄρα κύκλου ληφθῇ τι σημεῖον ἐκτός, καὶ ἀπ α ὐτοῦ πρὸς τὸν
κύκλον προσπίπτωσι δύο εὐθεῖαι, καὶ ἡ μὲν α ὐτῶν τέμνῃ τὸν
κύκλον, ἡ δὲ ἐφάπτηται, ἔσται τὸ ὑπὸ ὅλης τῆς τεμνούσης
καὶ τῆς ἐκτὸς ἀπολαμβανομένης μεταχὺ τοῦ τε σημείου καὶ
τῆς κυρτῆς περιφερείας ἴσον τῷ ἀπὸ τῆς ἐφαπτομένης τετραγώνῳ;
ὅπερ ἔδει δεῖχαι.}}

\ParallelRText{
\begin{center}
{\large Proposition 36}\end{center}

If some point is taken outside a circle, and  two straight-lines
radiate from it towards the circle, and (one) of them cuts the circle, and the (other)
touches (it), (then) the (rectangle contained) by the whole  (straight-line) cutting (the circle), and the (part of it) 
cut off outside (the circle),
between the point and the convex circumference, will be equal to the
square on the tangent (line).

For let some point $D$ be taken outside circle $ABC$, and let two
straight-lines, $DC[A]$ and $DB$, radiate from $D$ towards circle $ABC$. And
let $DCA$ cut circle $ABC$, and let $BD$ touch (it). I say that the rectangle
contained by $AD$ and $DC$ is equal to the square on $DB$.

\epsfysize=1.8in
\centerline{\epsffile{Book03/fig36e.eps}}

$[D]CA$ is surely either through the center, or not. Let it first of all be
through the center, and let $F$ be the center of circle $ABC$, and let 
$FB$ be joined. Thus, (angle) $FBD$ is a right-angle [Prop.~3.18].
And since straight-line $AC$ is cut in half at $F$, let $CD$ be added to it.
Thus, the (rectangle contained) by $AD$ and $DC$ plus the (square) on $FC$ is
equal to the (square) on $FD$ [Prop.~2.6]. And $FC$ (is) equal to $FB$.
Thus, the (rectangle contained) by $AD$ and $DC$ plus the (square) on $FB$
is equal to the (square) on $FD$. And the (square) on $FD$ is equal to
the (sum of the squares) on $FB$ and $BD$ [Prop.~1.47]. Thus, the
(rectangle contained) by $AD$ and $DC$ plus the (square) on $FB$ is equal
to the (sum of the squares) on $FB$ and $BD$. Let the (square) on $FB$ be subtracted
from both. Thus, the remaining (rectangle contained) by $AD$ and $DC$ 
is equal to the (square) on the tangent $DB$.

And so let $DCA$ not be through the center of circle $ABC$, and let the center $E$
be found, and let $EF$ be drawn from $E$, perpendicular to
$AC$ [Prop.~1.12]. And let $EB$, $EC$, and $ED$ be joined. (Angle)
$EBD$ (is) thus a right-angle [Prop.~3.18]. And since some straight-line, $EF$, through the
center, cuts some (other) straight-line, $AC$, not through the center, at
right-angles, it also cuts it in half [Prop.~3.3]. Thus, $AF$ is equal to
$FC$. And since the straight-line $AC$ is cut in half at point $F$, let $CD$ 
be added to it. Thus, the (rectangle contained) by $AD$ and
 $DC$ plus the
(square) on $FC$ is equal to the (square) on $FD$ [Prop.~2.6]. Let
the (square) on $FE$ be added to both. Thus, the (rectangle contained)
by $AD$ and $DC$ plus the (sum of the squares) on $CF$ and $FE$ is equal to
the (sum of the squares) on $FD$ and $FE$. But  the (square) on $EC$  is equal to the (sum of the squares)
on $CF$ and $FE$. For [angle] $EFC$ [is] a
right-angle [Prop.~1.47]. And the (square) on $ED$ is equal to the (sum of the squares) on $DF$ and
$FE$   [Prop.~1.47]. Thus, the (rectangle
contained) by $AD$ and $DC$ plus the (square) on $EC$ is equal to the (square)
on $ED$. And $EC$ (is) equal to $EB$. Thus, the (rectangle
contained) by $AD$ and $DC$ plus the (square) on $EB$ is equal to the (square)
on $ED$. And the (sum of the squares) on 
$EB$ and $BD$  is equal to  the (square) on $ED$. For $EBD$ (is) a right-angle [Prop.~1.47]. Thus, the (rectangle
contained) by $AD$ and $DC$ plus the (square) on $EB$ is equal to the (sum of the squares)
on $EB$ and $BD$. Let the (square) on $EB$ be subtracted from both. Thus, the remaining
(rectangle contained) by $AD$ and $DC$ is equal to the (square) on $BD$.

Thus, if some point is taken outside a circle, and  two straight-lines
radiate from it towards the circle, and (one) of them cuts the circle, and (the other) touches (it), (then) the (rectangle contained) by the whole  (straight-line) cutting (the circle), and the (part of it) 
cut off outside (the circle),
between the point and the convex circumference, will be equal to the
square on the tangent (line). (Which is) the very thing it was required to show.}
\end{Parallel}

%%%%%%
% Prop 3.37
%%%%%%
\pdfbookmark[1]{Proposition 3.37}{pdf3.37}
\begin{Parallel}{}{} 
\ParallelLText{
\begin{center}
{\large \ggn{37}.}
\end{center}\vspace*{-7pt}

{\greekfont Ἐὰν κύκλου ληφθῇ τι σημεῖον ἐκτός, ἀπὸ δὲ τοῦ σημείου πρὸς τὸν
κύκλον προσπίπτωσι δύο εὐθεῖαι, καὶ ἡ μὲν α ὐτῶν τέμνῃ
τὸν κύκλον, ἡ δὲ προσπίπτῃ, ᾖ δὲ τὸ ὑπὸ [τῆς] ὅλης τῆς
τεμνούσης καὶ τῆς ἐκτὸς ἀπολαμβανομένης μεταχὺ τοῦ τε σημείου καὶ
τῆς κυρτῆς περιφερείας ἴσον τῷ ἀπὸ τῆς προσπιπτούσης, ἡ
προσπίπτουσα ἐφάψεται τοῦ κύκλου.}

{\greekfont Κύκλου γὰρ τοῦ ΑΒΓ εἰλήφθω τι σημεῖον ἐκτὸς τὸ Δ, καὶ ἀπὸ
τοῦ Δ πρὸς τὸν ΑΒΓ κύκλον προσπιπτέτωσαν δύο εὐθεῖαι αἱ
ΔΓΑ, ΔΒ, καὶ ἡ μὲν ΔΓΑ τεμνέτω τὸν κύκλον, ἡ δὲ ΔΒ προσπιπτέτω,
ἔστω δὲ τὸ ὑπὸ τῶν ΑΔ, ΔΓ ἴσον τῷ ἀπὸ τῆς ΔΒ. λέγω,
ὅτι ἡ ΔΒ ἐφάπτεται τοῦ ΑΒΓ κύκλου.}\\~\\

\epsfysize=1.75in
\centerline{\epsffile{Book03/fig37g.eps}}

{\greekfont Ἤθθω γὰρ τοῦ ΑΒΓ ἐφαπτομένη ἡ ΔΕ, καὶ εἰλήφθω τὸ κέντρον
τοῦ ΑΒΓ κύκλου, καὶ ἔστω τὸ Ζ, καὶ ἐπεζεύθθωσαν αἱ ΖΕ, ΖΒ, ΖΔ.
ἡ ἄρα ὑπὸ ΖΕΔ ὀρθή ἐστιν. καὶ ἐπεὶ ἡ ΔΕ ἐφάπτεται τοῦ
ΑΒΓ κύκλου, τέμνει δὲ ἡ ΔΓΑ, τὸ ἄρα ὑπὸ τῶν ΑΔ, ΔΓ
ἴσον ἐστὶ τῷ ἀπὸ τῆς ΔΕ. ἦν δὲ καὶ τὸ ὑπὸ τῶν
ΑΔ, ΔΓ ἴσον τῷ ἀπὸ τῆς ΔΒ; τὸ ἄρα ἀπὸ τῆς ΔΕ ἴσον ἐστὶ τῷ ἀπὸ
τῆς ΔΒ;
ἴση ἄρα ἡ ΔΕ τῇ ΔΒ.
ἐστὶ δὲ καὶ ἡ ΖΕ τῇ ΖΒ ἴση; δύο δὴ αἱ ΔΕ, ΕΖ δύο ταῖς ΔΒ, ΒΖ
ἴσαι εἰσίν; καὶ βάσις α ὐτῶν κοινὴ ἡ ΖΔ; γωνία ἄρα ἡ ὑπὸ
ΔΕΖ γωνίᾳ τῇ ὑπὸ ΔΒΖ ἐστιν ἴση. ὀρθὴ δὲ ἡ ὑπὸ ΔΕΖ; ὀρθὴ
ἄρα καὶ ἡ ὑπὸ ΔΒΖ. καί ἐστιν ἡ ΖΒ ἐκβαλλομένη διάμετρος;
ἡ δὲ τῇ διαμέτρῳ τοῦ κύκλου πρὸς ὀρθὰς ἀπ
ἄκρας ἀγομένη ἐφάπτεται τοῦ κύκλου; ἡ ΔΒ ἄρα ἐφάπτεται
τοῦ ΑΒΓ κύκλου. ὁμοίως δὴ δειθθήσεται, κἂν τὸ κέντρον  ἐπὶ
τῆς ΑΓ τυγθάνῃ.}

{\greekfont Ἐὰν ἄρα κύκλου ληφθῇ τι σημεῖον ἐκτός, ἀπὸ δὲ τοῦ σημείου πρὸς τὸν
κύκλον προσπίπτωσι δύο εὐθεῖαι, καὶ ἡ μὲν α ὐτῶν τέμνῃ
τὸν κύκλον, ἡ δὲ προσπίπτῃ, ᾖ δὲ τὸ ὑπὸ  ὅλης τῆς
τεμνούσης καὶ τῆς ἐκτὸς ἀπολαμβανομένης μεταχὺ τοῦ τε σημείου καὶ
τῆς κυρτῆς περιφερείας ἴσον τῷ ἀπὸ τῆς προσπιπτούσης, ἡ
προσπίπτουσα ἐφάψεται τοῦ κύκλου; ὅπερ ἔδει δεῖχαι.}}

\ParallelRText{
\begin{center}
{\large Proposition 37}
\end{center}

If some point is taken outside a circle, and two straight-lines radiate
from the point towards the circle, and one of them cuts the circle, and the
(other) meets (it), and the (rectangle contained) by the whole (straight-line)
cutting (the circle), and the (part of it) 
cut off outside (the circle),
between the point and the convex circumference,  is equal to the (square)
on the (straight-line) meeting (the circle), (then) the (straight-line)
meeting (the circle) will touch the circle.

For let some point $D$ be taken outside circle $ABC$, and let two
straight-lines, $DCA$ and $DB$, radiate from $D$ towards circle $ABC$, and
let $DCA$ cut the circle, and let $DB$ meet (the circle). And let the (rectangle
contained) by $AD$ and $DC$ be equal to the (square) on $DB$. I say that $DB$ touches circle $ABC$.

\epsfysize=1.75in
\centerline{\epsffile{Book03/fig37e.eps}}

For let $DE$ be drawn touching $ABC$ [Prop. 3.17], and let the center of the circle
$ABC$ be found, and let it be (at) $F$. And let $FE$, $FB$, and $FD$ 
be joined. (Angle) $FED$ is thus a right-angle [Prop.~3.18].
And since $DE$ touches circle $ABC$, and $DCA$ cuts (it), the (rectangle contained) by $AD$ and $DC$ is thus equal to the (square) on $DE$ [Prop.~3.36]. And the (rectangle contained) by $AD$ and $DC$ was also equal
to the (square) on $DB$.  Thus, the (square) on $DE$ is equal to the
(square) on $DB$. Thus, $DE$ (is) equal to $DB$. And $FE$ is also equal to
$FB$. So the two (straight-lines) $DE$, $EF$ are equal to the two (straight-lines)
$DB$, $BF$ (respectively).  And their base, $FD$, is common. Thus, angle $DEF$ is equal to
angle $DBF$ [Prop.~1.8]. And $DEF$ (is) a right-angle. Thus, $DBF$
(is) also a right-angle. And $FB$ produced is a diameter, And a (straight-line)
drawn at right-angles to a diameter of a circle, at its extremity, touches the
circle [Prop.~3.16~corr.]. Thus, $DB$ touches circle $ABC$. 
Similarly, (the same thing) can be shown, even if the center happens to be
on $AC$.

Thus, if some point is taken outside a circle, and two straight-lines radiate
from the point towards the circle, and one of them cuts the circle, and the
(other) meets (it), and the (rectangle contained) by the whole (straight-line)
cutting (the circle), and the (part of it) 
cut off outside (the circle),
between the point and the convex circumference,  is equal to the (square)
on the (straight-line) meeting (the circle), (then) the (straight-line)
meeting (the circle) will touch the circle. (Which is) the very thing it was
required to show.}
\end{Parallel}

\newpage~\thispagestyle{empty}
